\documentclass[10pt,a4paper]{article}

% ----------------- Paquetes -------------------%
\usepackage[T1]{fontenc}
\usepackage[utf8]{inputenc}
\usepackage[a4paper,margin=2.5cm]{geometry}
\usepackage{mathptmx}             
\usepackage{changepage} 
\usepackage{setspace}
\usepackage[spanish]{babel}
\usepackage[labelfont=bf,labelsep=space]{caption}
\usepackage{amssymb}
\usepackage{amsmath}
\usepackage{ebgaramond}
\usepackage{placeins}
\usepackage{etoolbox}
\usepackage{setspace}
\usepackage{parskip} 
\usepackage{tcolorbox}
\usepackage{needspace}
\usepackage{xparse}
\usepackage{ocgx2}
\usepackage{hyperref}
\usepackage{hypcap}


%%%%%%%%%%%%%%%%%%%%%%%%%%%%%%%%%%%%%%%%%%%%%%%%%%%%%%%%%%%%%%%%
% --- Fija primero tus valores globales "buenos" ---
\setstretch{1.2}                 % Interlineado global
\setlength{\parskip}{0.7\baselineskip}
\setlength{\parindent}{0pt}
\newlength{\GlobalParskip}
\setlength{\GlobalParskip}{\parskip}

\newcounter{eqnum}[subsection]
\renewcommand{\theeqnum}{\arabic{eqnum}}
\newcounter{alphaitem}
\newcounter{numitem}

%% ------------------- Recuadros----------------------%%
\newcommand{\puntosclave}[1]{%
  \begin{tcolorbox}[
    colframe=black,
    title={Puntos clave},
    before upper={\setlength{\parindent}{0.5cm}\setlength{\parskip}{0.5\baselineskip}}
  ]
  #1
  \end{tcolorbox}
}

\newcommand{\modificaciones}[1]{
  \begin{tcolorbox}[
    colback=white,
    colframe=red,
    title={Modificaciones a realizar},
    before upper={\setlength{\parindent}{0.5cm}\setlength{\parskip}{0.5\baselineskip}}
  ]
  #1
  \end{tcolorbox}
}

%% ------------------- Preguntas TEST----------------------%%
% Opciones: radio (single-choice)
\NewDocumentCommand{\OptRadio}{m m m}{%
  \noindent\CheckBox[radio,name=#1,value=#2,width=1.2ex,height=1.2ex]{}%
  \;\textbf{#2.}~#3\par
}

% Opciones: checkbox (multi-choice)
\NewDocumentCommand{\OptCheck}{m m}{%
  \noindent\CheckBox[name=#1,width=1.2ex,height=1.2ex,bordercolor={0 0 0}]{}%
  \;#2\par
}

% Pregunta single-choice (radio)
% Args: 1=id, 2=enunciado, 3=opciones, 4=solución+justificación, 5=imagen (o {}), 6=origen
\NewDocumentCommand{\PreguntaTestSingle}{m m m m m m}{%
  \par\medskip
  \noindent\textbf{#1.}~#2\par\smallskip

    % Imagen (si no es {})
    \IfBlankTF{#5}{}{%
      \begin{center}
        \includegraphics[width=0.75\linewidth]{#5}
      \end{center}
    }

    \begin{Form}
      #3
    \end{Form}

    \par\smallskip
    \switchocg{sol-#1}{\fbox{Mostrar / ocultar solución}}%
    \hfill{\footnotesize #6}\par\smallskip

    \begin{ocg}{Solución}{sol-#1}{0}
      \noindent #4
    \end{ocg}
  \par\medskip
}

% Pregunta multi-choice (checkboxes)
\NewDocumentCommand{\PreguntaTestMulti}{m m m m m m}{%
  \par\medskip
  \noindent\textbf{#1.}~#2\par\smallskip

    % Imagen (si no es {})
    \IfBlankTF{#5}{}{%
      \begin{center}
        \includegraphics[width=0.75\linewidth]{#5}
      \end{center}
    }
    \begin{Form}
      #3
    \end{Form}

    \par\smallskip
    \switchocg{sol-#1}{\fbox{Mostrar / ocultar solución}}%
    \hfill{\footnotesize #6}\par\smallskip

    \begin{ocg}{Solución}{sol-#1}{0}
      \noindent #4
    \end{ocg}
  \par\medskip
}

%% ------------------- Listas----------------------%%
    % --- Lista alfabética ---
    \newenvironment{la}
      {\begin{list}{\alph{alphaitem})}{%
          \usecounter{alphaitem}%
          \setlength{\leftmargin}{1cm}%
          \setlength{\parskip}{3pt}% 
          \setlength{\itemsep}{0pt}% ← espacio entre ítems
          \setlength{\parsep}{0pt}%  ← párrafos dentro de un ítem
      }}
      {\end{list}}
      
    % --- Lista numérica ---
    \newenvironment{lnum}
      {\begin{list}{\arabic{numitem}.}{%
          \usecounter{numitem}%
          \setlength{\leftmargin}{1cm}%
          \setlength{\parskip}{3pt}% 
          \setlength{\itemsep}{0pt}% ← espacio entre ítems
          \setlength{\parsep}{0pt}%  ← párrafos dentro de un ítem
      }}
      {\end{list}}
      
% ------------------- Imágenes----------------------%
    \graphicspath{{Img/}}% Rutas por defecto 
    % --- Imagen adaptativa ---
    % Registros de dimensión definidos una sola vez
    \newdimen\imgcap
    \newdimen\imgmin
    \newdimen\imgmax
    \newdimen\imgremain
    \newdimen\imgh
    \newdimen\imgneed
    
    \newcommand{\imagenfit}[3]{%
      \addvspace{10pt}% separación antes
      \par\begingroup
        % Parámetros de composición (asignaciones, no \newdimen)
        \imgcap=1\baselineskip            % alto estimado de título+fuente
        \imgmin=0.15\textheight
        \imgmax=0.15\textheight
        \imgremain=\pagegoal
        \ifdim\pagegoal=\maxdimen \imgremain=0pt\fi
        \advance\imgremain by -\pagetotal
        \advance\imgremain by -\imgcap
        % Decisión de altura destino
        \ifdim\imgremain<\imgmin
          \newpage
          \imgh=0.20\textheight
        \else
          \imgh=\imgremain
          \ifdim\imgh>\imgmax \imgh=\imgmax \fi
        \fi
        % Reservar espacio para evitar cortes del bloque completo
        \imgneed=\imgh \advance\imgneed by \imgcap
        \Needspace{\imgneed}
        % Cuerpo
        \noindent\begin{minipage}{\textwidth}
          \centering
          \captionof{figure}{\textemdash\ #2}\par\vspace{0.3em}% título arriba
          \includegraphics[width=\linewidth,height=\imgh,keepaspectratio]{\detokenize{#1}}\par
          \vspace{0.3em}\small\textit{Fuente: }#3% fuente abajo
        \end{minipage}%
      \endgroup
      \par\addvspace{10pt}%
    }

%------------ Bloque de RECUADRO ESPECIAL -------------%
    \newenvironment{specialblock}[1]{%
      \begin{quote} % crea sangría a izquierda y derecha
      \small % texto más pequeño
      \noindent\textbf{#1}\\[-0.3em] % título en negrita
      \rule{\linewidth}{0.4pt}\\[0.6em]}{
      \end{quote}}
      
%-------------------------- Portada -----------------------%
\newcommand{\oldtitle}[1]{\def\OldTitle{#1}}
\newcommand{\changerationale}[1]{\def\ChangeWhy{#1}}

\newcommand{\changebox}{%
\begin{tcolorbox}[title={Historial de cambios del tema}]
\textbf{Título anterior:} \OldTitle\par
\textbf{Motivo del cambio:} \ChangeWhy
\end{tcolorbox}
}

%-------------------------- Author Font -----------------------%
    \newcommand{\authorfont}[1]{{\fontfamily{EBGaramond-TLF}\selectfont #1}}

%-------------------------- Anexos -----------------------%
    \makeatletter
    \newcounter{anexo}
    \renewcommand{\theanexo}{\Roman{anexo}}
    \newcommand{\anexo}[1]{%
      \clearpage
      \phantomsection
      \refstepcounter{anexo}%
      \noindent\textbf{Anexo \theanexo.- }#1\par\medskip
      \addcontentsline{stc}{section}{Anexo \theanexo.- #1}%
    }
    \makeatother

    %-------------------------- Lista de figuras (personal) -----------------------%
\makeatletter
\newcommand{\MyListOfFigures}{%
  \clearpage
  \phantomsection
  {\centering\bfseries\Large Índice de figuras\par}%
  \medskip
  % Entrada en nuestro índice de contenido (stc)
  \addcontentsline{stc}{section}{Índice de figuras}%
  % Recuperación del listado (sin cabecera propia)
  \begingroup
    \parindent=0pt\parskip=2pt
    \@starttoc{lof}%
  \endgroup
}
\makeatother

%-------------------------- Bibliografía -----------------------%
\makeatletter
% Contador para las referencias
\newcounter{myref}
% Cabecera de la bibliografía, entra en el índice 'stc'
\newcommand{\MyBibliography}{%
  \clearpage
  \phantomsection
  {\centering\bfseries\Large Bibliografía y referencias\par}%
  \medskip
  \addcontentsline{stc}{section}{Bibliografía y referencias}%
  \setcounter{myref}{0}%
}
\newcommand{\MyBibItem}[4]{%
  \refstepcounter{myref}%
  \noindent[\themyref]\ #1.\ \emph{#2}. #3. #4\par
}
\makeatother

%--------- Establecer cuatro niveles de índice --------%
    \setcounter{tocdepth}{4}   
    \setcounter{secnumdepth}{4}% numera hasta \paragraph

%%%%%%%%%%%%%%%%%%%%%%%%%%%%%%%%

\makeatletter

    %--------- Interlineado Global --------%
    \edef\GlobalBaseLineStretch{\baselinestretch}
    
    %--------- Espaciado de seccinones --------%
    \renewcommand\section{\@startsection{section}
      {1}{\z@}
      {2.5ex \@plus 1ex \@minus .2ex} % espacio antes
      {1.5ex \@plus .2ex}             % espacio después
      {\normalfont\Large\bfseries}    % formato del título
    }
    \renewcommand\subsection{\@startsection{subsection}{2}{\z@}
      {3ex \@plus 0.8ex \@minus .2ex}
      {1.5ex \@plus .2ex}
      {\normalfont\large\bfseries}
    }
    \renewcommand\subsubsection{\@startsection{subsubsection}{3}{\z@}
      {3ex \@plus 0.5ex \@minus .2ex}
      {1ex \@plus .2ex}
      {\normalfont\normalsize\bfseries}
    }
    \renewcommand\paragraph{\@startsection{paragraph}{4}{\z@}%
    {-2ex \@plus -0.5ex \@minus -.2ex}% espacio antes
    {0.8ex \@plus .2ex}% espacio después
    {\normalfont\normalsize\itshape}}% ← cursiva en lugar de negrita

    %--------- Ecuaciones --------%   
    % Variables globales para almacenar títulos y notas
    \gdef\leq@current@section{}%
    \gdef\leq@current@subsection{}%
    \gdef\leq@last@printed@section{}%
    \gdef\leq@last@printed@subsection{}%
    
    % Comando para añadir nota (invisible en el cuerpo)
    \newcommand{\eqnote}[1]{%
      \addcontentsline{leq}{leqnote}{#1}%
    }
    
    % Comando para ecuaciones
    \newcommand{\eqblock}[2]{%
      % Verificar si necesitamos imprimir el título de sección
      \ifx\leq@current@section\leq@last@printed@section\else
        \addcontentsline{leq}{leqsection}{\leq@current@section}%
        \global\let\leq@last@printed@section\leq@current@section
        \gdef\leq@last@printed@subsection{}% Reset subsección al cambiar sección
      \fi
      % Verificar si necesitamos imprimir el título de subsección
      \ifx\leq@current@subsection\leq@last@printed@subsection\else
        \ifx\leq@current@subsection\@empty\else
          \addcontentsline{leq}{leqsubsection}{\leq@current@subsection}%
          \global\let\leq@last@printed@subsection\leq@current@subsection
        \fi
      \fi
      % Ecuación
      \refstepcounter{eqnum}%
      \begin{equation*}
        #1\tag{E.\theeqnum}
      \end{equation*}%
      \addcontentsline{leq}{leqt}{[E.\theeqnum] -- #2}%
      \addcontentsline{leq}{leqm}{#1}%
    }
    
    %%% ----- Índice de ecuaciones ----------%%%
    % Tamaño para []
    \newcommand{\leqIndexFont}{\normalsize}
    \newcommand{\leqTitleFont}{\tiny}
    \newcommand{\leqNoteFont}{\tiny}
    % Cabecera de sección 
    \newcommand*\l@leqsection[2]{%
      \addvspace{25pt}% Mayor separación antes de nueva sección
      \noindent\leqIndexFont
      \makebox[\linewidth]{\rule{.38\linewidth}{0.4pt}\ \ \textbf{  \MakeUppercase{#1}}\ \ \rule{.38\linewidth}{0.4pt}}%
      \par\vspace{-10pt}
    }
    % Cabecera de subsección 
    \newcommand*\l@leqsubsection[2]{%
      \par\addvspace{15pt}%
      \noindent\leqIndexFont #1
      \par\vspace{-7pt}
    }
    % Título de ecuación 
    \newcommand*\l@leqt[2]{%
    \par\addvspace{10pt}%
      \par\noindent\leqTitleFont #1\par
    }
    % Miniatura de ecuación
    \newcommand*\l@leqm[2]{%
      \vspace{0.3\baselineskip}%
      \par\scriptsize\noindent\hspace*{1.5em}\(#1\)\par%
    }
    % Nota de ecuación 
    \newcommand*\l@leqnote[2]{%
      \vspace{0.2\baselineskip}%
      \par\noindent\leqNoteFont\hspace*{1.5em}\textbf{Nota.--} #1\par\vspace{0.5\baselineskip}%
    }
    % Generación del índice
    \newcommand{\mylistofequations}{%
      \clearpage
      \phantomsection
      % Título visual de la lista (ajústalo a tu gusto):
      {\centering\bfseries\Large Índice de ecuaciones\par}
      % Entrada en nuestro índice 'stc':
      \addcontentsline{stc}{section}{Índice de ecuaciones}%
      % Llamada a tu lista real (usa el nombre exacto de tu macro):
      \@starttoc{leq}%
    }
    
    %--------- Captura de títulos de sección --------%
    \let\leq@original@section\section
    \let\leq@original@subsection\subsection
    % Flag para saber si estamos en una sección asterisco
    \newif\ifLeqStarSection
    \LeqStarSectionfalse
    
    % Redefinir \section CON CONTROL DE ASTERISCO
    \renewcommand{\section}{%
      \@ifstar{\leq@section@star}{\@dblarg\leq@section@nostar}%
    }
    \def\leq@section@star#1{%
      \global\LeqStarSectiontrue
      \gdef\leq@current@section{#1}%
      \gdef\leq@current@subsection{}%
      \leq@original@section*{#1}%
      \global\LeqStarSectionfalse
    }
    \def\leq@section@nostar[#1]#2{%
      \global\LeqStarSectionfalse
      \gdef\leq@current@section{#2}%
      \gdef\leq@current@subsection{}%
      \leq@original@section[#1]{#2}%
    }
    % Redefinir \subsection
    \renewcommand{\subsection}{%
      \@ifstar{\leq@subsection@star}{\@dblarg\leq@subsection@nostar}%
    }
    \def\leq@subsection@star#1{%
      \global\LeqStarSectiontrue
      \gdef\leq@current@subsection{#1}%
      \leq@original@subsection*{#1}%
      \global\LeqStarSectionfalse
    }
    \def\leq@subsection@nostar[#1]#2{%
      \global\LeqStarSectionfalse
      \gdef\leq@current@subsection{#2}%
      \leq@original@subsection[#1]{#2}%
    }

    % ----- Restauración de valores Globales de estilo ----------%
    % Flags
    \newif\ifInFootnote
    \InFootnotefalse
    \pretocmd{\@footnotetext}{\InFootnotetrue}{}{}
    \apptocmd{\@footnotetext}{\InFootnotefalse}{}{}
    % Profundidad de listas (soporta anidadas)
    \newcount\MyListDepth \MyListDepth=0
    \AtBeginEnvironment{la}{\advance\MyListDepth by 1}
    \AfterEndEnvironment{la}{\advance\MyListDepth by -1}
    \AtBeginEnvironment{lnum}{\advance\MyListDepth by 1}
    \AfterEndEnvironment{lnum}{\advance\MyListDepth by -1}
    \AtBeginEnvironment{itemize}{\advance\MyListDepth by 1}
    \AfterEndEnvironment{itemize}{\advance\MyListDepth by -1}
    \AtBeginEnvironment{enumerate}{\advance\MyListDepth by 1}
    \AfterEndEnvironment{enumerate}{\advance\MyListDepth by -1}
    % Restauración diferida: hazlo al INICIO del próximo párrafo real
    \newif\if@pendingRestore
    \@pendingRestorefalse
    \newcommand{\RequestRestore}{\global\@pendingRestoretrue}
    \everypar{%
      \if@pendingRestore
        \global\@pendingRestorefalse
        \ifInFootnote\else
          \ifnum\MyListDepth=0
            \setlength{\parskip}{\GlobalParskip}%
            \setstretch{\GlobalBaseLineStretch}%
          \fi
        \fi
      \fi
    }
    % Al final de bloques:
    \AfterEndEnvironment{equation}{\RequestRestore}
    \AfterEndEnvironment{equation*}{\RequestRestore}
    \AfterEndEnvironment{la}{\RequestRestore}
    \AfterEndEnvironment{lnum}{\RequestRestore}
    % Al final de macros:
    \apptocmd{\eqblock}{\RequestRestore}{}{}
    \apptocmd{\imagenfit}{\RequestRestore}{}{}
    
\makeatother

% ===================== ToC propio con 4 niveles (único bloque) =====================
\makeatletter

% --- Escritores: añadimos una línea al .stc en cada cabecera numerada ---
% Guardamos las versiones actuales para llamar al comportamiento normal
\let\MyTOC@orig@section\section
\let\MyTOC@orig@subsection\subsection
\let\MyTOC@orig@subsubsection\subsubsection
\let\MyTOC@orig@paragraph\paragraph

% SECTION
\renewcommand{\section}{%
  \@ifstar{\MyTOC@section@star}{\@dblarg\MyTOC@section@nostar}%
}
\def\MyTOC@section@star#1{%
  \MyTOC@orig@section*{#1}% sin entrada al .stc
}
\def\MyTOC@section@nostar[#1]#2{%
  \MyTOC@orig@section[{#1}]{#2}% comportamiento normal (escribe al .toc)
  % Escribimos también nuestra línea al .stc (texto con \numberline)
  \addcontentsline{stc}{section}{\protect\numberline{\thesection}#2}%
}

% SUBSECTION
\renewcommand{\subsection}{%
  \@ifstar{\MyTOC@subsection@star}{\@dblarg\MyTOC@subsection@nostar}%
}
\def\MyTOC@subsection@star#1{%
  \MyTOC@orig@subsection*{#1}%
}
\def\MyTOC@subsection@nostar[#1]#2{%
  \MyTOC@orig@subsection[{#1}]{#2}%
  \addcontentsline{stc}{subsection}{\protect\numberline{\thesubsection}#2}%
}

% SUBSUBSECTION
\renewcommand{\subsubsection}{%
  \@ifstar{\MyTOC@subsubsection@star}{\@dblarg\MyTOC@subsubsection@nostar}%
}
\def\MyTOC@subsubsection@star#1{%
  \MyTOC@orig@subsubsection*{#1}%
}
\def\MyTOC@subsubsection@nostar[#1]#2{%
  \MyTOC@orig@subsubsection[{#1}]{#2}%
  \addcontentsline{stc}{subsubsection}{\protect\numberline{\thesubsubsection}#2}%
}

% PARAGRAPH
\renewcommand{\paragraph}{%
  \@ifstar{\MyTOC@paragraph@star}{\@dblarg\MyTOC@paragraph@nostar}%
}
\def\MyTOC@paragraph@star#1{%
  \MyTOC@orig@paragraph*{#1}%
}
\def\MyTOC@paragraph@nostar[#1]#2{%
  \MyTOC@orig@paragraph[{#1}]{#2}%
  % Si no numeras los \paragraph, cambia \theparagraph por texto plano sin \numberline
  \addcontentsline{stc}{paragraph}{\protect\numberline{\theparagraph}#2}%
}

% --- Impresor del índice propio (lee el .stc) ---
\newcommand{\MyTOC}{%
  \par\bigskip
  {\centering\bfseries\Large Índice de contenido\par}%
  \medskip
  \begingroup
    \parindent=0pt\parskip=2pt
    \@starttoc{stc}%
  \endgroup
}

\makeatother
% ===================== fin ToC propio =====================


%%%%%%%%%%%%%%


% --- Datos del tema ---
\title{Crisis financieras y pánicos bancarios\\
\large Especial referencia a la crisis financiera internacional iniciada en 2007-2008. Gestión de riesgos de las instituciones financieras}
\author{Marcelo Ortega}
\date{Marzo 2026}
\newcommand{\TemaCode}{3B26}

\oldtitle{
    B.24. Mercados financieros internacionales: emisores, instrumentos y mercados de renta fija y variable
}
\changerationale{
    Se opta por un enfoque analítico del funcionamiento de los mercados financieros (crisis y prevención de crisis) en vez de uno esencialmente descriptivo, que además causaba ciertos solapamientos con varios temas del Tercer y del Cuarto Ejercicio.
}

\begin{document}


\maketitle
\changebox

\newpage

% --- Página de resumen y modificaciones ---
\puntosclave{ %% Escribir puntos clave Apuntes CECO
     
    }
\vspace{1cm}

\modificaciones{
    
    }

\newpage
\MyTOC
\newpage

\mylistofequations
\clearpage

\MyListOfFigures
\clearpage


\pagenumbering{arabic}
\setcounter{page}{1}


\section{Introducción}





\subsection{Contextualización}


\subsubsection{Relevancia}




\subsubsection{Problemática}



\subsection{Estructura de la exposición}




\clearpage
\section{Fundamento económico: ¿qué es el Sistema Financiero?}
\puntosclave{
    La función económica del Sistema Financiero es \textbf{facilitar la asignación y utilización de los recursos económicos, tanto espacial como temporalmente, en un entorno incierto}.

    Dentro deel desarrollo institucional de los mercados financiero, la caraterística central de los sujetos es la intensidad de las asimetrías informativas entre agentes privados y agentes institucionales (financieros). En este sentido, ROSS (1989) distingue entre mercados opacos, traslúcidos y transparentes. 

    El Sistema Financiero es el corazón de la economía: bombea, distribuye y canaliza todos los recursos nominales de la economía. De toda transacción o intermediación podríamos encontrar un punto de conexión con el Sistema Financiero, constituyendo de esta forma el \textbf{nodo central} de la actividad económica.
    
    La toma de riesgos es inherente a toda actividad de carácter financiero, y en especial a la bancaria. 
}

Existen dos marcos de referencia fundamentalmente distintos para el análisis del Sistema Financiero. Una perspectiva toma como dada la estructura institucional existente y considera que el objetivo de Política Económica es ayudar a las instituciones establecidas a sobrevivir y prosperar y evitar la inestabilidad. En definitiva, que presten sus servicios de manera más eficiente y rentable. Una alternativa a la perspectiva institucional es la perspectiva funcional (MERTON R., 1995). La perspectiva funcional toma como dadas las funciones económicas desempeñadas por los intermediarios financieros y se pregunta cuál es la mejor estructura institucional para desempeñar dichas funciones.\footnote{Para una discusión en profundidad de la perspectiva funcional como herramienta de análisis, véanse Zvi Bodie y Robert C. Merton (1993), así como Robert C. Merton y Zvi Bodie (1993 y 1995). Este modo de análisis es consistente con los enfoques desarrollados por Fischer Black (1985), Michael J. Brennan (1993), Didier Cossin (1993), Michael C. Jensen y William H. Meckling (1976), Hayne E. Leland y David H. Pyle (1977), James L. Pierce (1991), Myron S. Scholes y Mark A. Wolfson (1992), y Oliver E. Williamson (1985 y 1988). Las funciones financieras se utilizan dentro de un marco analítico diferente en los trabajos de Douglas W. Diamond y Philip H. Dybvig (1986), así como en Stuart I. Greenbaum y Richard C. Higgins (1983).
} Por tanto, a diferencia de la perspectiva institucional, esta perspectiva funcional no presupone que las instituciones existentes, ya sean operativas o regulatorias, vayan necesariamente a preservarse.

\textbf{¿Qué queremos decir?} A lo largo de este primer apartado, vamos a adoptar la perspectiva funcional con el objetivo de entender la función económica del Sistema Financiero, los fallos de mercado que éste presenta, las formas de categorizar la actividad financiera y los riesgos inherentes que presenta. Esto nos permitirá abordar los diferentes mecanismos de gestión de riesgo con unos fundamentos sólidos y, posteriormente, entender las razones por las que una mala gestión o una exposición inadecuada puede conducir al Sistema a una situación de crisis inevitable.

\subsection{Perspectiva Funcional del Sistema Financiero: ¿Cómo definir los objetivos de Política Económica?}
La perspectiva funcional del Sistema Financiero parte de dos premisas fundamentales: (i) las funciones financieras son más estables que las instituciones financieras; es decir, las funciones cambian menos con el tiempo y varían menos entre distintas fronteras geopolíticas; y, (ii) la competencia provocará que los cambios en la estructura institucional evolucionen hacia una mayor eficiencia en el desempeño del sistema financiero.

Por tanto, ¿cómo podría diseñarse un sistema financiero completamente nuevo para un país? Al construir un sistema financiero desde cero, resulta natural comenzar definiendo su papel central: \textbf{facilitar la asignación y utilización de los recursos económicos, tanto espacial como temporalmente, en un entorno incierto}.\footnote{MERTON (1994), va un paso más alla. A partir de este nivel más agregado de una única función principal de asignación de recursos, podemos distinguir seis funciones básicas desempeñadas por el sistema financiero:
\begin{lnum}
    \item \textbf{Función 1}. Un sistema financiero proporciona un sistema de pagos para el intercambio de bienes y servicios;
    \item \textbf{Función 2}. Un sistema financiero proporciona un mecanismo para la agrupación de fondos que permita emprender proyectos indivisibles de gran escala;
    \item \textbf{Función 3}. Un sistema financiero proporciona una forma de transferir recursos económicos a través del tiempo y entre distintas regiones geográficas e industrias;
    \item \textbf{Función 4}. Un sistema financiero proporciona un medio para gestionar la incertidumbre y controlar el riesgo;
    \item \textbf{Función 5}. Un sistema financiero proporciona información sobre precios que ayuda a coordinar la toma descentralizada de decisiones en diversos sectores de la economía;
    \item \textbf{Función 6}. Un sistema financiero proporciona un mecanismo para abordar los problemas de información asimétrica e incentivos que surgen cuando una de las partes de una transacción financiera posee información de la que la otra carece.
\end{lnum}
Dwight B. Crane, Kenneth A. Froot, Scott P. Mason, André F. Perold, Robert C. Merton, Zvi Bodie, Erik R. Sirri y Peter Tufano (1995) presentan un desarrollo ampliado del enfoque funcional, con seis capítulos dedicados a ilustrar cada una de estas seis funciones. Desgloses alternativos de las funciones financieras pueden encontrarse en los trabajos de R. Glenn Hubbard (1994), Meir Kohn (1994), Peter S. Rose (1994) y Anthony Sanford (1993).
} 

\textbf{DIBUJAR GRÁFICO: NEXO DE UNIÓN}

Por diversas razones, la estructura institucional más eficiente para desempeñar las funciones del sistema financiero cambia generalmente con el tiempo y difiere entre distintas jurisdicciones geopolíticas.\footnote{Además, incluso cuando la identidad corporativa de las instituciones es la misma, las funciones que desempeñan suelen diferir considerablemente. Por ejemplo, los bancos en Estados Unidos en 1995 eran muy distintos de los que existían en 1925 o en 1955, del mismo modo que difieren notablemente de las instituciones denominadas bancos en Alemania o en el Reino Unido en la actualidad.

Los mercados financieros de New York City, London o Tokyo son hoy muy diferentes de lo que eran tan recientemente como en 1980, antes de la introducción generalizada de contratos de futuros, opciones y permutas financieras (swaps) sobre renta fija e índices bursátiles.
} Por el contrario, las funciones básicas de un sistema financiero son esencialmente las mismas en todas las economías, pasadas y presentes, orientales y occidentales. Y, dado que las funciones del sistema financiero son mucho más estables que la identidad y la estructura de las instituciones que las desempeñan, una perspectiva funcional proporciona un marco de referencia más fiable y duradero que una perspectiva institucional, especialmente en un entorno financiero caracterizado por cambios rápidos.

\subsubsection{Tipos de ISF (ROSS, 1989): Las asimetrías informativas}
Ahora bien, es obvio que los mercados financieros son, al menos estáticamente, ciertamente institucionales. Por tanto, una vez conocemos el objeto o función económica del Sistema Financiero, sería propio conocer los sujetos (o como mínimo, la gama de sujetos) que lo conforman. 

En este sentido, podemos adoptar la clasificación presentada por ROSS (1989) que nos proporciona valiosas implicaciones acerca de la relación entre mercados y, sobretodo, las características del nexo entre agentes privados y la naturaleza de sus relaciones. Bajo una clasificación ampliada de las \textbf{instituciones} financieras distinguimos:

\begin{center}
\begin{tabular}{|p{0.09\linewidth}|p{0.09\linewidth}|p{0.09\linewidth}|p{0.09\linewidth}|p{0.07\linewidth}|p{0.07\linewidth}|p{0.07\linewidth}|p{0.09\linewidth}|p{0.09\linewidth}|}
\hline
\multicolumn{3}{|c|}{\textbf{Transparente}} &
\multicolumn{4}{|c|}{\textbf{Traslúcido}} &
\multicolumn{2}{c|}{\textbf{Opaco}}\\
\hline
\parbox{1\linewidth}{\centering\scriptsize Mercados de deuda\par\tiny(Bonos y otros instrumentos)} &
\parbox{1\linewidth}{\centering\scriptsize Mercados de valores\par\tiny(Acciones)} &
\parbox{1\linewidth}{\centering\scriptsize Mercados de instrumentos derivados\par\tiny(permuta, a plazo, etc.)} &
\parbox{1\linewidth}{\centering\scriptsize Fondos de Inversión Colectiva\par\tiny(Unit Trusts)} &
\parbox{1\linewidth}{\centering\scriptsize Fondos Mutualistas\par\tiny(Mutual Funds)} &
\parbox{1\linewidth}{\centering\scriptsize Fondos de Pensiones\par\tiny(Pension Funds)} &
\parbox{1\linewidth}{\centering\scriptsize Compañías Financieras\par\tiny(Financial Companies)} &
\parbox{1\linewidth}{\centering\scriptsize Compañías de Seguros\par\tiny(Insurance Companies)} &
\parbox{1\linewidth}{\centering\scriptsize Bancos Comerciales\par\tiny(Commercial Banking)}\\
\hline
\end{tabular}
\end{center}

¿Qué nos quiere decir ROSS con esta clasificación? Como podemos intuir, ROSS pretende enfatizar las asimetrías informativas presentes entre los agentes privados y las instituciones financieras. Como es evidente, caracterizaremos como instituciones transparentes aquellas que actúan como vehículos de transmisión (pass-throughs) y reflejan de forma bastante directa las fuerzas minoristas. Por el contrario, caracterizaremos como instituciones opacas aquellas ocultas y veladas a la vista de los participantes y donde solo se refleja débilmente las fuerzas minoristas, en el mejor de los casos. 

En esencia, todas las instituciones financieras están regidas por relaciones de agencia. Como sabemos, esto será especialmente relevante a la hora de abordar la tarea regulatoria, pues estas asimetrías traen por causa la existencia de un fallo de mercado que justificaría, a priori, la intervención del Sector Público. Por ejemplo, un depositante apenas tiene conocimiento de los préstamos específicos que concede la institución: ¿son locales o nacionales? ¿están garantizados o no? Muy poco de esto es visible para el participante. Lo mismo ocurre con una compañía de seguros que invierte las primas antes de pagar los siniestros, y la situación es aún peor para las empresas manufactureras cuyas actividades financieras constituyen un aspecto secundario de su negocio principal. 

No obstante, estas asimetrías y relaciones de agencia son también cruciales a la hora de entender cómo son las dinámicas del Sector Financiero, qué reisgos emergen y cómo estos pueden acaar provocando las tan témidas \textit{crisis financieras}. Además, pocos discutirían que a lo largo de los últimos cincuenta años, el volumen nocional del Sistema Financiero se ha traslado desde mercados opacos, donde prima la intermediación bancaria, a mercados más trasnparentes, en particular los mercados de deuda. 

\subsection{Centralidad del Sistema Financiero: Riesgo gestionado y riesgo desbordado}
Estamos en buena posición para abordar la relevancia del Sistema Financiero y las razones por las que debe ser, y es, fuente de preocupación generalizada. 

A primera vista, la presencia de información asimétrica no parece resolver adecuadamente esta cuestión: ¿no existen estas justificaciones en cualquiera otros sectores? ¿por qué no se intervienen en la misma medida?.\footnote{En el marco de la TOI, la regulación tiene como objetivo reducir las consecuencias derivadas de la presencia de fallos de mercado, tales como la competencia imperfecta o la información incompleta. Esto, evidentemente, puede ser perfectamente uno de los objetivos de la regulacíon financiera. Sin embargo, no parece suficiente para explicar la magnitud e importancia que tanto el legislador como la sociedad en su conjunto le atribuye a la regulación financiera. Por otro lado, una justificación funcional de la regulación excedería por mucho los límites que el regulador estaría dispuesto a reconocer: ¿quién atribuye esa función deseada/deseable o natural al sistema financiero? ¿cómo podemos garantizar que este corpus regulatorio concuerda con estos objetivos funcionales, incluso si todos estuvieramos de acuerdo con ellos? Como vemos, simplemente la incomodidad de exponerse a estas preguntas supondría un rechazo frontal a esta justificación regulatoria.
} De hecho, el poco contenido normativo que se desprende de la función económica que cumple tampoco nos ayuda a resolverlo. Sin duda necesitamos algo más. ¿Cuál es \textbf{realmente} la razón que hace del Sistema Financiero un objeto esencial de nuestras economías? Tenemos dos respuestas que se complementan. 

En primer lugar, el Sistema Financiero no es simplemente un canalizador. No actúa de caja negra a la que entran unos inputs, o recursos económicos, y se reasignan temporal y espacilamente hacia outputs más productivos. 

Todo lo contrario. El Sistema Financiero es el corazón de la economía: bombea, distribuye y canaliza todos los recursos nominales de la economía. Constituye el nodo central de la economía puesto que no reasigna únicamente recursos (crédito), sino que los gestiona (inversión), los trasnforma (riesgo y función), los almacena (depósito) y los canaliza (sistema de pagos): de toda transacciones o intermediación podríamos encontrar un punto de conexión con el Sistema Financiero, constituyendo de esta forma el \textbf{nodo central} de la actividad económica.

\textbf{INSERTAR GRÁFICO}

En segundo lugar, la esencia misma de su actividad supone un arma de doble filo: no solo se expone a riesgos de diverso origen (\textit{incoming risk}), sino que expone al resto de la economía a los riesgos asumidos por su sector (\textit{spillover risk}); generando un círculo de responsabilidades que debe ser adecuadamente controlado ya que, bajo una óptica tradicional, \textbf{el Sistema Financiero no es capaz de internalizar adecuadamente el coste social que supondría el colapso de su actividad}.

Por tanto, y en definitiva, es este par riesgo-centralidad del Sistema Financiero lo que constituye la fuente de preocupación e interés respecto al Sistema Financiero; siendo no solo la razón que se esgrime a la hora de justificar su intervención, sino también la razón por la que su observación resulta crucial para garantizar la buena salud de la economía.

En los 80, otros autores, como STIGLITZ y WEISS (1981)5, y GERTLER (1988) incluyeron el concepto de información asimétrica, es decir situaciones en las que una parte de los intervinientes en el mercado carece de toda la información que necesitaría conocer sobre la otra parte para tomar las decisiones correctas. Estos problemas de información asimétrica se pueden dar antes de que la transacción financiera se produzca –dando lugar a problemas de selección adversa, modelizados por primera vez por AKERLOF (1970)- o después de que la transacción financiera se produzca -dando lugar a problemas de riesgo moral, presentados en el trabajo seminal de ARROW (1970)-.

El problema de selección adversa en el sistema financiero surge porque el intermediario financiero es incapaz de determinar la calidad crediticia de aquellos dispuestos a participar en el mercado. Y, en determinadas situaciones esto podría llevarle a rechazar operaciones con un riesgo de impago bajo, por no ser capaz de distinguirlas de aquéllas con alto riesgo, dando lugar al racionamiento de crédito o incluso a no conceder ningún préstamo. Por su parte, el riesgo moral ilustrado por HOLMSTROM y TIROLE (1997) provoca que el agente (prestatario) tenga incentivos a invertir en proyectos arriesgados en los que si no obtiene rentabilidad el coste de las pérdidas es soportado en parte por el prestamista (principal). Este conflicto de interés entre prestatarios y prestamistas (problema de agencia) puede conllevar un nivel de crédito más bajo del que sería óptimo.

Así, MISHKIN (1999) considerando los problemas de información asimétrica como los desencadenantes de las crisis financieras, amplió la definición de crisis financieras definiéndolas como perturbaciones de los mercados financieros que agravan los posibles problemas de selección adversa y de riesgo moral, hasta tal punto que no es posible canalizar eficientemente los recursos hacia quienes tienen las oportunidades de inversión más productivas. Así, una crisis financiera haría a la economía pasar de un equilibrio en el que los mercados financieros funcionan bien a otro equilibrio en el que la producción disminuye drásticamente porque los mercados financieros no canalizan eficientemente los recursos.


\clearpage
\section{Instrumentos de Gestión del riesgo: cuantificación y mitigación del riesgo}
\puntosclave{
    Los riesgos financieros más importantes a los que se enfrentan las entidades de crédito son: el riesgo de crédito, el riesgo de mercado, el riesgo de liquidez y el riesgo operacional.

    Para evitar las consecuencias negativas para la economía y el sistema financiero de la mala gestión de riesgos de las entidades, éstas están sometidas a regulación. Sin embargo, esta regulación tan sólo supone un “suelo” legal que las entidades han de cumplir: la gestión de riesgos es un proceso interno de cada entidad en materia de evaluación, control, monitorización y mitigación de riesgos. Su objetivo es por tanto garantizar a la solvencia de la entidad, así como maximizar la rentabilidad del negocio.
    
    Se considera que el riesgo de crédito es el principal factor de riesgo que puede amenazar la solidez de un banco, de la mayoría de las instituciones financieras y del sistema financiero en su conjunto. La gestión del riesgo de crédito se hace tanto de forma previa a la adopción de una exposición con riesgo de crédito (con el fin de reducir problemas de información asimétrica) como a posteriori (mediante la monitorización de las posiciones de la entidad y la cobertura de las mismas ante el posible impago de parte de su cartera).
    
    Los modelos VaR, permiten realizar una estimación de la peor perdida monetaria en un determinado horizonte temporal. La crisis de 2007-08 puso de manifiesto las limitaciones de los modelos VaR como mecanismo de identificación y control de riesgos debido, entre otras razones, a que estos modelos no captan, los riesgos en la cola de la distribución de pérdidas e infravaloran el riesgo, al reflejar solo el riesgo percibido y no el real, debido a su dependencia de datos históricos.
}

Aunque es de sobra conocida la noción abstracta del riesgo, sobretodo en la actividad financiera, debemos ir un paso allá en nuestra labor de economistas: ¿cuáles son las fuentes de riesgo? ¿cómo podemos cuantificarlo? ¿por qué surge y se desborda este riesgo? El Comité de Basilea señala siete riesgos particulares, de entre los que trabajremos tres: riesgo de crédito, de mercado y liquidez.\footnote{En el anexo se tratan los demás riesgos así como sus principales instrumentos de gestión [\authorfont{Ver Anexo \ref{anx:riesgos_basilea}}]. 
}

\textbf{¿Qué queremos decir?} A continuación veremos las maneras de cuantificar y, sobretodo, gestionar estas fuentes de riesgo. Estas vías de gestión se recogen bajo el paraguas de los que conocemos como \textbf{Instrumentos de Gestión del Riesgo}. Cada uno de los riesgos tendrá su paleta de instrumentos adecuadamente pensada para este. Sin embargo, aunque puedan parecer compartimentos estancos, con la visión presentada deberíamos tener claro que se trata de instrumentos complementarios y no alternativos ni subsidiarios. Hemos prestado especial atención a los fundamentos del Sistema con este fin: comprender que se trata de un \textbf{sistema} y no un \textbf{mercado} parcial. Como tal, tiene riesgos que pueden ser óptimamente mitigados de forma individualizada, y otros que necesariamente requieren una labor de abstracción, observación y análisis del comportamiento conjunto, con el fin de identificar lo que conocemos como \textbf{riesgos sistémicos}.

\subsection{Riesgo de crédito: Incumplimiento de obligaciones}
El primero de los riesgos que trataremos será el riesgo de crédito: la posibilidad de que un prestatario o contraparte incumpla sus obligaciones contractuales. El Comité de Basilea lo define exactamente en esos términos y exige que sea gestionado tanto a nivel de operación individual como de cartera, porque la concentración, la correlación y el deterioro cíclico convierten pérdidas idiosincráticas en pérdidas agregadas.\footnote{Es relevante tener presente que las fuentes de este riesgo pueden ser individuales (empeora la calidad crediticia de un deudor, o ésta no fue evaluada correctamente inicialmente) o sistémicas (empeora la situación económica general, lo que afecta a la capacidad crediticia de los agentes económicos). Para la mayoría de los bancos, la concesión de préstamos es la mayor fuente de riesgo de crédito, pero hay otras actividades que también pueden dar lugar a la aparición de este riesgo como, por ejemplo, las inversiones en instrumentos de renta fija, las operaciones con derivados, los préstamos con tarjeta de crédito o en descubierto, los depósitos en otras instituciones financieras, las garantías financieras o las liquidaciones.
}

A priori podríamos pensar que la solvencia crediticia constituye una cuestión crítica para todas las empresas y también para los hogares. Sin embargo, para los intermediarios financieros cuya actividad principal consiste en emitir contratos de pago contingente a sus clientes, la solvencia constituye la cuestión financiera central. La posibilidad de un incumplimiento futuro por parte de un intermediario respecto a sus obligaciones con los clientes puede reducir significativamente la eficiencia ex ante de dichos contratos y, en consecuencia, disminuir sustancialmente la eficacia de la principal función económica desempeñada por el intermediario. Por el contrario, la posibilidad de incumplimiento de la deuda en manos de inversores emitida por una empresa típica puede tener poco o ningún impacto sobre la efectividad de dicha empresa para desempeñar su función económica principal. Por ello, la gestión de este riesgo es casi siempre una actividad de primer orden para el funcionamiento del Sistema Financiero, mientras que, en general, no tiene por qué serlo para las empresas no financieras. 

De esta forma, el riesgo de crédito es la fuente principal de de riesgo, que puede llegar a amenazar la solidez de una institución financiera. Por tanto, ¿cuál sería la mejor manera de gestionarlo? Si consiste en la probabilidad de incumplimiento de una obligación contraída, la manera más natural de abordarlo sería: (i) evitando posiciones arriesgadas; y, (ii) mitigando posibles contingencias futuras. Precisamente en esta linea se dividen los instrumentos de gestión del riesgo de crédito (en dos momentos diferentes):
\begin{la}
    \item \textbf{Ex ante: Asimetrías informativas}. Antes de contraer cualquier obligación, las entidades realizan labores de screening con el fin de determinar la viabilidad y las condiciones de la obligación (notablemente, qué tipo de interés ha de remunerar el riesgo que asume la entidad, dado el coste de oportunidad que suponen otras opciones en el mercado). Uno de los objetivos, sino el principal, es tratar de superar los problemas de información asimétrica: en lo referente a la selección adversa, a través de la recopilación de información sobre la calidad crediticia del deudor; en lo referente al riesgo moral, a través del diseño de contratos que generen incentivos al futuro deudor para actuar de manera diligente; y,
    \item \textbf{Ex post: Control de las posiciones y dotación}. Por otro lado, la entidad también debe monitorizar constantemente sus propias posiciones y dotar adecuadamente las mismas ante el posible default de parte de su cartera crediticia.\footnote{Este segundo caso está estrechamente condicionado por la regulación microprudencial y, por tanto, se estudia en el tema 3B27.
    }
\end{la}

Veamos como se llevarían a cabo en la práctica.

\subsubsection{Ex ante: Calidad creditica y asimetrías informativas}
En un primer momento, podemos pensar que la cuantificación del riesgo crediticio sería equivalente a determinar cómo cuantificaría/valoraría el mercado dicha probabilidad, teniendo en cuenta todas las posibilidades de diversificación y cobertura proporcionadas por los mercados financieros.\footnote{Naturalmente, el nivel de riesgo debe ser adecuadamente ajustado por la presencia o no de garantías, saldos compensatorios y endosos entre otros. Pero también serán relevantes otras características del mercado crediticio: ¿Comparten los bancos información sobre sus acreedores? ¿Cómo se resuelve el proceso de quiebra? 

Es fácil comprobar que el rendimiento (aleatorio) de un préstamo dependerá de todas estas características.
} Sobre esta linea operan precisamente las primeras aproximaciones al respecto: cuantificar el riesgo crediticio plasmado en el \textit{pricing} de los préstamos (i.e. a la determinación del tipo de interés) a través de la prima de riesgo (BIERMAN y HASS, 1975; y, YAWITZ, 1977). 

Para desarrollar esta aproximación simplemente debemos suponer dos cosas: (i) que el riesgo de crédito es diversificable, de tal manera que lo único que importa es la probabilidad de incumplimiento individual; y, (ii) los métodos de calificación creditica (primos hermanos de las técnicas actuariales), permiten a los bancos estimar a priori esta probabilidad de incumplimiento, basándose en las características observables del solicitante de préstamo. 

Bajo este marco de hipótesis, podemos aplicar los mismos métodos de valoración de los instrumentos en los mercados de deuda (por su propia naturaleza) a las obligaciones financieras. Por tanto, consideremos que el precio del préstamo es igual a su valor presente: la suma descontada de los flujos futuros. 

¿Cómo podemos cuantificar el valor del riesgo en el préstamo? Si el precio del préstamo es igual a la suma descontada de los flujos futuros, la diferencia entre ese valor (sin riesgo) y ese mismo valor ajustando por la posibilidad de que dicho(s) flujo(s) no se materialicen nos dará (con riesgo), precismante, el \textbf{valor nominal del riesgo}. Analíticamente:

\eqblock{
\begin{aligned}
    \text{CR}_j=\underset{{\scriptsize{\text{Valor esperado del préstamo sin riesgo}}}}{\sum_{t_j=0}^{T_j}\frac{\text{Cupón}_j}{(1+\delta)^{t_j}}}-\underset{{\scriptsize{\text{Valor esperado del préstamo con riesgo}}}}{\sum_{t_j=0}^{T_j}\frac{\text{Cupón}_j}{(1+\delta)^{t_j}}\cdot \mathrm P_j(X=1)}\quad\Longrightarrow\quad\boxed{\text{CR}_j=\sum_{t_j=0}^{T_j}[1-\mathrm P_j(X=1)]\frac{\text{Cupón}_j}{(1+\delta)^{t_j}}}
\end{aligned}
}{}

Como vemos, la diferencia da el valor nominal del riesgo, ajustando por la posibilidad de que el deudor pague su cupón en cada periodo (se conoce a menudo como la probabilidad de no-default). ¿Qué podemos hacer con este resultado? Resulta obvio que esto no nos permite, a priori, obtener lo que nos interesa: la prima de riesgo del deudor. Por tanto, ¿cómo podemos pasar del coste del riesgo a la prima de riesgo?

El razonamiento que sigue es sencillo. Si la entidad pretende mitigar su riesgo de crédito lo más natural sería conseguir trasladar dicho riesgo al deudor. Entonces, ¿cómo compensa el banco esa pérdida esperada? Lo que debe hacer la entidad es incorporar dicha pérdida esperada al precio del préstamo mediante una \textbf{prima de riesgo} ($s$). Es decir, en lugar de prestar al tipo libre de riesgo $\delta$, prestará a un tipo superior: $r=\delta+s$. Analíticamente, deberemos resolver la siguiente igualdad:

\eqblock{
\begin{aligned}
    &\underset{\scriptsize{\text{Valor esperado del préstamo con riesgo}}}{\sum_{t_j=0}^{T_j}\frac{\text{Cupón}_j}{(1+\delta)^{t_j}}\cdot \mathrm P_j(X=1)}=\underset{\scriptsize{\text{Valor del préstamo con prima de riesgo}}}{\sum_{t_j=0}^{T_j}\frac{\text{Cupón}_j}{(1+r)^{t_j}}},\qquad \text{donde}\,\,\,r=\delta+s\\[2em]
    &\text{Si,}\,\,\,X\sim\text{Poisson}(\lambda)\qquad\Longrightarrow\qquad r=\delta+\lambda
\end{aligned}
}{}

Con esto la entidad garantiza que el establecimiento de una prima de riesgo cubra el riesgo de crédito (se lo traslada al deudor): \textbf{se fija un pricing tal que el valor del préstamo con prima es igual al valor esperado descontado con un tipo de interés sin riesgo}. Además, si aumimos que la probabilidad de no-incumplimiento sigue una distribución de POISSON de intensidad $\lambda$, es sencillo comprobar que la ecuación tiene una única solución: $r=\delta+\lambda$. Es decir, si la probabilidad de default del deudor sigue una distribución de POISSON con intensidad $\lambda$, la prima de riesgo $s$ es independiente de las características del préstamo: es igual a la intensidad del proceso de POISSON.\footnote{En otras palabras, si se considera una deuda determinada, de manera sencilla, el spread puede considerarse como la probabilidad instantánea de impago $\lambda$. Por ejemplo, un spread de 50 p.b. indica aproximadamente una probabilidad de fallo del $1 – e^{(-λ)} \sim λ = 0.5\%$ por año. 
}

A pesar de todas sus virtudes, este enfoque presenta grandes limitaciones. En particular, se fundamenta sobre hipótesis muy restrictivas: (i) que la probabilidad instantánea de impago es constante y exógena; (ii) que el riesgo crediticio es completamente diversificable; y, (iii) que en caso de impago, el valor residual de la empresa es cero. 

\paragraph{Extensiones: MERTON (1974)}
Para salvar estos problemas, MERTON (1974) desarrolla una extensión más completa de la valoración del riesgo basada en la valoración de opciones. Su razonamiento es sencillo. 

Considerando una empresa con un colateral $V(t)$ (valor de los activos de la empresa, por ejemplo), que planea pedir $D_0$ prestado en la fecha $t = 0$, y devolver $D$ en la fecha $t = T$. Al transcurso de la ventana temporal pueden ocurrir dos cosas en el momento $T$ del vencimiento del préstamo:
\begin{la}
    \item \textbf{Solvencia $\to$ Amortización}. Si $D < V(T)$, la empresa será solvente y capaz de amortizar el capital adeudado al banco, tal y como se comprometió; o,
    \item \textbf{Insolvencia $\to$ Liquidación}. Si $D > V(T)$, la empresa será insolvente, más bien estará en quiebra, y el Banco liquidará sus activos con el objetivo de recuperar la mayor parte posible del capital adeudado a la fecha. 
\end{la}

Es decir, las variables economícas relaventes de la obligación de pueden resumir de la siguiente manera: 

\eqblock{
\begin{aligned}
    \text{MERTON}:\qquad\underbrace{\min\{D,V(T)\}}_{\scriptsize{\text{Recuperación del Banco}}}\qquad \underbrace{\max\{0,V(T)-D\}}_{\scriptsize{\text{Patrimonio de la empresa en $T$}}}
\end{aligned}
}{}

Como se puede intuir, esta fórmulación sobre el del riesgo y rendimiento de los préstamos consiste en la traslación del Modelo BLACK-SCHOLES de valoración de opciones: un préstamo sería el equivalente al pago de una opción de compra (call option) sobre los activos de la empresa con un precio de ejercicio igual a $D$. En consecuencia, desde un punto de vista puramente financiero, otorgar un préstamo con riesgo a una empresa de responsabilidad limitada es similar a comprar los activos de la empresa y vender una opción de compra a sus accionistas. Esta formulación permite trabajar con una probabilidad endógeno, a diferencia del caso anterior, y muestra dependencia a otros factores interesantes desde una óptica económica como son la volatilidad del activo, el tiempo y/o el nivel de apalancamiento de la empresa deudora. Se alcanzan así los siguientes resultados:
\begin{lnum}
    \item El spread/prima de riesgo s aumenta proporcionalmente al grado de apalancamiento del deudor. Por lo tanto, las empresas más endeudadas pagarán tasas de interés más elevadas;
    \item El spread s aumenta con la volatilidad de los activos de la empresa deudora. Por lo tanto, las empresas con actividades de mayor riesgo pagarán tasas de interés más altas; y,
    \item La prima de riesgo global s*T aumenta con el vencimiento del préstamo. Por lo tanto, los préstamos a plazos más largos serán más caros. 
\end{lnum}

Por supuesto, este enfoque es algo simplista porque ignora, entre otros, posibles pagos intermedios, costes de liquidación, emplea un tipo de interés / tasa de refinanciación $r$ determinista o asume que la empresa no tiene más deudas. Sin embargo, sus resultados están alineados con la evidencia empírica y permite la evaluación aproximada del coste del riesgo crediticio de una forma muy razonable. 

\subsubsection{Ex post: Calidad de Balance y dotación frente a contingencias}
Ahora bien, aunque el riesgo de crédito o incumplimiento puede preverse y, en principio, trasladarse adecuadamente al cliente, existe la posibilidad de que el deudor se vea incapaz de cumplir con su obligación de pago. Esta situación dejaría a cualquier entidad financiera con una exposición notable de su balance por el deterioro de los créditos incobrables. 

Por razón de su función, esta exposición no puede ser ignorada. Imaginemos que pasaría si no tuvieramos que gestionar el riesgo de crédito ya presente en nuestro balance. La materialización progresiva y sostenida de estas péridad por deterioro podrían acabar provocando que la entidad se vea incapaz de cumplir ella misma sus oblgaciones a corto, situación de insolvencia que pondría en duda su viabilidad. Por esta razón, es necesario imponer restricciones o precauciones ante estas posibles contingencias. 

Se trata de instrumentos de gestión de crédito que entran en juego dentro del balance de la propia entidad y permiten que el Banco se aproviosone/dote lo suficiente como para poder afrontar sus obligaciones teniendo en cuenta las pérdidas que probablemente se materialicen. Es decir, consiste en reconocer previamente la posibilidad de que tenga lugar el default y aumentar la resiliencia ante impagos, minimizando la volatilidad. Distinguimos dos instrumentos que se ditinguen por su rareza y fuente de financiación.

\paragraph{Provisiones contables: Pérdidas esperadas}
El primero son las provisiones contables por pérdidas esperadas. Considerando en la actividad financiera siempre habrá un impago, lo normal es que estas pérdidas esperadas sean tratadas como un coste normal del negocio bancario y, por tanto, financiadas con ingresos corrientes mediante provisiones contables. (Excepcionalmente, con capital regulatorio si no hay suficiente para efectuar la provisión).

¿Cómo podemos cuantificar la provisión por pérdidas esperadas? Se puede cuantificar mediante diversos indicadores, de entre las presentamos las más comunes:\footnote{También mediante exposiciones al incumplimiento, pérdidas dado incumplimiento, provisiones, morosidad, concentración sectorial y activos ponderados por riesgo. 
}

\eqblock{
\begin{aligned}
    &{\textbf{EL}}=\overset{{\scriptsize\substack{\text{Probabilidad}\\\text{Incumplimiento}}}}{\text{PD}}\cdot \underset{{\scriptsize\substack{\text{Pérdida dado}\\\text{Impago}}}}{\text{LGD}}\cdot\overset{{\scriptsize\substack{\text{Exposición}\\\text{por Impago}}}}{\text{EAD}}={\text{PD}\cdot (1-\overset{{\scriptsize\substack{\text{Tasa}\\\text{Recuperación}}}}{\text{RR}})\cdot \text{EAD}}\\[1em]
\end{aligned}
}{}

Con una distribución de probabilidad de default conocida, las entidades deben provisionar las pérdidas esperadas teniendo en cuenta la probabilidad de que suceda (la media de la distribución, PD), la pérdida que generaría para la entidad (LGD) y la exposición por impago (EAD).\footnote{La cobertura ante la ECL, como se verá, no es a través de requerimientos de capital. Se incluye en este apartado por su relación estrecha con la UL, y porque, pese a que en este texto no se entrará en su estudio detallado (por limitaciones de espacio), se trata de una cuestión muy relevante en la literatura y la práctica de la regulación financiera. Para conocimiento del opositor interesado: la norma contable internacional IFRS9, entre otras cosas, introdujo un modelo de ECL para dotar el deterioro de activos financieros. Este modelo requiere que las entidades reconozcan pérdidas esperadas de crédito en instrumentos financieros basándose en información razonable y respaldada, que incluya datos históricos, condiciones económicas actuales e información prospectiva. La alternativa podría ser dotar los deterioros siguiendo una lógica de pérdida incurrida (como de hecho se hacía con anterioridad a IRFS9); sin embargo, ello genera balances menos estables. De este análisis prospectivo resulta una clasificación de activos financieros en diferentes “stages” con el propósito de contabilizar las pérdidas esperadas por deterioro de la calidad del crédito:
\begin{lnum}
    \item \textbf{STAGE 1}. Los activos financieros se clasifican aquí cuando se reconocen inicialmente. Las entidades reconocen una ECL de 12 meses en el activo financiero; es decir, la parte de la ECL a lo largo de la vida útil que se espera que resulte de eventos de incumplimiento en los próximos 12 meses;
    \item \textbf{STAGE 2}. Un activo financiero se reclasifica aquí si ha habido un aumento significativo en el riesgo de crédito desde el reconocimiento inicial, pero el riesgo de crédito aún no se ha vuelto \textit{inmanejable}. La ECL que reconocen las entidades reconocen para estos activos recogen toda la vida útil del activo: las entidades se tienen que dotar una pérdida. Son los denominados créditos en vigilancia especial;
    \item \textbf{STAGE 3}. Finalmente, los activos se reclasifican aquí cuando hay evidencia objetiva de deterioro en la fecha de presentación de informes (por ejemplo, dificultades financieras significativas del emisor o un incumplimiento de contrato). En estos casos, la PD se iguala a 1. Son los denominados créditos dudosos o non-performing loans (NPL).
\end{lnum}
}\textsuperscript{,}\footnote{Los elementos principales y sus cálculos pueden entenderse de la siguiente forma:
\begin{lnum}
    \item \textbf{Probabilidad de incumplimiento} (Probability of Default). La la probabilidad de que un deudor o contraparte incumpla en un horizonte temporal determinado. Refleja, por tanto, la probabilidad estimada de incumplimiento durante un período específico, como un año. Por ejemplo, si es del 2\%, significa que hay un 2\% de probabilidad de que el deudor incumpla dentro del período especificado;
    \item \textbf{Pérdida dado el impago} (Loss Given Default). Representa la proporción de una posición que no se espera recuperar en caso de incumplimiento. Suele expresarse como un porcentaje de la posición/cuantía en el momento del incumplimiento. Por ejemplo, si la exposición en el momento del incumplimiento es de 1,000,000€ y el LGD es del 40\%, la pérdida esperada dado el incumplimiento sería de $400,000$€;
    \item \textbf{Exposición en caso de impago} (Exposure at Default). Es la pérdida que un banco enfrentará en caso de incumplimiento. Para su cálculo debemos tener en cuenta la exposición actual en el momento del incumplimiento y puede incluir una exposición futura potencial. Es una medida de la pérdida máxima potencial en caso de incumplimiento. Por ejemplo, si un banco tiene una exposición de crédito de $500,000$€ a una contraparte en el momento del incumplimiento, la EAD es de $500,000$€;
    \item \textbf{Tasa de recuperación} (Recovery Rate). Es el porcentaje de la exposición que el banco espera recuperar en caso de que un deudor incurra en incumplimiento. Representa la proporción de la deuda pendiente que se puede recuperar a través de garantías u otros medios después de que se haya producido un incumplimiento. Por tanto, se expresa como un porcentaje y se calcula dividiendo el monto real de recuperación entre la exposición total en el momento del incumplimiento.
\end{lnum}
} Como es normal, la probabilidad por impago coincidirá con la media de la distribución de POISSON anterior, pues se trata del mismo componente. De hecho, es por eso que a menudo la provisión por pérdidas esperadas se cargan sobre el cliente incluyendolo implícitamente en el margen financiero que propone la entidad (la prima de riesgo). 

\paragraph{Capital regulatorio: Pérdidas inesperadas}
El segundo instrumento es la cobertura por pérdidas extremas, o pérdidas inesperadas. Considerando que la actividad financiera está sometida a incertidumbre radical, existe la posibilidad de que suceden eventos extremos que deterioren mucho el balance de la entidad. Ante esta escenario, y atendiendo a su escasa probabilidad, las entidades deben contar con el suficiente capital como para afrontar estas pérdidas extremas o inesperadas. 

¿Cómo podemos cuantificar esta cantidad? Lo más común es utilizar el método de ponderación de activos por riesgo:

\eqblock{
\begin{aligned}
    &
    \left.\begin{aligned}
        &(1)&&\text{RWA}_j=\left(\text{LGD}_j\cdot F(X_j)-\text{LGD}_j\cdot \text{PD}_j\right)\cdot \tau_j \cdot \overset{{\scriptsize\substack{\text{Ratio Capital}\\\text{Total}}}}{\text{RCT}}\cdot \underset{{\scriptsize\substack{\text{Ratio ajuste}\\\text{positivo}}}}{\Gamma}\\[1em]
        &(2)&&\text{RWA}=\sum_{j=1}^N\text{RWA}_j\cdot A_j
    \end{aligned}\right.
\end{aligned}
}{}

Para calcular la cuantía correspondiente a la pérdida inesperada, debemos tener en cuenta varios factores, de entre los cuales destaca el RWA, o Risk-Weighted Assets.\footnote{No explicaremos cómo se calcula. Sin embargo, constituye un indicador esencial introducido por el Comité de Basilea. La fórmula de cálculo está basada en el modelo ASFR (Asymptotic Single Risk Factor), desarrollado por GORDY (2003), y en las aportaciones de VASICEK, que permiten modelar las tasas de incumplimiento de una cartera formada por un número muy elevado de exposiciones de igual importe (carteras granulares) y con similares características en cuanto a riesgo.
} En un primer momento calculamos la ponderación de riesgo para cada activo y, una vez tenemos este indicador, calculararíamos el global de la identidadcon la segunda fórmula. Es sencillo.\footnote{
Brevemente, para obtener el peso de riesgo de cada activo, debemos tener en cuenta algunos de los elementos antes presentados pero obtenidos desde un punto de vista individualizado: la pérdida por impago (valor absoluto); la probabilidad de impago; la función de distribución para esa posición en particular ($F_j$); el parámetro $\tau_j$, que consiste en un ajuste por vencimiento (cuánto más cerca, menos probable); el parámetro $\text{RCT}$, la inversa de la ratio de capital total; y, el parámetro $\varepsilon$, que consiste en un factor de ajuste positivo, para que los RWA alcancen un nivel de riesgo adecuado.

A partir de ahí, los activos ponderados por riesgo se obtienen como una combinación lineal de todos los activos ponderados por su peso de riesgo RW.
} Sin embargo, el cálculo de las RW siguiendo la fórmula anterior puede resultar muy costoso para las entidades. Por ello, el marco de Basilea reconoce dos posibilidades: el enfoque estándar (SA) y el enfoque de los modelos internos (IRB). El primero consiste en que la propia regulación establece el RW para cada tipo de exposición. El segundo es empleado generalmente por las entidades de mayor tamaño, ya que disponen de medios suficientes para este análisis más granular. Sin embargo, existe evidencia de que el enfoque IRB da lugar a RWA más bajos, y por tanto a menores necesidades de emitir capital regulatorio.

De forma parecida a la fórmula anterior, pero individualizando por la distribución de cada activo, la pérdida inesperada o activos ponderados por riesgo se calcula teniendo en cuenta, la pérdida que generaría a la entidad (LGD), la probabilidad de default (PD) y la cola derecha de la función ($F(X_j)$). Además, suelen añadirse parámetros de precaución, como el ajuste positivo $\tau$ y la Ratio de Capital total de la entidad. 

En síntesis, partiendo de una distribución de probabilidad de default conocida, $X\sim \text{Poisson}(\lambda)$:

\imagenfit{Fig5.png}{Distribución de las pérdidas para un riesgo y descomposición: pérdida esperada y pérdida inesperada}{HULL (2018)}

La pérdida esperada (EL) sería la esperanza de la distribución, por tanto su probabilidad es igual al al valor acumulado: $1-\mathbb E[X]$. Por su parte, cubrir toda la cola de la derecha sería excesivamente costoso y dada la baja probabilidad de los sucesos de cola, parte de ésta no se incorpora al cálculo de la pérdida inesperada. Por eso, lo normal para calcular la probabilidad necesaria para el capital regulatorio por pérdidas inesperadas suele hacerse estableciendo un percentil y restándole la probabilidad de la pérdida esperada: $0.999-(1-\mathbb E[X])$.\footnote{
De hecho, se suele descomponer el area a la derecha de la pérdida esperada en dos: (i) el riesgo de crédito \textbf{poco probable}, $\text{EL}+\text{UL}$; y, (ii) el área a la derecha de la recta discontinua como riesgo de crédito \textbf{muy poco probable}. El riesgo muy poco probable equivale, por tanto, a la probabilidad de que el banco no pueda cumplir con sus obligaciones usando sus provisiones y capital regulatorio, por lo que constituye un parámetro importante que se denomina \textbf{pérdida por tensión} (stress loss).

En terminología financiera, el 100\% menos esta probabilidad se denomina valor en riesgo (VaR, value at risk, por sus siglas en ingles), que veremos más adelante. Si el capital se fija en función de la diferencia entre la EL y el VaR, la probabilidad de que el banco siga siendo solvente en un horizonte temporal dado es igual al nivel de confianza.
}

\subsubsection{Crisis financiera de 2008: Limitaciones}
\textcolor{red}{\textbf{OJO -- (ChatGPT)} Fuentes: Bernanke (2015), Mishkin (2022), Hull (2018), Gorton (2010), Financial Crisis Inquiry Commission (2011).

\textbf{Riesgo de crédito (3.1.3):} se infravaloró la probabilidad de impago de los activos subyacentes.}

La crisis financiera internacional iniciada en 2007 constituye uno de los ejemplos más ilustrativos de las limitaciones de los mecanismos tradicionales de gestión del riesgo de crédito. El origen inmediato de la crisis se encuentra en el mercado hipotecario estadounidense. Durante los años previos, el prolongado crecimiento del precio de la vivienda favoreció una expansión extraordinaria del crédito hipotecario, incluyendo préstamos concedidos a prestatarios con escasa capacidad de pago (subprime). La premisa implícita era que una eventual dificultad de pago podría resolverse mediante la refinanciación o la venta del inmueble a un precio superior.

Sin embargo, el problema no residía únicamente en la calidad de los deudores. El periodo de la desregulación que se inició en los 80's y culminó en 1999, con la revocación de la Glass-Steagall Act, permitió transformar cientos de miles de préstamos individuales en títulos negociables mediante procesos de titulización. De este modo, el riesgo de crédito dejó de permanecer en los balances de las entidades originadoras y pasó a distribuirse entre inversores de todo el mundo mediante instrumentos como los Mortgage-Backed Securities (MBS) y las Collateralized Debt Obligations (CDO).

En teoría, esta dispersión del riesgo debía aumentar la estabilidad del sistema. En la práctica, produjo el efecto contrario. La presencia de fuertes asimetrías informativas, poca prudencia y la complejidad de los instrumentos dificultó la valoración real de las exposiciones y debilitó los incentivos para mantener estándares rigurosos de concesión de crédito. Muchas entidades confiaron excesivamente en las calificaciones crediticias otorgadas por las agencias de rating y subestimaron la posibilidad de un deterioro simultáneo de un gran número de préstamos.

Cuando el precio de la vivienda comenzó a caer y aumentó la morosidad hipotecaria, las pérdidas dejaron de ser idiosincráticas y pasaron a estar fuertemente correlacionadas. Las probabilidades de impago estimadas a partir de datos históricos resultaron demasiado optimistas y las provisiones y requerimientos de capital acumulados fueron insuficientes para absorber el volumen de pérdidas efectivamente registrado.

La principal lección de la crisis fue que el riesgo de crédito no depende únicamente de la probabilidad de impago de cada prestatario individual. También depende de factores sistémicos, como la evolución del ciclo económico, la correlación entre incumplimientos y la calidad de los incentivos existentes en la cadena de intermediación financiera. En consecuencia, la gestión moderna del riesgo crediticio ha evolucionado hacia enfoques más prospectivos, incorporando pruebas de resistencia, análisis macroprudenciales y modelos basados en pérdidas esperadas a lo largo del ciclo económico.

\subsection{Riesgo de mercado: Deterioros de balance}
El segundo es riesgo que nos interesa es el de mercado: pérdidas en posiciones dentro o fuera de balance causadas por movimientos en precios de mercado, tipos de interés, tipos de cambio, acciones, commodities o spreads. Es decir, el riesgo que responde a la posibilidad de registrar pérdidas derivadas de variaciones en los precios de un activo mantenido por la entidad en su balance. 

Para medir el riesgo de mercado se recurre a diferentes métodos. Algunos de los más notables son el análisis o test de sensibilidad, el análisis por valores en riesgo (VaR y Expected Shortfall) o el análisis por correlaciones. Como puede tener un impacto material en la rentabilidad y viabilidad de una entidad, éstas deben mantener posiciones sólidad para cubrir este riesgo.\footnote{Sin embargo, se apunta que una parte fundamental de la gestión de riesgo de mercado incorpora otras metodologías como son:
\begin{la}
    \item \textbf{Análisis de sensibilidad}. Uno de los métodos más utlizados en el análisis por sensibilidades es el Modelo de BSM. En especial, el análisis de \textit{las griegas}, que miden la sensibilidad de un activo a variaciones en algunos de sus parámetros. [\authorfont{Ver Tema 3.B.25}];
    \item \textbf{Volatilidad}. Debemos reconocer que los mercados financieros no presentan únicamente riesgos cuantificables, sino que, a menudo, la incertidumbre radical hace que su desarrollo futuro sea imposible de medir con exactidud. En este sentido, es importante tratar de averiguar métodos que nos permitan cubrirnos ante la periodos de extrema volatilidad, y no solo de riesgo. Para ello existen numerosas metodologías, entre las más destacadas: (i) los EWMA (Exponentially Weighted Moving Average); (ii) los ARCH (Autoregressive Conditional Heteroscedasticity); y/o, (iii) los GARCH (Generalized Autoregressive Conditional Heteroscedasticity). Todos ellos son modelos estadísticos avanzados cuyo elemento más característico es poder tratar con volatilidades no constantes (a diferencia de los modelos dinámicos $\text{ARIMA}(p,d,q))$.
    \item \textbf{Correlaciones}. Otro método tradicional en la gestión de riesgo de mercado es mediante el análisis de las correlaciones. Como sabemos, este método es uno de los pilares dunamentales de la Teoría de gestión de carteras y busca neutralizar la volatilidad de los activos en portfolio mediante la diversificación. Sencillamente: si los cambios en las dos variables tienen una correlación positiva alta, la exposición total de la empresa será muy alta (aditividad); si tienen una correlación de cero, la exposición será menor pero excesiva; lo óptimo, por tanto, será que tengan una correlación negativa alta, permitiendo que la exposición se minimice al poder compensar la pérdida en un activo con las ganancias en otro; y,
    \item \textbf{Cópulas estadísticas}. Por último, y quizás el más complicado en términos estadísticos, es la gestión del riesgo mediante cópulas. En la teoría de la probabilidad y estadística, una cópula es una función de distribución multivariada cuyas distribuciones marginales para cada variable son distribuciones uniformes. En este contexto, las cópulas describen la estructura de dependencia entre variables aleatorias. Han sido usadas ampliamente en el área de finanzas y en aplicaciones de optimización de carteras. Por ejemplo, en ciencia actuarial, donde diversos procesos en los seguros de accidentes involucran pares de variables correlacionadas, se podría utilizar en la pérdida y los gastos que se le asignan a una reclamación en una empresa de seguros. En el área financiera son utilizadas en el modelado de los activos y la gestión de riesgos.
\end{la}
} 

\subsubsection{Value at Risk (VaR): ¿Cuánto puedo llegar a perder?}
Veamos como se gestionaría este riesgo mediante los modelos de valor en riesgo (VaR, por sus siglas en inglés, \textit{value at risk}). 

La pregunta que intentamos responder es: ¿qué pérdida máxima podemos enfrentar en un periodo $t$ con un intervalo de confianza de $1-\alpha$? Por tanto, esta metodología permite calcular el riesgo, acotando la pérdida máxima posible según las condiciones del mercado e indicando una referencia temporal para dicha perdida:\footnote{Actualmente, a nivel regulatorio, los bancos pueden utilizar dos metodologías para determinar su capital regulatorio por riesgo de mercado: el método estándar y el método basado en modelos internos, sujeto el uso de estos últimos a la aprobación y supervisión de las autoridades nacionales.
} es decir, el método VaR nos permite \textbf{determinar la probabilidad de sufrir la pérdida máxima en un determinado horizonte temporal}.

\imagenfit{Fig1.png}{Distribución de los rendimientos y VaR}{Apuntes ICEX-CECO. Tema 3.B.26 | HULL (2018)}

Aunque se puede utilizar para la gestión de de diferentes riesgos (mercado, liquidez y crédito), nos centraremos en cómo se utiliza en la gestión del riesgo de mercado. Para ello, deberemos especificar un nivel de confianza y un horizonte de riesgo (una semana, un mes, etc.). Disponemos de varios procedimientos para calcularlo, pero los más comunes son: (i) el procedimiento de varianzas y covarianzas (VCV); (ii) el procedimiento de simulación histórica (HM); o, (iii) la simulación de Montecarlo (MCS).

\paragraph{Métodos parámetricos: Estimación por varianzas y covarianzas (VCV)}
El procedimiento por varianzas y covarianzas es un método de estimación paramétrico que asume que los rendimientos de un activo son variables aleatorias con distribución normal (centrada en $0$). De esta forma, para calcular el VaR aplicaríamos la siguiente fórmula:

\eqblock{
\begin{aligned}
    \boxed{\text{VaR}_{1-\alpha,t}=\left(z_{1-\alpha}\,\,\sigma\,\,\sqrt{t}-\mu\,\,t\right)V_0}
\end{aligned}
}{VaR: Procedimiento por varianzas y covarianzas (VCV). Instrumentos de Gestión del riesgo}

Donde: (i) $z_{1-\alpha}$ es el estadístico asociado a ese nivel de confianza $1-\alpha$; (ii) $V_0$ es el valor del activo o posición; (iii) $\sigma$ es la desviación estándar de la distribución, en la unidad de medida diaria (ya que $\sigma_t=\sigma\sqrt{t}$); y, (iii) $t$ es el horizonte temporal sobre el cual se está estimando el VaR. 

No obstante, este método presenta series limitaciones: (i) suponer que los rendimientos financieros siguen una normal es excesivamente fuerte, ya que las series financieras suelen presentar lo que se conoce como colas gruesas; y, (ii) por otro lado, suponer que los rendimientos diarios son variables aleatorias independientes e idénticamente distribuidas también es fuerte puesto que los mercados financieros suelen presentar fenómenos de clusterización, periodos de concentrada volatilidad. 

\paragraph{Métodos no paramétricos (I): Simulación histórica (HS)}
El método de simulación histórica es un procedimiento no paramétrico que evita asumir una forma funcional para la distribución de los rendimientos. 

De esta forma, el enfoque de simulación histórica consiste en calcular el VaR a partir de datos históricos sobre el activo (o activos) en cuestión: asumimos que, al desconocer la verdadera distribución, la distribución histórica (de la muestra) es una buena aproximación a la distribución de probabilidad de rendimientos (de la población). Para calcular el VaR a un día utilizando este enfoque se realizarían los siguientes pasos: 
\begin{lnum}
    \item \textbf{Recopilación}. Recopilar datos diarios sobre el precio de cierre de la inversión durante un periodo de muestra y calcular los rendimientos diarios;
    \item \textbf{Tratamiento}. Construir una distribución (es decir, ordenarlos de menor a mayor rendimiento);
    \item \textbf{Definición}. Decidir el intervalo de confianza deseado; y,
    \item \textbf{Análisis}. Para calcular el VaR con un nivel de confianza del 99\%, se elige el valor de rendimiento del percentil deseado (sería el dato que dejara a su izquierda el 1\% de los datos en la muestra) y se multiplica por el valor de mercado de la inversión. 
\end{lnum} 

\imagenfit{Fig6.jpeg}{Tratamiento de los datos y simulación histórica: Rendimientos S\&P 500 a un día}{\href{https://www.youtube.com/watch?v=mJNAM6gbyiE}{Prytex Academy. Youtube}}

A pesar de sus ventajas, el método por simulación histórica es muy sensible a la muestra. En especial, al construirse con datos históricos, por muy larga que sea la serie temporal recogida, la muestra será condicional a las circunstancias económicas del momento.

\paragraph{Métodos no paramétricos (II): Simulación de MONTECARLO (MCS)}
El último de los métodos para obtener el VaR es mediante simulaciones de MONTECARLO. 

Las simulación de MONTECARLO es un método de estimación no paramétrico que permite al analista simular generar tantas simulaciones como quiera de las variables objetivos durante un periodo temporal definido. Es decir, asumiendo cualquier forma funcional para la distribución de los rendimientos, con la posibilidad de no ser idénticamente distribuidos y dependientes, permite generar un gran número de simulaciones que construyan per se una nueva distribución de probabilidad, en base a los outputs observados de la simulación. Una vez generadas las simulaciones, el VaR se estimaría siguiendo la misma lógica que la simulación histórica pero en base a los datos generados por la simulación.\footnote{La extensión de este método a toda la cartera de activos es también muy frecuente. En este caso, se seguiría el mismo proceso que para un activo aunque encontraríamo diferencias a la hora de generar la simulación:
\begin{la}
    \item \textbf{Rendimientos no correlacionados}. Si los activos del portfolio no están correlacionados, deberíamos simplememte realizar una simulación para cada instrumento objetivo; 
    \item \textbf{Rendimientos correlacionados}. No obstante, si los instrumentos si estuviesen correlacionados, la simulación de los retornos debe considerar tales covarianzas, lo cual complica el procedimiento de generación de procesos estocásticos. Lo normal en estas situaciones es recurrir a mecanismos de identificación de la estructura de la matriz de varianzas y covarianzas por medio de descomposición, por ejemplo mediante el método de descomposición de CHOLESKY.
\end{la}
}

\subsubsection{Crisis financiera de 2008: Limitaciones}
\textcolor{red}{\textbf{OJO -- (ChatGPT)} Fuentes: Hull (2018); Jorion (2007); Gorton (2010); Financial Crisis Inquiry Commission (2011); Basel Committee on Banking Supervision (2016).

\textbf{Riesgo de mercado (3.2.2):} se infravaloró cuánto podían caer de valor esos activos una vez que el mercado empezó a dudar de ellos y, sobre todo, se infravaloró la dependencia entre riesgos aparentemente diversificados.}

Nuevamente, la crisis financiera puso de manifiesto cómo y por qué los modelos tradicionales de gestión del riesgo de mercado presentaban importantes limitaciones cuando los mercados abandonaban condiciones normales de funcionamiento. En particular, reveló que herramientas ampliamente utilizadas como la presentada (VaR) podían proporcionar una falsa sensación de seguridad al estimar adecuadamente las pérdidas ordinarias, pero no los riesgos asociados a episodios de estrés sistémico. 

Durante los años previos a la crisis, numerosas entidades financieras mantenían exposiciones significativas a productos estructurados vinculados al mercado hipotecario estadounidense. La valoración de estos instrumentos descansaba sobre modelos estadísticos que utilizaban datos históricos para estimar probabilidades de impago, volatilidades y correlaciones entre activos, pero cuando comenzaron a aumentar los impagos hipotecarios, las pérdidas registradas superaron ampliamente las previstas por dichos modelos. La principal razón fue que los supuestos utilizados dejaron de cumplirse simultáneamente:
\begin{lnum}
    \item \textbf{Eventos extremos: Colas gruesas}. En primer lugar, las distribuciones de pérdidas mostraron comportamientos extremos incompatibles con los escenarios contemplados por muchos modelos. Como hemos visto, el VaR sería muy apropiado para saber cuánto perderíamos el, pongámos, 99\% de los días. Sin embargo, no nos dice cuánto podríamos llegar a perder ese 1\% restante, y este podría ser muy superior a la pérdida esperada el 99\% de los días. Este hecho es especialmente importante cuando los rendimientos no siguen una distribución normal y, por ejemplo, presentan colas gruesas (como lo hacen los rendimientos financieros). Esto evidenció una de las principales limitaciones del VaR: puede subestimar significativamente la exposición de una entidad a eventos de cola, precisamente aquellos que resultan más relevantes durante una crisis financiera. Por esta razón una medida complementaria que se ha impuesto desde enctonces es el \textbf{expected shortfall}, que nos dirá cuánto podemos esperar perder fuera de ese intervalo de confianza: ${\text{ES}_{\alpha }(X)=-\mathbb E[X\,|\,X<x_{\alpha }]}$; y,
    \item \textbf{Correlaciones: Clusterización}. En segundo lugar, las correlaciones entre activos aumentaron bruscamente: riesgos que parecían diversificados comenzaron a deteriorarse de forma conjunta, reduciendo drásticamente los beneficios esperados de la diversificación. Hemos visto como el VaR es una medida de riesgo subaditiva bajo supuestos de normalidad (es decir, $\rho(X+Y)\leq \rho(X)+\rho(Y)$). Por lo tanto, una cartera adecuadamente diversificada no debería parecer más arriesgada que la suma de sus partes. Sin embargo, cuando se experimentan fenómenos de clusterización, cambios bruscos en los niveles de correlación entre activos, el VaR puede llegar a \textbf{castigar la diversificación}, lo que es conceptualmente problemático y no había sido previamente contemplado.
\end{lnum}

La situación se agravó aún más en los mercados de derivados de crédito. Los Credit Default Swaps (CDS), diseñados originalmente como instrumentos de cobertura frente al riesgo de impago, generaron una compleja red de exposiciones cruzadas entre instituciones financieras. Muchas entidades consideraban cubiertas sus posiciones porque habían transferido el riesgo mediante contratos CDS. Sin embargo, esta protección dependía de la solvencia de las contrapartes que habían vendido dichas coberturas.\footnote{El caso más paradigmático fue el de AIG, una de las mayores aseguradoras del mundo. A través de su división financiera había vendido protección mediante CDS sobre grandes volúmenes de deuda titulizada sin mantener reservas suficientes para afrontar un deterioro simultáneo de los activos subyacentes. Cuando las pérdidas comenzaron a materializarse, surgieron dudas sobre su capacidad de pago, poniendo en riesgo la validez de numerosas coberturas existentes en el sistema financiero internacional.
}

Esto materializó otra de las grandes limitaciones del VaR: \textbf{¿Qué pasa si no podemos vender?}. El marco VaR ignora la posibilidad de enfrentarse a mercados ilíquidos. Pensemos que la interpretación es: ¿cuál es la pérdida máxima esperada durante los próximos 10 días con un nivel de confianza del 99\%? Para que esa pregunta tenga sentido como medida del riesgo real, se suele asumir que durante un marco temporal (p.e. 10 días) el banco puede reducir, cerrar o cubrir sus posiciones a precios cercanos a los observados en el mercado. Pero, ¿qué pasa si no? Por ejemplo, si tenemos un bono cotiza a 100 y el VaR estima que, con un 99\% de confianza, no perderás más de 5, nuestro valor en riesgo será de 5. Pero si estalla una crisis y nadie quiere comprar ese bono, quizá sólo puedas venderlo a 85, es decir, nuestra pérdida efectiva será 15, no 5. El problema no es que el VaR esté mal calculado respecto a la distribución de precios utilizada. El problema es que la distribución de pérdidas relevante debería incluir también el coste de deshacer la posición en un mercado ilíquido. Esto ocurrió en 2008, cuando muchos activos dejaron de ser líquidos, habiendo muy pocos compradores, grandes diferenciales bid-ask que se dispararon y resultaba imposible vender grandes posiciones sin mover el precio fuertemente en contra del vendedor.

En su conjunto, estas debilidades del modelo VaR no se habían valorado adecuadamente desde el punto de vista regulatorio.\footnote{Aunque sí que habían sido ampliamente tratadas y puestas de manifiesto desde el ámbito académico (DANIELSSON, 2002; DANIELSSON et al, 2004; DANIELSON \& ZIGRAND, 2006)
} Por tanto, si bien el VaR sigue siendo un componente esencial en la gestión del riesgo de mercado (por ejemplo, al determinar la cuantía de capital regulatorio), se debe complementar con otras medidas adicionales que mejore la gestión de riesgos, entre otros:\footnote{También se prestó especial atención, dentro de la categoría de los instrumentos de negociación de crédito, a las titulizaciones.
} (i) el expected shortfall; (ii) modelos internos que introduzcan, por ejempol, el VaR estresado (S-VaR), que hemos visto antes; (iii) el capital de riesgo incremental (IRC, por sus siglas en inglés, Incremental Risk Charge); o, (iv) la Medida de riesgo Global (CRM, Comprehensive Risk Measure, por sus siglas en ingles.

\begin{specialblock}{Extensiones - Refinamiento de los métodos VaR}
    \textbf{Expected Shortfall (ES)}. El expected shortfall, también denominada en ocasiones riesgo condicional, expectativa de cola condicional o pérdida esperada en la cola, consistiría en preguntarse: \textit{Si las cosas realmente se ponen mal, ¿cuál es la pérdida esperada?}. Por tanto, como podemos intuir, se trata de calcular el valor esperado de la distribución condicionada a que se encuentra dentro del intervalo del VaR: donde las cosas están muy mal. Para su cálculo, al igual que el VaR, necesitamos únicamente definir dos parámetros: el horizonte temporal ($t$) y el nivel de confianza ($\alpha$). Si suponemos que $R_t\sim N(\mu_t,\sigma_t)$, entonces $\text{ES}_{1-\alpha}=-\mathbb E[R_t \mid R_t > \text{VaR}_{1-\alpha}]$. O, de forma directa (si sigue una normal): $\text{ES}_{1-\alpha,t}=\left(-\mu t+\sigma\sqrt{t}\frac{\varphi(z_{1-\alpha})}{\alpha}\right)V_0$, donde $\varphi(\cdot)$ es la función de densidad.
    
    \textbf{VaR estresado} (S-VaR). El VaR estresado o VaR en situaciones de tensión, utiliza los rendimientos de un año malo, en un intento de capturar valores extremos y tener más en cuenta el riesgo de cola. Por ello, el SVaR suele ser mayor que el VaR, arrojaría una cantidad menor de activos ponderados por riesgo y, en consecuencia, mayores requerimentos de capital reglamentario.
    
    \textbf{Marco de Capital de Riesgo Incremental} (Incremental Risk Charge, IRC). El VaR tradicional se centra en las fluctuaciones de precios a corto plazo, pero apenas captura las pérdidas asociadas al deterioro de la calidad crediticia de los emisores. El IRC surge para cubrir esta carencia. En esencia, responde a la pregunta: ¿qué pérdidas podría sufrir la cartera si, durante el próximo año, algunos emisores incumplen sus obligaciones o ven deteriorada significativamente su calificación crediticia? A diferencia del VaR, el IRC utiliza un horizonte temporal de un año y un nivel de confianza del 99,9\%, por lo que está diseñado para capturar eventos poco frecuentes pero potencialmente muy graves. Además, incorpora explícitamente el concepto de horizonte de liquidez, es decir, el tiempo necesario para reducir o cubrir una posición en condiciones de tensión en los mercados. De este modo, el IRC extiende la medición del riesgo desde las variaciones de precios a corto plazo hacia las pérdidas derivadas del riesgo de crédito y de la falta de liquidez.
    
    \textbf{Medida Global de Riesgo} (Comprehensive Risk Measure, CRM). La crisis financiera puso de manifiesto que determinados productos titulizados, especialmente los más complejos, estaban expuestos a riesgos que no eran capturados adecuadamente por el VaR. La CRM fue diseñada específicamente para medir el riesgo de estas posiciones. Mientras que el IRC se centra principalmente en el riesgo de impago y en las migraciones de calificación crediticia, la CRM incorpora además otros factores relevantes para los productos titulizados, como la posibilidad de impagos múltiples, la variación de los diferenciales de crédito (credit spreads), los riesgos de base, los cambios en las correlaciones implícitas o la incertidumbre sobre las tasas de recuperación tras un incumplimiento. Al igual que el IRC, utiliza un horizonte temporal de un año, lo que permite evaluar pérdidas derivadas de escenarios de estrés prolongados que difícilmente quedarían reflejados en un VaR convencional de diez días.

    En esencia, una forma simple de entender esta sección sería:
    \begin{lnum}
        \item \textbf{VaR}: ¿cuánto puedo perder en condiciones normales de mercado?
        \item \textbf{S-VaR}: ¿cuánto podría perder si se repitiese un periodo histórico de tensión?
        \item \textbf{ES}: si supero el VaR, ¿cuánto perderé de media?
        \item \textbf{IRC}: ¿qué ocurre si empeora la calidad crediticia de los emisores durante el próximo año?
        \item \textbf{CRM}: ¿qué ocurre si los riesgos específicos de los productos titulizados se materializan en un escenario de estrés?
    \end{lnum}

    De esta forma, es más fácil entender por qué Basilea fue añadiendo capas sucesivas al VaR tras la crisis de 2007-08.
\end{specialblock}

\subsection{Riesgo de liquidez: Desfase temporal del Balance}
El último riesgo que vamos a tratar es el \textbf{riesgo de liquidez} y sus diferentes instrumentos de gestión.

El riesgo de liquidez puede definirse como la \textbf{incapacidad de financiar aumentos de activos o atender obligaciones vencidas sin incurrir en pérdidas inaceptables}, por no disponer de activos líquidos o por no poder acceder a los mercados para obtener financiación a un precio razonable.\footnote{Cuando se alude al concepto de iliquidez es importante distinguirlo de la insolvencia. El primero se refiere a una situación en la que los activos de una entidad son mayores que los pasivos, pero el banco no tiene suficientes activos líquidos como para asumir todos sus compromisos de pago, ni puede vender sus activos a precios razonables, para poder hacerlo en el corto plazo. Por el contrario, la insolvencia implica que el valor de los activos es inferior a los compromisos de pago son mayores que se poseen.
} Por esta razón, este riesgo suele descomponerse en dos factores: el riesgo de liquidez en la financiación, cuando una entidad no es capaz de obtener el efectivo que requiere, y/o el riesgo de liquidez de mercado, la posibilidad de ser incapaz de revertir una posición en el mercado de valores sin afectar a las condiciones de éste (precio/spread bid/ask). También, por esta misma razón, puede traer por causa motivos sistémicos o particulares de la entidad.

Como vemos, este riesgo es especialmente característico de las entidades que operan en el Sector Financiero, por ello Basilea subraya especialmente la liquidez como condición de supervivencia aun para entidades solventes. Veamos algunos de los instrumentos de gestión más comunes.

\begin{specialblock}{Caso particular - Suspensión de la convertibilidad}
    Un instrumento de gestión común en cualquier entidad financiera es la \textbf{suspensión de la convertibilidad}. Evidentemente, cuando la estructura temporal del balance pone en riesgo la solvencia de la entidad debido a una crisis de liquidez, la suspensión de la convertibilidad de los depósitos o liquidación de activos es una forma drástica pero efectiva de mitigar este riesgo:
    \begin{la}
        \item \textbf{Sistema Financiero No Bancario}. En el Sistema Financiero No Bancario, este instrumento de gestión se conoce con el nombre genérico de \textbf{suspensión de los reembolsos}. Supone mantener las aportaciones de los partícipes indisponibles para estos, por un período de tiempo. No es infrecuente que la suspensión desencadene en una liquidación y disolución del fondo: los partícipes habrían quedado \textit{atrapados};
        \item \textbf{Sistema Financiero Bancario}. En el Sistema Financiero Bancario, esta medida se conoce bajo el nombre genérico de \textbf{suspensión de la convertibilidad}. Este mecanismo fue utilizado por los bancos antes de que existieran los seguros de depósitos gubernamentales (corralito). DIAMOND y DIRVIG señalan, como una extensión de su modelo original, que estableciendo una política donde se permite solo retirar una fracción ($\theta$) de los depósitos el Sistema Bancario permanece estable, pues se fuerza el equilibrio de NASH deseado. Ahora bien, esto funcionará siempre y cuando no exista incertidumbre sobre la proporción de depositantes que desean consumir anticipadamente su ahorro; si existiera, este tipo de medidas no serían efectivas para evitar el equilibrio de NASH no óptimo puesto que no se podría fijar de manera óptima la fracción disponible que hiciese óptima la intermediación bancaria para todos los agentes. 
    \end{la}

    Una forma suavizada de esta medida, es la liquidación o convertibilidad bajo descuento (perdiendo una fracción). Si las entidades son capaces de establecer un descuento tal que minimice la utilidad de la retirada anticipada, entonces el riesgo de liquidez quedaría también mitigado. Otra cuestión a tener en cuenta es los posibles costes que se deriven de la suspensión de la convertibilidad. No solo en términos económicos, sino también sobre la credibilidad, inestabilidad institucional de ningún tipo, etc. 
\end{specialblock}

\subsubsection{Instrumentos de Gestión de Liquidez: SFNB}
Empezando por los instrumentos más característicos del Sector Financiero No Bancario (SFNB).\footnote{Respecto a las causas particulares que pueden ocasionar este problema, la literatura aún es muy recientes. Sin embargo, algunos (\textcolor{red}{BlackRock, 2020}) distinguen, entre otros, los siguientes factores como posibles fuentes de riesgo de iliquidez: 
\begin{la}
    \item \textbf{Instrumentos derivados (OTC)}. Los derivados OTC, en especial, las margin calls que requieren las cámaras de contrapartida para compensar las posiciones en derivados;
    \item \textbf{Hipersensibilidad}. La hipersensibilidad de los precios de las participaciones de los ETF (fondos de inversión cotizados) y la retirada de los intermediarios necesarios para la correcta formación de precios en estos activos;
    \item \textbf{Primas de riesgo sistémico}. Los movimientos en el precio del activo libre de riesgo, a través de los ajustes de cartera asociados a éstos. En especial, los realizados por los fondos del mercado monetario, MMF;
    \item \textbf{Señalización}. Las calificaciones crediticias de las agencias de rating, que en ocasiones dan lugar a profecías auto-cumplidas, al tener lugar antes la caída de rating que la pérdida de confianza de los acreedores;
    \item \textbf{Capacidad operativa}. Capacidad operacional, incluyendo la de las infraestructuras de mercado; y,
    \item \textbf{Crisis de liquidez}. Los reembolsos masivos en los fondos de inversión.
\end{la} 
}

Imaginemos una situación en la que se experimentan peticiones de reembolso masivo en los fondos de inversión. ¿Qué situación puede acabar provocando este movimiento? ¿En qué consisten los reembolsos masivos? ¿Por qué es frágil la estructura y el modelo de negocio de los fondos de inversión?\footnote{Los fondos de inversión agrupan capital de diverso origen y lo canalizan hacia destinos concordantes con las preferencias de sus depositantes, generando rentabilidad y sirviendo a la diversificación de su exposición. Como vehículos de inversión colectiva, los fondos de inversión emiten participaciones y utilizan los ingresos de suscripción para financiar la compra de una cartera de activos que el gestor del fondo ajusta con el tiempo de acuerdo con la estrategia de inversión declarada. 
}

Cuando los inversores suscriben o reembolsan (compran o venden) participaciones de un fondo, pagan/reciben el precio por unidad, que se basa en el valor liquidativo de los activos del fondo (NAV, \textit{net asset value}).\footnote{De manera simplificada, el NAV se obtiene de dividir el valor de la cartera del fondo por el número de participaciones. Es decir, puede ser mayor o menor a la cantidad aportada inicialmente por el partícipe, y varía de manera continua.

En esto difieren de otros vehículos para canalizar el ahorro, como los depósitos bancarios o las participaciones en ciertos fondos del mercado monetario (MMF): los fondos de inversión no garantizan que las inversiones sean redimibles al valor nominal. Por tanto, pese a compartir con éstos algunos elementos comunes a toda actividad de transformación de vencimientos o de reasignación de riesgos, la modalidad de fragilidad difiere de un pánico bancario o de un reembolso masivo en MMF.
} Por tanto, \textbf{existe una diferencia entre el plazo de reembolso de algunos fondos de inversión} (para algunos, diario) \textbf{y el tiempo que puede llevar liquidar sus tenencias de manera ordenada para satisfacer las solicitudes de reembolso} (especialmente en el caso de los activos menos líquidos, como ciertos bonos corporativos o posiciones en real estate). Surge así un first-mover-advantage: cada partícipe tiene incentivos para reembolsar antes que los demás: (i) porque según el fondo recibe peticiones de reembolso y por tanto liquida su activo, el precio de los activos cae, disminuyendo el valor de las participaciones de los inversores que se quedan; y, (ii) porque el fondo tiende a liquidar primero los activos más líquidos, por lo que según va liquidando los menos líquidos después, las caídas en el NAV son cada vez mayores.\footnote{Es fácil comprobar que, en caso de un episodio de tensión, esta lógica genera efectos retroactivos circulares y potencialmente problemas de expectativas auto-cumplidas.
} 

Ante este escenario, ¿qué herramientas permiten atajar este problema? Evidentemente, esto dependerá de la jurisdicción en que nos encontremos, puesto que no todos los países permiten las mismas decisiones o comportamiento, pero podemos distinguir esencialmente dos tipos/grupos.

\paragraph{Grupo LMT (I): Penalizaciones sobre NAV}
Dentro del primer grupo de instrumentos de gestión de riesgo encontraríamos los que establecen algún tipo de penalización sobre la liquidación anticipada de los activos, normalmente incidiendo sobre el net asset value (NAV).

Por ejemplo, la \textbf{fijación de precios swing}, que consiste en ajustar el NAV para trasladar al menos parte de los costes de transacción del reembolso de las participaciones a los partícipes que solicitan dichos reembolsos. La fijación de precios swing tiene el potencial de reducir el first-mover-advantage al asegurar que todo o parte del coste de satisfacer los reembolsos recae en los inversores que realizan esos reembolsos, en lugar de los inversores que permanecen en el fondo. De manera estilizada, el swing price ($P^s$) se puede expresar como función del factor swing ($f^s$), los costes de transacción ($C$) y el volumen neto de reembolsos ($R$):

\eqblock{
\begin{aligned}
    \boxed{P^s=NAV\cdot f^s}\qquad f^s=\frac{\mathbb E[C]}{R+\mathbb E[C]}<1
\end{aligned}
}{Riesgo de liquidez (SFNB): Fijación de precios swing. Instrumentos de Gestión de riesgo}

Como es de esperar, la estimación de los costes es la parte más compleja ya que incorpora numerosos componentes, muy heterogéneos y a menudo con distribuciones de probabilidad correlacionadas (por ejemplo, impuestos y gravámenes, comisiones de brokerage, spreads bid-ask, y el impacto en el precio de mercado). Además del factor de swing, el gestor habrá de determinar el umbral a partir del cual se activa éste, si se establece.

Otro instrumento en esta misma linea sería la fijación de \textbf{comisiones de reembolso y gravámenes contra la dilución} (ADL, Anti-Dilution Levies). Con un espíritu similar al anterior en tanto que actúa sobre la variable precio y pretende desincentivar los reembolsos, la calibración será fundamental ya que debe fijarse una comisión suficiente para que el partícipe prefiera mantener su posición en el fondo antes que liquidarla. Evidentemente, esto resulta difícil de estimar dada la heterogeneidad de expectativas de los distintos inversores, así como sus diferentes preferencias (aversión al riesgo, preferencia por la liquidez). Por ejemplo, la obtención de la comisión óptima se podría modelizar a través de una función cóncava, que refleje la existencia de un trade-off: en un primer momento, mayores comisiones afectarán negativamente a la demanda de reembolsos (al encarecerlas) pero, a partir de un punto, comisiones excesivamente elevadas incrementarán la demanda de reembolsos. Esto podría deberse o bien a que la pérdida de liquidez de la participación la hace menos atractiva para el inversor, o a que comisiones muy elevadas pueden interpretarse como señalizadora de serios problemas dentro del fondo.

Por último, tedríamos la \textbf{liquidación en especie} que permite al gestor del fondo reembolsar al partícipe la cuantía en valores en cartera. Por tanto: (i) implica menores costes de transacción para el fondo, así como una mayor preservación de la liquidez del activo del fondo; (ii) el partícipe soporta el riesgo de mercado incluso tras la salida (lo que puede tener un efecto desincentivador para solicitar reembolsos); y, (iii) no se requiere que el fondo liquide sus activos en el mercado, lo que reduce el riesgo de ventas masivas y de caída de los precios. Es decir, el fondo no contribuye a los efectos de segunda ronda en el mercado de activos. Sin embargo, esta LMT presenta dificultades, como encontrar una subcartera que sea representativa del riesgo y liquidez de la cartera agregada del fondo.

\paragraph{Grupo LMTs (II): Desfase temporal de la liquidación}
Otra gama de instrumentos de gestión serán aquellos que tienen por objeto establecer un desfase temporal entre las órdenes de liquidación anticipada y su materialización.

En este sentido, el más común sería el \textbf{cambio en la frecuencia de operaciones del fondo}, permitiendo al gestor del fondo calibrar la frecuencia de publicación del NAV así como establecer ventanas de reembolsos con menor periodicidad en caso de tensión en los mercados:

\imagenfit{Fig6.png}{Efecto de la reducción de la liquidez en los reembolsos de un fondo}{Banque de France (2020)}

Una frecuencia menor en el cálculo del NAV permite: (i) proporcionar evaluaciones más sólidas del valor de la cartera del fondo, al recopilar más información sobre las transacciones realmente ejecutadas en los activos subyacentes; y, (ii) limitar el efecto first-mover-advantage, acumulando órdenes de compra y venta entre dos NAV (ya que se liquidarán las participaciones al valor al segundo NAV del rango, que es el más bajo).

Por otro lado tenemos instrumentos de carácter más jurídico, como el establecimiento de \textbf{períodos de aviso y temporary gates}. La lógica es la misma: se busca introducir una rigidez en los precios (no se permite la liquidación continuada de las participaciones), para evitar la formación de sendas explosivas. Sin embargo, en este caso el periodo de preaviso de reembolso supone el establecimiento de un lapso de tiempo requerido entre la solicitud de reembolso y su pago efectivo, con el fin de dar al gestor del fondo tiempo para obtener la liquidez necesaria, que se pagará al valor neto de los activos calculado en el día de la ejecución de la orden (no en el de la solicitud). Cuando el mercado no es muy líquido, el preaviso permite al gestor del fondo distribuir la orden a lo largo de varios días, para no distorsionar la cartera en detrimento de los restantes titulares de participaciones. De esta manera, el gestor del fondo puede anticipar las transacciones futuras y tener más flexibilidad para cumplir con las solicitudes de reembolso en condiciones óptimas. Las temporary gates son un caso extremo en el que, una vez superado un umbral de reembolsos determinado (para los fondos con cálculo de NAV diario, generalmente, el 5\% de los activos netos totales), el fondo aplaza los reembolsos. De este modo, se garantiza que la salida de los inversores es escalonada y, a diferencia de los períodos de preaviso estándares, el fondo cuenta con discrecionalidad para ordenar las peticiones y elegir los intervalos para ejecutarlas.

\textbf{LMT (VI): Side pockets}. En caso de que el episodio de tensión afecte a un segmento de mercado específico, los gestores pueden decidir segregar estos activos en cuestión,\footnote{Un precedente famoso fue empleado tras la revelación del fraude de Madoff por parte de fondos que tenían valores litigiosos.
} dividiéndose en dos: (i) un fondo \textit{saludable} o réplica, que asume toda la cartera del fondo original excepto los activos problemáticos; y, (ii) un fondo de \textbf{cancelación de deudas}, gestionado hacia la extinción, pero que puede esperar un aumento en la valoración de las inversiones afectadas que permita la compensación de los titulares de participaciones. Cada inversor en el fondo inicial recibe una asignación proporcional tanto del fondo réplica como del compartimento problemático. El fondo réplica luego puede seguir suscribiéndose o reembolsándose, a diferencia de las participaciones del fondo de cancelación de deudas (sin suscripción ni reembolso). A diferencia de la mayoría de LMT presentadas hasta ahora, en la mayoría de jurisdicciones el gestor puede implementar side-pockets de ser necesario sin necesidad de una habilitación expresa en el folleto informativo del fondo.

\subsubsection{Instrumentos de Gestión de Liquidez: Sistema Bancario}
Veamos ahora los intrumentos de gestión de liquidez disponibles para el Sistema Bancario. 

Imaginemos que una situación en la que un rápido deterioro del balance mermase la confianza de los depositantes y estos desearan retirar, en masa, sus depósitos. Como es normal, dado el sistema de banca fraccionaria, una entidad, o en este caso todo el sistema de intermediación bancaria, podría no tener suficiente liquidez como para hacer frente a tal volumen de liquidaciones, cuestión que fue estudiada a raíz del modelo propuesto por DIAMOND y DYBVIG (1983), que veremos más en detalle en la sección correspondiente.

¿Cómo podemos gestionar este riesgo? Es decir, ¿cómo podemos prepararnos ante esta posible contingencia? Existen varias soluciones.\footnote{Existe la opción de optar por un modelo de \textit{narrow banking}, como el propuesto por FISHER en el Plan Chicago frente a la Gran Depresión. Un banco en este modelo sería una institución en cuyo activo sólo figuraría un conjunto muy limitado de activos, todos ellos de muy bajo riesgo y altamente líquidos (un ejemplo serían letras del Tesoro). En pocas palabras, este modelo supone separar el negocio de la captación de depósitos del de la provisión de préstamos a la economía. Éste segundo sería realizado por otros intermediarios financieros, no habilitados para la captación de depósitos. Como ventaja, este modelo hace casi irrealizable la insolvencia para la entidad, por lo que los depósitos estarían seguros. Sin embargo, implicaría una menor oferta de capital para la economía, acabar con el negocio de transformación de plazos bancario y potencialmente un riesgo para la transmisión de la Política Monetaria. 
} 

\paragraph{Seguro de depósitos: Fondo de Garantía de Depósitos}
En primer lugar, y como se puede intuir, la manera más simple de gestionar este riesgo es a través de un \textbf{seguro}. Es decir, constituir un fondo o seguro de depósito que, en base a las aportaciones de la teoría de seguros (Teoría de la Agencia), permite \textbf{agrupar el riesgo} (\textit{risk pooling}) y gestionarlo adecuadamente para todas las entidades del Sistema a partir de un fondo común.\footnote{Una nota sobre terminología. La gestión de riesgos puede realizarse internamente o conjuntamente entre varios agentes. En lo que respecta a la primera vía, existen numerosas formas (varias de ellas presentadas en el tema 3B26, y otros del temario, incluido éste), entre otras: la diversificación de las exposiciones, la gestión de las griegas (por ejemplo, a través de la contratación de derivados), modelos de cobertura basados en los primeros y segundos momentos de las distribuciones de pérdidas (por ejemplo, el VaR o el ECL), modelos de predicción basados en las varianzas (por ejemplo, los modelos GARCH o los EWMA), la gestión de la duración, etc. Nos interesa no obstante la segunda: el caso en el que varios agentes colaboran para compartir parte o la totalidad de sus riesgos: \textit{risk sharing}. La idea central del risk sharing es que, bajo una serie de supuestos, el riesgo no es linealmente aditivo, lo que permite reducir el riesgo total. Aunque la forma más evidente es aquella en la que cooperan un agente averso al riesgo con uno neutral, existen, presuntamente, ganancias paretianas en términos de utilidad incluso en aquellos casos en que todos los agentes son aversos al riesgo (Teorema de ARROW-LIND). 
} 

La lógica sería sencilla: cada banco debería pagar periódicamente por una cobertura igual a sus depósitos de tal manera que, de tener lugar un evento crítico, se pudieran utilizar los fondos acumulados para garantizar la liquidez suficiente a la entidad para poder gestionar las peticiones de reembolso.

Presentemos un modelo simple de cómo funcionaría. Supongamos un Sistema Bancario compuesto por $n$ entidades, todas idénticas en cuanto al riesgo y el activo inicial ($A_0$), y con un escenario posible que suponga una pérdida de volumen $L$ que generaría la quiebra de la entidad. Asumimos también que la probabilidad $p$ de que acaezca este suceso de pérdida en un período determinado es idéntico para todas ellas e independiente. De este modo, al final de cada período, el valor esperado del activo y la varianza para cada entidad serán respectivamente:

\eqblock{
\begin{aligned}
    &\mathbb E[A_1]=p(A_0-L)+(1-p)A_0=A_0-p\,L\\
    &\mathrm{Var}(A_1)=\mathbb E[(A_1-\mathbb E[A_1])^2]
\end{aligned}
}{}

Ahora bien, aunque la constitución de un fondo común no altera la pérdida esperada de cada entidad, sí reduce significativamente la volatilidad de las pérdidas soportadas por el conjunto. ¿Por qué? Porque cuando se agregan los riesgos de múltiples entidades independientes, la pérdida esperada total del sistema es simplemente la suma de las pérdidas esperadas individuales ($n\cdot L$). Por tanto, al repartir posteriormente esa pérdida agregada entre todas las entidades, la pérdida esperada por participante permanece inalterada y coincide con la que tendría cada entidad de forma aislada ($\frac{n\cdot L}{n}\to L$). Sin embargo, la dispersión de los resultados no se comporta de la misma manera. Mientras que la esperanza matemática crece proporcionalmente con el número de entidades agregadas, la volatilidad relativa de la pérdida media disminuye conforme aumenta el tamaño del grupo. Esto se debe a que las fluctuaciones idiosincráticas de cada entidad tienden a compensarse entre sí cuando los riesgos son independientes. Como consecuencia, las pérdidas efectivamente observadas se concentran cada vez más alrededor de su valor esperado a medida que aumenta el número de participantes.\footnote{

Ahora bien, la agregación de riesgos no modifica la pérdida esperada por entidad, pero sí reduce la volatilidad de dicha pérdida. Sea $(X_i)$ la pérdida de la entidad ($i$), con: esperanza, $\mathbb E[X_i]=pL$; y, varianza $\mathrm{Var}(X_i)$. Dado que las pérdidas son independientes, la pérdida agregada del sistema es: $X_n=\sum_{i=1}^{n}X_i$. Por la linealidad de la esperanza, $\mathbb E[S_n]=\sum_{i=1}^{n}\mathbb E[X_i]=npL$. Por tanto, la pérdida esperada por entidad sigue siendo: $\mathbb E!\left[\frac{S_n}{n}\right]=\frac{1}{n}\mathbb E[S_n]=pL$, idéntica a la pérdida esperada individual. Sin embargo, el comportamiento de la varianza es diferente. Bajo independencia, $\mathrm{Var}(S_n)=\sum_{i=1}^{n}\mathrm{Var}(X_i)=n,\mathrm{Var}(X_i)$. Por consiguiente, la varianza de la pérdida media por entidad es $\mathrm{Var}!\left(\frac{S_n}{n}\right)=\frac{1}{n^2}\mathrm{Var}(S_n)=\frac{\mathrm{Var}(X_i)}{n}$.

Así, aunque la pérdida esperada permanece inalterada, la dispersión de la pérdida media disminuye proporcionalmente a ($1/n$). Equivalentemente, su desviación típica disminuye a razón de ($1/\sqrt{n}$).
} Este es el resultado fundamental de la diversificación: \textbf{la agregación de riesgos independientes no altera el primer momento de la distribución (la esperanza), pero reduce su segundo momento (la varianza), haciendo que las pérdidas efectivas se concentren cada vez más en torno a su valor esperado conforme aumenta el número de entidades participantes}.

Ahora bien, este enfoque presenta serias lmiitaciones. Nos centraremos en dos problemas particulares. 

En primer lugar, ¿qué efecto podría tener el establecimiento de estos fondos sobre el grado de riesgo en las estrategias de las entidades? Esta es una cuestión fundamental: cómo afecta el FGD a los incentivos de las entidades financieras. Aunque estos mecanismos contribuyen a estabilizar el sistema y a prevenir retiradas masivas de depósitos, también pueden generar problemas de riesgo moral al alterar la distribución de pérdidas entre accionistas, depositantes y fondo de garantía.

Para ilustrarlo, considérese un banco que opera durante dos períodos. En el primero capta depósitos ($D$) y paga una prima de seguro al FGD ($\pi$); en el segundo, sus activos son liquidados ($V$). Si el valor recuperado de los préstamos ($P'$) resulta insuficiente para cubrir los depósitos, el FGD compensa la diferencia ($S$). La indemnización recibida puede expresarse como:

\eqblock{
\begin{aligned}
    S=\max{D-P’,0},
\end{aligned}
}{}

Donde ($D$) representa los depósitos y ($P’$) el valor recuperado de la cartera de préstamos. Supongamos ahora que el resultado de la cartera crediticia es incierto. Con probabilidad ($\theta$) el proyecto tiene éxito y genera un valor ($X$), mientras que con probabilidad ($1-\theta$) fracasa y no genera recuperación alguna. En estas condiciones, el beneficio esperado para los accionistas puede escribirse como:

\eqblock{
\begin{aligned}
    \mathbb E[R]=\underset{\scriptsize{\substack{\text{Rendimiento esperado}\\\text{cartera de préstamos}}}}{(\theta X-P)}+\underset{\scriptsize{\substack{\text{Valor esperado}\\\text{protección del FGD}}}}{\bigl((1-\theta)D-\pi\bigr)}
\end{aligned}
}{}

La expresión pone de manifiesto que el valor esperado de la entidad tiene dos componentes: (i) el primero corresponde al rendimiento esperado de la cartera de préstamos; y, (ii) el segundo refleja el valor esperado de la protección proporcionada por el FGD, neto de la prima pagada. Si la prima fuese actuarialmente justa, este segundo término sería nulo y se obtendría el resultado estándar de MODIGLIANI-MILLER: la estructura financiera no alteraría el valor de la entidad. Sin embargo, en la práctica las primas suelen fijarse de manera relativamente uniforme y no reflejan plenamente el riesgo específico de cada banco.

En ese contexto surge el problema de riesgo moral. Si dos proyectos ofrecen el mismo rendimiento esperado, pero uno de ellos incorpora una menor probabilidad de éxito y, por tanto, una mayor exposición al riesgo, los accionistas tendrán incentivos a preferir este último. La razón es que una parte de las pérdidas asociadas a los escenarios desfavorables es absorbida por el FGD, mientras que los beneficios de los escenarios favorables permanecen en manos de los accionistas. En otras palabras, la existencia de la garantía introduce una asimetría que reduce el coste privado del riesgo asumido por la entidad.

La raíz del problema reside en considerar una prima de seguro exógena e independiente del perfil de riesgo de la cartera crediticia. Cuando la contribución al fondo no aumenta con el riesgo asumido, las entidades no internalizan completamente el coste de sus decisiones. Por ello, una de las principales respuestas regulatorias consiste en vincular las primas del seguro de depósitos al riesgo efectivo de cada entidad, de modo que el coste de la cobertura refleje adecuadamente la probabilidad de generar pérdidas para el sistema.\footnote{No entramos en la obtención de la prima óptima por parte del FGD, pero la fundamentación teórica es la misma que la obtención de la prima que aplican las entidades a los préstamos de su cartera crediticia en el marco de la gestión del riesgo de crédito (MERTON 1974, estudiado en el tema 3B26). No debería sorprender esto al opositor, ya que, en esencia, el problema es el mismo.
} De hecho, ése es en espíritu el sistema seguido por la mayoría de economías avanzadas hoy.\footnote{Cuestiones adicionales que se consideran a la hora de diseñar este tipo de instrumentos en la práctica son, entre otras, las siguientes: ¿Deben de cubrirse el 100\% de los depósitos o sólo parte de ellos (hasta un importe determinado, sólo algunos depositantes, etc.)? ¿Es suficiente el volumen del fondo (en especial, de tener lugar crisis sistémicas)? ¿Debe participar el Estado en la financiación del fondo?
}

El establecimiento de un Fondo de Garantía de Depósitos de esta índole permite proteger a los depositantes y, de hecho, ha proporcionado un grado de estabilidad al sistema: con un correcto pricing de las primas a pagar por los bancos, se alcanza a minimizar los problemas de riesgo moral. 

Sin embargo, para algunas stuaciones este sistema no es suficiente. Y aquí entra el segundo punto, ¿qué pasaría si experimentásemos un fenómeno de clusterización de los eventos de pérdida? Es decir, ¿qué pasaría si la volatilidad no disminuyese, sino que aumentase con el número de entidades pertenecientes al FGD?

Como podemos intuir, este es una de las propiedades emergentes más características de las crisis sistémicas: cuando todo el sistema se derrumba, a la vez y en todas partes. Naturalmente, este instrumento resulta del todo insuficiente para hacer frente a estos sucesos. 

\paragraph{Prestamista de último recurso}
Con este escenario en mente, presentamos otro instrumento, posiblemente el único capaz de hacer frente a estos eventos sistémicos de riesgo de liquidez: la figura del \textbf{prestamista de último recurso}.\footnote{La génesis de esta vertiente fue propuesta por BAGEHOT (1873). El autor sugirió que, para prevenir el fracaso de instituciones financieras solventes pero ilíquidas, el banco central debería prestar: (i) sin límite, (ii) a cambio de activos de garantía de calidad; y, (iii) a un tipo penalizador frente al tipo de interés formado en el mercado monetario. El primero de los elementos significa que la entidad recibiría toda la liquidez (reservas) que solicitase, por lo que ella misma evaluaría su situación financiera, explotando su ventaja informacional, y solicitaría el volumen adecuado. El segundo serviría como mecanismo para filtrar a las entidades solventes de las no-solventes, ya que las segundas no serían capaces de aportar dichos activos como garantía (colateral). El tercero supone la implementación de un desincentivo a la toma excesiva de riesgos, y por tanto al posible riesgo moral. Esta medida supone varias ventajas y, de hecho, es una herramienta que existe en la mayoría de sistemas bancarios. 

Sin embargo, los mecanismos propuestos inicialmente por BAGEHOT no parecen suficientes para superar los problemas de riesgo moral y de señalización. Además, cuenta entre otras con los siguientes defectos adicionales: puede afectar a la autonomía y a los objetivos de inflación de la autoridad monetaria, de no ser suficiente la inyección puede favorecer la aparición de contagio y de sendas explosivas dentro del sistema financiero y monetario, implica la socialización de pérdidas de las entidades. Ante las limitaciones de estas (y otras) alternativas, y ante la percepción de que la gestión de riesgos y el establecimiento de requerimientos microprudenciales pueden no resultar suficientes para evitar la caída de las entidades, el sistema más ampliamente empleado para proteger los ahorros de los depositantes son los fondos de garantía de depósitos. 
} 

La idea central de este instrumento es establecer una entidad que actúe como prestamista de último recurso ante situaciones extremas. Lo más natural, por su condición de monopolista de la emisión de dinero, es que este rol sea asumido por la Autoridad Monetaria. Veamos cómo podría funcionar.

Supongamos que un Banco Central anuncia que está dispuesto a prestar al banco a un tipo de interés penalizador, pero suficientemente moderado como para que las entidades financieras pudiesen repagar los préstamos sin mermar la confianza de sus depositantes más pacientes. En esta situación, las entidades podrían liquidar y reembolsar los depositos de aquellos agentes que acudiesen en primer lugar, pero manteniendo los suficientes márgenes como para poder reembolsar más adelante los costes del préstamo, el capital y los rendimientos prometidos a sus depositantes más pacientes. De esta forma, incluso con heterogeneidad de agentes privados, la desconfianza en el Sistema y el rápido deterioro de los balances de estas entidades podrían, a priori, ser sorteados. (Eliminádose, como luego se verá, el equilibrio de NASH que genera el pánico bancario).

Sin embargo, la presencia de un prestamista de último recurso no está exenta de limitaciones. En particular, encontramos los siguientes problemas que resulta díficil descartar:
\begin{la}
    \item \textbf{Dinero en circulación: Inflación}. En primer lugar, aunque su presencia pudiese estabilizar la confianza del sistema en términos de solencia y liquidez, no podría lidiar con la inestabilidad de precios que pudiese generar. Es decir, como su presencia no eleva per se la cantidad de bienes en la economía, si se introdujera dinero sería muy posible acabar en un escenario inflacionista. Esto afectaría jugaría un papel doble. En primer lugar, el aumento de la oferta monetaria podría elevar el nivel de precios reduciría los incentivos para restirarse anticipadamente por las mermas que supondría sobre el consumo y disminuiría el valor real de las exposiciones de las entidades. Sin embargo, la estabilidad de precios a largo plazo de la economía quedaría fuertemente dañada, puesto que su figura, per se, reduciría la confianza en sus compromisos inflacionarios;\footnote{Algunos autores, resuelven esta limitación señalando que, al estar respaldado por una autoridad fiscal con la capacidad de recaudar impuestos, permitiría contar con un mecanismo de contrabalance que actuase como freno en los precios y, además, permitiría contar con los recursos necesarios en caso de una eventual retirada anticipada de depósitos que pudiera conducir a un pánico bancario.
    } y,
    \item \textbf{Deaslineamiento de incentivos: Excesiva asunción de riesgos}. Por otro lado, la existencia del prestamista de último recurso, podría seguir generando problemas de desalieneamiento de incentivos por parte de las entidades. Si el Banco Central prestase a un tipo verdaderamente penalizador, las entidades no podrían devolver el capital e interés habiéndose materializado un pánico bancario, aún a escala reducida. De esta forma, el Banco Central se ve obligado a prestar a un tipo solo lo \textit{suficientemente penalizador}, por debajo de la rentabilidad de mercado, por lo que las entidades tendrían incentivos a pedir prestado para atender las demandas de todos sus solicitantes de forma que no liquidarían ningún proyecto en sus manos y acabarían obteniendo un beneficio a su costa.\footnote{Piénsese en el tipo de la facilidad marginal de crédito del Eurosistema, siempre por encima de los tipos overnight del mercado interbancario.
    } Para sortear este problema, una posible vía sería que el Banco Central obligase a las entidades a liquidar activos para atender al menos una proporción de las solicitudes, y prestarle solo cuando superen esa proporción. 
\end{la}

\subsubsection{Crisis Financiera 2008: Limitaciones}
\textcolor{red}{\textbf{OJO -- (ChatGPT)} Fuentes: Diamond y Dybvig (1983); Brunnermeier (2009); Gorton (2010); Bernanke (2015); Basel Committee on Banking Supervision (2013).

\textbf{Riesgo de liquidez:} cuando aparecieron las pérdidas y la incertidumbre, desapareció la financiación y colapsó la confianza.}

Una vez más, la crisis financiera internacional puso de manifiesto que la gestión de la liquidez no puede limitarse a garantizar la disponibilidad de efectivo en circunstancias normales. Incluso entidades aparentemente solventes pueden enfrentarse a graves dificultades cuando desaparece la confianza de los participantes del mercado y las fuentes habituales de financiación dejan de estar disponibles.

Durante las décadas previas a la crisis, muchas instituciones financieras habían modificado sustancialmente sus modelos de negocio. Junto a los depósitos tradicionales, una parte creciente de su financiación procedía de mercados mayoristas a corto plazo, especialmente mediante repos, papel comercial y préstamos interbancarios. Esta estructura permitía reducir costes de financiación, pero incrementaba la dependencia de la renovación continua de pasivos de muy corto vencimiento.

Mientras las condiciones de mercado fueron favorables, este modelo funcionó sin dificultades aparentes. Sin embargo, el deterioro de los activos hipotecarios y la creciente incertidumbre sobre la exposición real de las entidades provocaron un rápido aumento de la desconfianza entre intermediarios financieros. Como consecuencia, muchas instituciones comenzaron a restringir sus préstamos a otras entidades, acumulando liquidez por motivos precautorios.

El resultado fue una \textbf{paralización progresiva de los mercados de financiación mayorista}. Las entidades dejaron de enfrentarse únicamente a pérdidas derivadas de sus activos, para enfrentarse también a la imposibilidad de refinanciar sus pasivos a corto plazo. La crisis pasó así de ser un problema de calidad de activos a convertirse en un problema de liquidez sistémica.

La quiebra de Lehman Brothers en septiembre de 2008 constituyó el punto de inflexión más significativo. La decisión de permitir su insolvencia generó incertidumbre sobre el compromiso de las autoridades con el rescate de instituciones sistémicas y provocó un colapso generalizado de la confianza. A partir de ese momento, numerosas entidades financieras encontraron crecientes dificultades para acceder a financiación, independientemente de su situación patrimonial concreta.

La crisis también reveló las limitaciones de los mecanismos tradicionales de protección frente a problemas de liquidez. Los fondos de garantía de depósitos contribuyeron a evitar retiradas masivas por parte de los depositantes minoristas, pero resultaron insuficientes para contener la fuga de financiación institucional. Del mismo modo, los prestamistas de última instancia pudieron proporcionar liquidez de emergencia, pero solo mediante intervenciones extraordinarias de una magnitud sin precedentes.\footnote{La principal lección fue que la liquidez constituye una característica sistémica y no exclusivamente individual. Una entidad puede parecer suficientemente líquida mientras los mercados funcionan con normalidad, pero perder rápidamente esa condición cuando desaparece la confianza colectiva. Por esta razón, la regulación posterior a la crisis incorporó requisitos específicos de liquidez, como el Liquidity Coverage Ratio (LCR) y el Net Stable Funding Ratio (NSFR), destinados a limitar la dependencia de financiación inestable y reforzar la capacidad de resistencia de las entidades ante episodios prolongados de tensión financiera.
}

\begin{specialblock}
    Ahora bien, si estos mecanismos fueron tan satisfactorios, ¿dónde están los rescates? ¿por qué fueron, entonces, eso de los rescates que tanto permanece en la memoria colectiva? ¿Qué es y por qué surjen los rescates bancarios que experimentamos en la crisis financiera de 2008? ¿No se nos argumentó que era la única vía para evitar el colapso del sistema financiero? 

    En efecto, los rescates bancarios emergen como antagónicos a lo que se conoce como mecanismos de resolución. Estos suceden, y sucedieron, cuando en las entidades financieras (en particular, las instituciones de los mercados más opácos) y el conjunto del sistema se experimentan dos fénomenos paralelos: (i) por un lado, y más a nivel sistémico, las entidades financieras sufren deterioros de balance tales que no pueden afrontar nuevas cargas de pasivo (incluso con el prestamista de último recurso); y, (ii) determinadas entidades, a título particular, han adquirido una relevancia y magnitud tal dentro del propio Sistema Financiero que se han convertido en lo que se conoce como \textbf{entidades sistémicas}. 
    
    Dicho de otra forma, tiene \textbf{potencial de convertirse en una fuente de riesgos sistémicos}.\footnote{Importante: mientras que el aseguramiento de los depósitos se refiere exclusivamente a los bancos (como únicas entidades financieras que pueden captar depósitos), el concepto de sistemicidad aplica potencialmente a cualquier tipo de entidad financiera. En concreto, desde la crisis de 2007- 08, el FSB ha realizado trabajos para el establecimiento de mecanismos de resolución homogéneos para: (i) bancos, (ii) aseguradoras, (iii) cámaras de contrapartida central.
    } Con esto en cuenta, se generan dos problemas desde la óptica del policy-maker:
    \begin{lnum}
        \item \textbf{Expectativas: promesas no creíbles}. En primer lugar, la promesa de no rescate puede no ser creíble. Esto genera una situación de riesgo moral: la entidad se siente libre para asumir riesgos ya que da por descontado que el Estado habrá de acudir a su rescate en caso de situación crítica a través de un bail-out;
        \item \textbf{Coste: prima de riesgo}. El segundo problema consiste en que, de acudir el Estado al rescate, no sólo los contribuyentes pagarían el rescate de la entidad, sino que además el propio rescate podría suponer un impacto fiscal que afectase al activo libre de riesgo: incluso con rescate no se podría evitar el episodio de crisis sistémica. 
    \end{lnum}
    
    Así, las situaciones arriba descritas ponen de manifiesto la necesidad de contar con normas de resolución para que las autoridades puedan gestionar la crisis de una entidad de una manera reglada, rápida y segura, no sólo para garantizar la estabilidad financiera, sino también para que sean sus propietarios y acreedores (y no los ciudadanos con sus impuestos) quienes asuman las pérdidas.\footnote{¿Qué es la resolución de una entidad financiera? Se trata del marco y los procedimientos previstos para llevar a cabo una restructuración ordenada en una entidad en lugar de una liquidación concursal; es decir, en lugar de cerrar el negocio y hacer frente al pasivo con el activo disponible, se trata de salvar la mayor parte posible del negocio: se busca mantener la prestación de las funciones esenciales de la entidad (generalmente, a través de una nueva entidad). Para ello, existe un menú de herramientas de que disponen las autoridades de resolución. 
    } Es decir, se busca un bail-in que permita: (i) reducir el riesgo sistémico, (ii) repartir el coste del riesgo entre los agentes privados. 
    
    Lógicamente, éstas dependen del sector concreto. En el caso de los bancos ([\authorfont{Ver Tema 3.B.44}]), la UE ha ido más allá de la coordinación internacional de política financiera en materia de resolución y aplica este marco no sólo a las entidades sistémicas, sino a todos los bancos con licencia (si bien con diferentes grados de exigencia).
    
    Si el gobierno es capaz de garantizar que todo aquel que espere a retirar sus depósitos hasta $T=2$ obtendrá el rendimiento prometido, entonces se consigue evitar el equilibrio de NASH no deseado. Sin embargo, para que esto ocurra, el seguro de depósito gubernamental debe ser creíble, es decir el gobierno debe tener la forma de obtener los recursos para hacer frente a los posibles pagos. Los autores argumentan que el gobierno puede conseguir estos ingresos a través de impuestos diseñados para hacer frente a eventuales necesidades de ejercitar el seguro de depósitos.  Sin embargo, la existencia de correlaciones entre las posiciones (en este caso, entre los sucesos de quiebra) altera sustancialmente este resultado. En la práctica, los riesgos no son independientes dada la existencia de interrelaciones e interconexiones entre entidades. Por ese motivo, han de considerarse las influencias recíprocas en el cálculo de la varianza.
\end{specialblock}


\clearpage
\section{When everything goes wrong: Momentos de crisis}
\puntosclave{
    Existen varias aproximaciones teóricas a las crisis financieras, como el mecanismo circular deuda-deflación o las burbujas de activos. 
    
    La evidencia empírica sobre las crisis parece alcanzar el resultado de que la mayoría de estos eventos tienen varios patrones comunes. 
    
    El modelo de Diamond-Dybyig consiste en un juego con dos equilibrios de NASH. En el primer equilibrio, un depositante puede creer que se avecina una crisis bancaria y que todos los demás depositantes intentarán obtener fondos líquidos. Como resultado, su estrategia óptima es retirar sus propios activos líquidos de inmediato. Sigue un ataque especulativo y los bancos se quedan sin fondos líquidos. Un equilibrio alternativo es aquel en el que nadie cree que se producirá una corrida bancaria y los bancos tienen suficientes fondos para satisfacer las verdaderas demandas de liquidez, de modo que no se desarrolla ninguna crisis.
}

A lo largo del apartado anterior hemos presentado los principales riesgos a los que se enfrenta el Sistema Financiero y la paleta de instrumentos de gestión de los que dispone para mitigarlos. Sin embargo, como la memoria colectiva no falla, surgen aún dudas que quedan sin resolver. Como sabemos, a partir de septiembre de 2008, las autoridades monetarias y fiscales desplegaron medidas extraordinarias de apoyo al sistema financiero: inyecciones masivas de liquidez, garantías públicas, recapitalizaciones bancarias y programas de compra de activos. La mera magnitud de estas intervenciones pone de manifiesto una cuestión fundamental: si los instrumentos de gestión del riesgo estaban diseñados precisamente para evitar este tipo de situaciones, ¿por qué fue necesario recurrir a rescates de semejante escala? 

La respuesta exige abandonar el análisis aislado de cada riesgo e introducir una perspectiva sistémica sobre las crisis financieras.

\textbf{¿Qué queremos decir?} El término crisis financiera se utiliza para referirse a aquella situación en la que se produce una crisis económica que no tiene origen en la economía real, sino que está asociada a problemas originados en el sistema financiero o monetario. Por lo tanto, en primer lugar presentaremos dos enfoques teóricos que pretenden explicar los comportamientos y características que pueden originar una crisis financiera; centrándonodes en las teorías explicativas de las crisis en los mercados financieros (renta variable, renta fija y derivados: estallido de burbujas) y las crisis en el sistema bancario (el conocido \textit{pánico bancario}). A continuación, y en segundo lugar, nos alejaremos del enfoque circunscrito al Sistema Financiero para adoptar una visión más general que nos permita ver los canales y mecanismos a través de los cuales una situación de inestabilidad originada en algún segmento concreto del Sector Financiero puede acabar provocando una auténtica crisis financiera sistémica.

\subsection{Aproximación parcial: ¿cómo y por qué emergen situaciones de inestabilidad?}
Como podemos intuir, las crisis financieras no son un grupo homogeneo. La literatura identifica, entre otras: (i) las crisis bancarias; (ii) las crisis de deuda (interna/externa); (iii) las crisis cambiarias; y/o, (iv) el estallido de burbujas; claificación que responde, esencialmente, a la  segmento del Sistema Financiero donde la inestabilidad emerge en primer lugar. Por tanto, no debemos etenderla como categorías alternativas de crisis sino que, más bien, en muchas ocasiones lo que observamos es que un tipo de crisis puede verse precedida o exacerbada por otra.

Por razones de tiempo, vamos a centrarnos en dos enfoques que no son tratados en ninguna otra parte del temario: las crisis bancarias y el estallido de burbujas. Por este motivo, conviene aclarar también los riesgos que entran en juego en cada una de estas categorias:
\begin{la}
    \item \textbf{Riesgo de mercado y crédito}. Cuando hablamos de crisis generadas por el \textbf{estallido de burbujas}, nos referimos al súbito deterioro del valor de los activos en balance. De esta forma, estas crisis se originan cuando los instrumentos de gestión de \textbf{riesgos de mercado y riesgos de crédito} han resultado insuficientes para mitigar sus efectos adversos, conduciendo a las entidades a una situación de grave desequilibrio y peligro de quiebra; y,
    \item \textbf{Riesgo de liquidez}. Por otro lado, cuando hablamos de \textbf{crisis bancarias}, nos referimos a situaciones en las que, por motivos de diversa índole, existe un desmedido interes en liquidar y retirar los depósitos bancarios; situación que, debido a la estructura temporal del balance, estas entidades no son capaces de afrontar. Por tanto, se trata de una situación generada en el momento en el que, a priori, los instrumentos de gestión del \textbf{riesgo de liquidez} han resultado insuficientes, conduciendo a las entidades bancarias a una situación de grave insolvencia.
\end{la}

Empecemos con las crisis generadas por el estallido de burbujas. 

\subsubsection{Mercados transparentes: ¿Por qué no funcionan como nosotros decimos que funcionen?}
En primer lugar, ¿cómo podemos identificar una burbuja? ¿qué es una burbuja? Esta respuesta no está para nada clara y la razón es extremadamente sencilla: una burbuja se definen como una situación en las que el precio de cotización del activo supera su valor fundamental del activo. Así que, necesariamente debemos poder descomponer el precio del activo en:

\eqblock{
\begin{aligned}
    \boxed{p_t=v_t+b_t}
\end{aligned}
}{}

Como hacerlo es, naturalmente, conflictivo. En primer lugar, y quizás más importante, porque supone descartar, por definición, la Hipótesis de los Mercados Financieros, que le granjeó (lamentablemente) el Nobel de Economía a FAMA en 2013, y, brevemente, exige que el \textbf{precio sea el correcto} y \textbf{no existan comidas gratis} (inexistencia de arbitraje). 

A pesar de esta aparente dificultad, y dado que los precios de los activos afectan la asignación real de una economía, los economistas han dedicado tiempo a la observación y análisis de estos comportamientos excepcionales y aparentemente irracionales ya que es importante entender las circunstancias en las que estos precios podrían desviarse de su valor fundamental.\footnote{Esto puede ocurrir por distintos motivos, todos ellos relacionados con las expectativas de los inversores, como se verá a continuación. Ejemplos históricos famosos son la manía de los tulipanes en los Países Bajos (1634-37), la Burbuja del Misisipi (1719-20), la Burbuja de la Compañía del Mar del Sur (1720) y los "Felices años 20"que precedieron al colapso de 1929. Más recientemente, hasta marzo de 2000, los precios de las acciones de empresas relacionadas con internet (dot com) se dispararon antes de caer más del 75\% a finales de aquel año.
}

La literatura de burbujas financieras suele dividirse en cuatro grupos de modelos: los dos primeros analizan las burbujas dentro del paradigma de expectativas racionales, pero difieren en el supuesto acerca de si todos los inversores tienen la misma información o están asimétricamente informados; un tercer grupo de modelos se centra en la interacción entre inversores racionales y no racionales (conductuales); y, en el último, son el conjuto de creencias y expectativas de los operadores las que son heterogéneas, pero es un hecho que todos conocen y están de acuerdo en que sea así.  

Presetaremos brevemente dos modelos, uno perteneciente al tercer grupo y un segundo perteneciente al cuarto. En todo caso, tengamos presente que, como las burbujas suelen asociarse con dramáticos incrementos en el precio de los activos seguidos por un colapso, los modelos se centrarán únicamente en explicar cómo esto sucede y las razones por las que sucede, pero poco sabremos sobre las causas que inician y/o terminan la burbuja, así como de sus consecuencias.  

\paragraph{Burbujas: Límites del arbitraje}
Imaginemos un mercado formado por dos tipo de inversores. Un primer tipo, que llamaremos inversor racional, actúa de manera diligente: recabando información, proyectando y formando sus expectativas futuras racionalmente y valorando adecuadamente los activos financieros. El seguno tipo, al que llamaremos inversor conductual, actúa de manera semi-diligente, pues revisa con frecuencia la información del mercado y la conducta de los agentes, pero la mayoría de sus decisiones vienen dadas por un conjunto de reglas heurísticas que marcan su pauta de conducta. (Por ejemplo, si el inversor conductual observa un incremento pronunciado del precio de sus activos, es probable que se deje llevar por el efecto rebaño y acaba comprando más, aunque sepa que está ligeramente sobrevalorado).\footnote{Bajo el paraguas de unos mercados perfectamente eficientes (HME), este tipo de conductas no tendrían posibilidad de mantenerse en el tiempo, ya que los inversores racionales, que son perfectamente conocedores de que ese activo esta sobrevalorado, podrían vender en corto y acabarían devolviendo su precio al valor fundamental. (Los inversores racionales deberían ir en contra de la burbuja incluso antes de que surja). Ahora bien, ¿qué pasaría si esto no pudiese hacerse o si no tuviere el impacto previsto? Esto es precisamente lo que sostiene la literatura sobre los límites de arbitraje. 
}

Lo natural sería que los inversores racionales puediesen arbitrar en corto (o largo) y desinflar cualquier desviación de precios que pudiese observarse, incluso antes de que sucediera. Ahora bien, ¿qué pasaría si no pudiesen? Si los inversores racionales no pudieran llevar a cabo dichas operaciones, o no pudieran hacerlo con la intensidad tal que devolviésen siempre los precios a su valor fundamental, las desviaciones pronunciadas (¡tanto positivas como negativas!) sobre su valor fundamental pueden perfectamente  persistir e impidir que los arbitrajistas racionales corrijan la mala fijación de precios. 

Las razones por las que esta situación podría experimentarse son varias. Algunas de las más destacadas por la literatura son:
\begin{lnum}
    \item \textbf{Aversión al riesgo}. En primer lugar, el riesgo fundamental hace arriesgado vender en corto un activo de burbuja, ya que potenciales alteraciones posteriores en los fundamentos podrían deshacer la sobrevaloración inicial. La aversión al riesgo limita la agresividad de los operadores racionales si no hay sustitutos cercanos ni coberturas cercanas disponibles;
    \item \textbf{Potencia del ruido}. En segundo lugar, los operadores racionales también enfrentan el riesgo de operadores con ruido (DELONG et al., 1990). Oponerse a la burbuja es arriesgado incluso sin riesgo fundamental, ya que los operadores irracionales que generan ruido podrían aumentar el precio aún más en el futuro y ampliar temporalmente la determinación incorrecta de precios. Así, los operadores racionales incorporan en su toma de decisiones el valor fundamental a largo plazo, pero priorizan los horizontes cortos: como resultado, sólo corrigen parcialmente la mala fijación de precios;\footnote{Por ejemplo, en un mundo con gestión de carteras delegada, los gestores de fondos de inversión a menudo se preocupan por los movimientos de precios a corto plazo, porque las pérdidas temporales provocan salidas de fondos (SHLEIFER y VISHNY, 1997). Una sobrevaloración continuada en el tiempo y la posterior salida de fondos obligarían a los gestores de fondos a deshacer sus posiciones, paradójicamente, exactamente en el momento en que el mispricing es mayor. Anticipándose a este escenario posible, los gestores no se enfrentan agresivamente a la fijación incorrecta de precios.
    } y,
    \item \textbf{Sincronización}. En tercer lugar, los operadores racionales enfrentan el riesgo de sincronización (ABREU y BRUNNERMEIER, 2002, 2003). Dado que un solo operador no puede típicamente llevar el mercado a la baja por sí mismo, se requiere coordinación entre operadores racionales y surge un problema de sincronización. Cada operador racional enfrenta el siguiente dilema: si ataca la burbuja demasiado pronto, renuncia a las ganancias de la posterior subida causada por los operadores de impulso conductual; si ataca demasiado tarde y permanece invertido en el activo de burbuja, sufrirá el posterior colapso. Cada operador intenta prever cuándo actuarán los otros operadores racionales.\footnote{Este ejercicio es complejo, ya que los operadores se dan cuenta secuencialmente de la burbuja y no saben en qué orden en la cola están. Es precisamente esta falta de conocimiento común la que eliminaría el impacto del argumento estándar de inducción hacia atrás: dado que no hay un punto en el tiempo comúnmente conocido desde el cual se podría comenzar la inducción hacia atrás, incluso las burbujas con horizontes finitos pueden persistir.
    }
\end{lnum}

En sintésis, la razón principal por la que estos inversores racionales puedan enfrentar restricciones en términos operativos puede expresarse con el sonado aforismo de KEYNES: \textit{Los mercados pueden permanecer irracionales mucho más tiempo del que tu te puedes mantener solvente}.\footnote{Otra línea complementaria e interesante de la literatura apunta que son precisamente los agentes racionales quien construyen la burbuja. Ante noticias relativamente insignificantes se desencadenarían grandes movimientos de precios, porque dichas noticias permitirían solventar el problema de sincronización arriba explicado (ABREU y BRUNNERMEIER, 2003): los operadores racionales prefieren montar la burbuja en lugar de atacarla. Empíricamente, hay evidencia de apoyo a la \textbf{hipótesis de montar la burbuja}. Por ejemplo, entre 1998 y 2000, los fondos de cobertura estaban fuertemente inclinados hacia acciones tecnológicas altamente valoradas (BRUNNERMEIER y NAGEL, 2004). Contrariamente a la HME, los fondos de cobertura no fueron una fuerza correctiva de precios, a pesar de ser algunos de los inversores más sofisticados y tener un perfil cercano al de los arbitrajistas racionales de la teoría económica. De manera similar, TEMIN y VOTH (2004) documentan que Hoares Bank contribuyó a la burbuja de la South Sea en 1719-20, a pesar de dar numerosas indicaciones de que creía que la acción estaba sobrevalorada. Muchos otros inversores, incluido NEWTON, también intentaron montar la burbuja de la South Sea pero con menos éxito. Frustrado con su experiencia comercial, NEWTON concluyó: \textit{Puedo calcular los movimientos de los cuerpos celestiales, pero no la locura de la gente} (KINDLEBERGER 2005, p. 41).
}

\paragraph{Burbujas: Creencias heterogeneas}
Otra linea de investigación paralela razona que las desviaciones sostenidas pueden surgir cuando los inversores tienen creencias heterogéneas y enfrentan restricciones en su capacidad de realizar ventas en corto. 

Imaginemos un mercado atomizado en el que cada operador tiene sus propias creencias sobre el rendimiento futuro de un activo financiero, que puede verse motivado por infinitud de fenómenos, típicamente sesgos conductuales. En este escenario, las creencias de los operadores son heterogéneas. Por ejemplo, si un inversor tiene demasiada confianza en sus propias señales, contará con una distribución previa diferente (con menor varianza) sobre el término de ruido de las señales. De hecho, incluso con información completa y simétrica, los operadores podrían seguir estando en desacuerdo. La combinación de creencias heterogéneas con restricciones en las ventas en corto puede resultar en sobrevaloración, ya que los optimistas elevan el precio del activo, mientras que los pesimistas no pueden contrarrestarlo, dadas las mencionadas restricciones (MILLER, 1977). 

OFEK y RICHARDSON (2003) modelizan la burbuja \textit{punto com} empleando estos supuestos. En un modelo dinámico, el precio del activo puede incluso superar la valoración del inversor más optimista en la economía. Esto es posible, ya que los inversores actualmente optimistas, los propietarios actuales del activo, tienen la opción de revender el activo en el futuro a un precio alto siempre que se vuelvan menos optimistas. Se asume que el grado de optimismo oscila entre diferentes grupos de inversores, por lo que, en ese momento, existirán otros operadores optimistas que estarán dispuestos a comprar el activo (HARRISON y KREPS 1978).

En resumen, la literatura sobre burbujas ha avanzado enormemente desde la década de 1970 y ha llevado a varias clases de modelos con pruebas empíricas distintas. Sin embargo, \textbf{muchas preguntas siguen aún sin resolverse}. En especial, la literatura no ha logrado construir modelos convincentes que expliquen cuándo y por qué comienzan las burbujas. Además, mientras que en la mayoría de los modelos las burbujas estallan, \textbf{la regularidad empírica es que las burbujas parecen desinflarse al cabo de unas semanas o meses}. Por último, pese a que la introducción de elementos propios de las finanzas conductuales supone una mejora sustancial en el entendimiento de este fenómeno económico, los problemas de generabilidad limitan el alcance y la aplicación de los modelos basados en sesgos.

Ahora que conocemos algunos de los determinantes que explican la emergencia de riesgos, y crisis, en los mercados más transparentes del Sistema Financiero, veamos una linea de investigación paralela que pretende explicar la aparición de crisis en los mercados más ópacos. 

\subsubsection{Mercados opácos: Crisis bancarias}
Como hemos visto anteriormente, el principal riesgo que enfrenta el Sistema Bancario es el \textbf{riesgo de liquidez}, ya que el sistema de reserva fraccional es inherentemente frágil: mientras que los depósitos se pueden recuperar en plazos temporales cortos (depósitos a plazo) o incluso en cualquier momento (depósitos a la vista), los plazos de recuperación del dinero prestado por parte de los bancos suelen ser más largos; 

Esta labor de transformación de plazos que lo expone al riesgo depende crucialmente de la confianza para su estabilidad: cuando la confianza entre las partes se quiebra, y un porcentaje muy elevado de depositantes quiere recuperar su dinero, el banco puede quedarse sin liquidez, fenómeno que se conoce como pánico bancario (o corrida bancaria) y puede provocar que el sistema bancario en su conjunto colapse.\footnote{ La literatura sobre crisis bancarias se puede organizar en dos grandes grupos:
\begin{la}
    \item \textbf{Pánicos bancarios}. La primera visión se refiere a los mencionados pánicos bancarios, y sostiene que las crisis bancarias son eventos aleatorios, no relacionados con cambios en la economía real. Las crisis bancarias pueden surgir de expectativas autocumplidas, como lo modelaron entre otros Diamond y Dybvig (1983), como veremos a continuación; y,
    \item \textbf{Ciclo económico real}. Frente a los pánicos bancarios, la segunda visión sostiene que las crisis bancarias están relacionadas con el ciclo económico real y son desencadenadas por cambios bruscos en el riesgo agregado. Dado que es probable que una recesión económica reduzca el valor de los activos de un banco, las señales sobre una recesión inminente pueden inducir a los depositantes a considerar la retirada de sus fondos. Por ejemplo, KAUFMAN (1998) argumenta que los bancos fracasan debido a la exposición a un mismo shock común (por ejemplo, contracciones en la economía o en los mercados de valores e inmobiliario), en lugar de estar expuestos a otros fallos bancarios que fueron el resultado de factores idiosincráticos.
\end{la}

No obstante, y atendiendo a la categorización que ROSS propone de los segmentos del Sector Financiero, en ambas lineas juega un papel esencial la idea de que la información asimétrica es clave en el desencadenamiento de una crisis bancaria. Además, los problemas de información asimétrica generalmente se intensifican durante períodos de fuertes caídas en los precios de los activos, a medida que los valores colaterales de los prestamistas y el valor neto de las empresas prestatarias disminuyen. Esto, a su vez, aumenta la posibilidad de una crisis bancaria. MISHKIN (1994) sostiene que se pueden explicar mediante este marco de información asimétrica un gran número de crisis financieras en Estados Unidos que ocurrieron en los siglos XIX y XX, y que típicamente comenzaron con caídas en el mercado de valores.

Trabajos recientes combinan la visión del ataque especulativo y la visión del ciclo económico. CHARI y JAGANNATHAN (1988) consideran un modelo donde agentes informados observan una señal negativa sobre el rendimiento de una inversión arriesgada. Sin embargo, los agentes no informados no tienen claridad sobre la motivación de estas retiradas (es decir, si están basadas en información o no) y pueden tomar decisiones similares a las descritas por DIAMOND y DYBVIG (1983). Por lo tanto, las corridas bancarias pueden ocurrir porque los fundamentos parecen malos o porque los inversores creen que las demandas de liquidez son excesivas.

Para un trabajo actualizado sobre la cuestión, véase Alonso, Anguren, Manzano, Mochón (2023).
} 

Por razones de tiempo, presentaremos solamente el modeloque originó esta rama de la literatura.

\paragraph{Modelo: DIAMOND y DYBVIG (1983)}
Se trata del modelo de DIAMOND y DYBVIG, que pretende responder simultáneamente a dos preguntas:
\begin{lnum}
    \item ¿Por qué existen los bancos?
    \item ¿Por qué los bancos son vulnerables a las corridas bancarias incluso cuando son solventes?
\end{lnum}

Para hacerlo, presentaremos una versión simplificada del modelo. En una economía extremadamente sencilla, asumiremos:
\begin{lnum}
    \item Una ventana temporal de tres periodos: $T=\{0,1,2\}$;
    \item Todos los individuos poseen inicialmente una unidad de un bien homogéneo en $T=0$.
    \item Existe una tecnología de inversión con las siguientes característica: si la producción se interrumpe en $T=1$ se obtiene únicamente una unidad de consumo, si se mantiene, la inversión genera un rendimiento $R$. La secuencia tecnológica puede representarse de la siguiente manera: 
    $$T=0\Rightarrow -1;\quad T=1\Rightarrow\left\{\begin{aligned}0\\1\end{aligned}\right\};\quad T=2\Rightarrow\left\{\begin{aligned}R\\0\end{aligned}\right\}$$
    Por tanto, la inversión es ilíquida: genera rentabilidad únicamente si se mantiene hasta vencimiento.\footnote{Se considera un activo ilíquido si los ingresos que se generan por su venta en t (rt) son menores que el valor presente de su liquidación en t+n (rt+n). En el caso extremo, sería totalmente ilíquido si en t no puede venderse o liquidarse por una cantidad positiva (rt =0), pero tiene un valor positivo en t+n. Cuanto menor sea la fracción del valor presente del flujo de caja futuro que se puede obtener hoy por la venta de un activo, menos líquido es.
    }
    \item La incertidumbre relevante no afecta a la rentabilidad de la inversión sino a las necesidades de liquidez de los individuos: existen dos tipos de agentes con información privada acerca de sus preferencias temporales de consumo. Para simplificar, suponemos que existe: (i) una fracción $\theta,\qquad 0<\theta<1$, que descubre en $T=1$ que necesita consumir inmediatamente (\textit{early consumers}); y, (ii) una fracción restante $1-\theta$ que puede esperar hasta T=2 (late consumers).
\end{lnum}

Empecemos por plantear una situación en la que no existen Bancos y la asignación viene dada por un planificador benevolente. Las utilidades de los dos tipos vienen dadas por:\footnote{Esta forma funcional no es general, sino que solo sirve para ilustrar de forma sencilla y concreta el modelo dando lugar resultados de interés económico. Por ejemplo, el supuesto de que $\rho < 1$ tiene el efecto de que el individuo tipo early valore especialmente más que el tipo late el consumo en $T=1$.

El modelo formulado con estas funciones de utilidad implica que si se produce un pánico bancario y el individuo tipo \textit{late} recibe $0$, dada la forma logarítmica de la función de utilidad, obtendrá un bienestar de $-\infty$, lo que haría que los agentes no tuvieran incentivos a poner su dinero en el banco ante la posibilidad de un pánico bancario y privaría a la economía de los beneficios de pool que proporciona la existencia de bancos ante el riesgo de liquidez. D\&D proporcionan dos respuestas alternativas que generalizan el modelo:
\begin{la}
    \item La primera seria que, ante la posibilidad de pánico bancario, los individuos tienen la posibilidad de mantener una determinada cantidad del bien no depositado en el banco y de esta forma se garantizan siempre un consumo positivo en caso de quiebra. De forma que los bancos pueden atraer depositantes incluso existiendo la posibilidad de pánicos bancarios; y,
    \item Alternativamente una modelización diferente de las funciones de utilidad, logarítmicas en el rango $1$, $R$ y no infinitamente negativas cuando el consumo es cero, mantendría la validez del modelo de D\&D para explicar el papel que desempeñan los bancos y porqué son capaces de atraer depositantes incluso con probabilidad positiva de que se produzca una quiebra.
\end{la}
}

\eqblock{
\begin{aligned}
    &t=0\Rightarrow -1;\quad t=1\Rightarrow\left\{\begin{aligned}0\\1\end{aligned}\right\};\quad t=2\Rightarrow\left\{\begin{aligned}R\\0\end{aligned}\right\}\\[2em]
    &\theta :U_a(c)=\ln{c_1^a}\qquad\text{y,}\qquad (1-\theta):U_b(c)=\rho\ln{(c_1^b+c_2^b)}\\[1em]
    &\text{donde,}\qquad 0<\rho<1\qquad \text{y,}\qquad \rho R>0
\end{aligned}
}{}

Donde $c_t^i$ representa el consumo del individuo $i$ en el período $t$. Imaginemos que existe un planificador social que asigna recursos de forma eficiente. Dado que los individuos early únicamente valoran el consumo en $T=1$, nunca tiene sentido asignarles consumo en $T=2$: $c_2^a=0$. Asimismo, dado que los individuos late pueden esperar y la inversión es más productiva cuando madura, tampoco tiene sentido asignarles consumo en $T=1$: $c_1^b=0$. Por tanto, las únicas variables relevantes son: $c_1^a$ y $c_2^b$.

Supongamos que cada individuo \textit{early} consume $c_1^a$. Como la fracción de individuos early es $\theta$, la cantidad total de proyectos que debe liquidarse anticipadamente es: $\theta c_1^a$. Dado que cada liquidación temprana genera exactamente una unidad, esa cantidad permite financiar el consumo temprano. La fracción restante de proyectos ($1-\theta c_1^a$) permanece invertida hasta $T=2$ y cada uno de esos proyectos produce $R$.

Por tanto, la producción total disponible en $T=2$ es: $(1-\theta c_1^a)R$, producción que se distribuye entre los individuos \textit{late}, que constituyen una proporción $1-\theta$ de la población. De aquí obtenemos la restricción de recursos fundamental:

\eqblock{
\begin{aligned}
    &c_2^b=\frac{(1-\theta c_1^a)R}{1-\theta}\qquad\overset{\scriptsize{\text{Sustituyendo en}}}{\Longrightarrow}\qquad \mathbb E[U]=\theta \overset{\scriptsize\text{utilidad grupo A}}{\ln(c_1^a)}+(1-\theta)\underset{\scriptsize\text{Utilidad grupo B}}{\rho\ln(c_2^b)}\\[2em]
    &\boxed{\begin{aligned}
        \max_{c_1^a}\mathbb E[U]=\theta \ln(c_1^a)+(1-\theta)\rho\ln\left(\frac{(1-\theta c_1^a)R}{1-\theta}\right).
    \end{aligned}}\\[1em]
    &\underset{\scriptsize{\substack{\text{Derivando}\\\text{respecto $c_1^a$}}}}{\Longrightarrow}\qquad \text{CPO}:\frac{\theta}{c_1^a}=\frac{\rho\theta}{1-\theta c_1^a}\quad \Longrightarrow\quad 1-\theta c_1^a=\rho c_1^a,\\[1em]
    &\Longrightarrow\quad
    \boxed{\begin{aligned}
        &c_1^{a*}=\frac{1}{\theta+\rho}>1\\[1em]
        &c_2^{b*}=\frac{\rho R}{\theta+\rho}>c_1^{a*}
    \end{aligned}} 
\end{aligned}
}{}

Es decir, resolviendo el problema del planificador benevolente teniendo únicamente en cuenta el consumo que se experimenterá en el primer periodo (el que requiere de liquidez suficiente para hacer frente a las demandas), obtenemos el consumo óptimo que se asignará a los individuos del primer grupo y a los individuos del segundo grupo. De manera que, el supuesto $\rho R>1$ garantiza que $c_2^{b*}>c_1^{a*}$. Además, $c_1^{a*}>1$.

Éste es el resultado crucial: el mecanismo óptimo proporciona a los individuos que necesitan liquidez temprana más de lo que obtendrían simplemente liquidando por sí mismos la inversión. En otras palabras, \textbf{existe un auténtico seguro de liquidez}. ¿Y qué puede funcionar como seguro de liquidez en ausencia de planificador benevolente? Este esprecisamente una de las principales conclusiones: no es necesario conocer las características de cada individuo para alcanzar un resultado óptimo, un intermediario financiero puede servir a este propósito.

¿Cómo? Muy sencillo. Si todos los individuos depositan inicialmente su unidad de riqueza, el banco puede invertir colectivamente los recursos y ofrecer un contrato de depósito, que prometa: un pago $c_1^{a*}$ a quien retire en $T=1$ y un pago $c_2^{b*}$ a quien espere hasta $T=2$. Obsérvese la lógica económica: el banco sabe que únicamente una fracción $\theta$ necesitará retirar fondos anticipadamente. Por ello no necesita mantener toda la riqueza líquida y puede invertir la mayor parte en proyectos de largo plazo y aprovechar la rentabilidad $R$. Lo que está haciendo realmente es compartir el riesgo de liquidez entre todos los individuos. La función económica del banco es, por tanto, proporcionar un seguro frente a necesidades imprevisibles de liquidez. Ahora bien, el mismo contrato bancario genera dos equilibrios de Nash:
\begin{la}
    \item \textbf{Equilibrio eficiente (sin pánico)}. Si todos creen que únicamente los individuos early retirarán fondos en $T=1$, entonces los early retiran y los late esperan. La fracción de retiradas es exactamente \theta. El banco sólo liquida la cantidad necesaria de proyectos. Las inversiones restantes llegan a vencimiento. Todos reciben exactamente: c_1^{a*} y c_2^{b*}. Nadie tiene incentivos a desviarse. Por tanto, éste es un equilibrio de Nash. Es además el equilibrio eficiente y coincide con el óptimo social;
    \item \textbf{Equilibrio ineficiente} (corrida bancaria). Supongamos ahora que cada individuo late cree que todos los demás depositantes retirarán en $T=1$. ¿Qué ocurre? El banco prometió: $c_1^{a*}>1$, pero sus inversiones sólo generan $1$ si se liquidan prematuramente. Por tanto, no puede pagar simultáneamente a todos los depositantes. En este escenario, si demasiados individuos acuden al banco éste tendrá que liquidar todos sus proyectos, desapareciendo la producción futura y liquidándose la porción reservada para repartir en $T=2$. En consecuencia, un individuo late razona: si espero hasta $T=2$, probablemente recibiré $0$, si corro al banco inmediatamente, al menos tengo alguna probabilidad de recibir $c_1^{a*}$. Evidentemente su mejor respuesta es retirar. Pero exactamente el mismo razonamiento lo realizan todos los demás individuos late, lo que acaba conduciendo a que todos corren al banco, intenten retirar y el banco enfrente insolvencia por falta de liquidez.
\end{la}

Como vemos, en este segundo escenario, la expectativa inicial provoca el resultado que temían los depositantes: la creencia genera el hecho y por eso se habla de una profecía autocumplida, que además es también un equilibrio de NASH.

A pesar de la importante contribución, este modelo ha sido sometido a numerosas críticas, muchas relacionadas con los supuestos sobre los que descansa, su simplificación del funcionamiento de la economía y del sistema bancario, con un único banco, donde no existe dinero, así como acerca de las razones por las que se producen los pánicos bancarios y los mecanismos que plantean para evitarlos.\footnote{Muchos autores han trabajado sobre la base de este modelo para sortear estas limitaciones:
\begin{lnum}
    \item GOLDSTEIN y PAUZNER (2005) incorporaron, al modelo presentado, información sobre el estado de la economía, con lo que consiguen que el retorno de la inversión depende de éste y con ello es posible determinar cuál de los dos equilibrios de NASH es más probable. Motivados, principalmente, por críticas acerca de que el modelo no provee de herramientas analíticas para determinar cuál de los dos equilibrios es más probable;
    \item GREEN y LIN (2000) señalan que con contratos de depósito más sofisticados no tendrían lugar pánicos bancarios, frente a la idea de que los bancos están sometidos a la eventual posibilidad de pánicos bancarios únicamente debido a su labor de transformación de plazos.
\end{lnum}

En trabajos como los de GORTON (1985, 1988), JACKLIN y BHATTACHARYA (1988), CHARI y JAGANNATHAN (1988), ALONSO (1996), ALLEN y GALE (1998) y SAMARTÍN (2002), se señala que son los ciclos económicos y la información adquirida por algunos depositantes del banco, con relación a la solvencia del mismo, las razones que explican la retirada masiva de depósitos, mientras que en el modelo presentado estos surgen por un cambio de expectativas, que podría depender de cualquier cosa, es decir ser puramente aleatorias. En cuanto a las soluciones que plantean para evitar los pánicos, BHATTACHARYA y THAKOR (1993) y ENGINEER (1989) señalan que ni el seguro de depósitos ni la suspensión de la convertibilidad son soluciones de primer orden para prevenirlos. CHARI (1989) por su parte diseña un modelo, en el cual la combinación efectiva de requerimientos de reserva, políticas de redescuento y restricciones ocasionales de pagos en efectivo puede llevar a la consecución de un sistema bancario estable, sin necesidad de que exista garantía de depósitos. También DE NICOLÒ (1996), señala que la habilidad de los seguros de depósitos para prevenir los pánicos bancarios no es un hecho empíricamente demostrado y elabora un modelo en el cual los contratos de depósitos no son susceptibles a sufrir pánicos bancarios, independientemente de la existencia o no de un seguro de depósitos o de un mecanismo de suspensión de convertibilidad. Por su parte, DOWD (1993) critica el modelo al considerar que no da una respuesta satisfactoria a por qué los fondos de garantía de depósito deben gestionarse de forma pública, y señala que aseguradoras privadas también podrían acumular los mismos fondos que el gobierno acumularía a través de impuestos.
}


\subsection{Aproximación general: Crisis contenidas y Crisis sistémicas}
Hemos visto cómo determinadas características o comportamientos de los agentes del Sistema Financiero pueden generar propiedades emergentes que ponen en riesgo al mismo Sistema. Sin embargo, para entender la importancia de las crisis financieras debemos adoptar una visión más general del Sistema y sus interconexiones con la economía real. 

En primer lugar, cualquier crisis financiera puede categorizarse como: (i) crisis financiera contenidad; o, (ii) crisis financiera sistémicas. Una crisis financiera contenida se originará y finalizará, normalmente, en los mercados financieros más transparentes. Sus efectos tendrán poca o ninguna traslación a la economía real, y la salud de la economía en su conjunto no se verá mermada. Como podemos intuir, las causas comunes a todas ellas serán la posible existencia de burbujas especulativas que situan el precio de los activos, o los activos de un segmento concreto de los mercados financieros, muy por encima de su teórico valor fundamental. Por ello es frecuente denominarlas como \textit{crisis correctivas} o \textit{shocks bursátiles}. 

Ahora bien, desde la óptica de Política Económica, lo que nos resultará especialmente interesante será la otra categoría: las \textbf{crisis financieras sistémicas}. La razon es que sus efectos desbordarán los límites del Sistema Financiero y desbordarán sobre la economías real, provocando recesiones económicas, aumentos del desempleo y una desaceleración del crecimiento. 

Durante el preludio de una crisis financiera, cuando tanto analistas como policy makers observan un incremento sostenido de la exposición al riesgo y comportamientos/propiedades emergentes que hagan sospechar acerca de la estabilidad del Sistema Financiero, es muy complejo determinar, no solo si se materializará la crisis en cuestión, sino también (y quizás más importante) la categoría en la que esta caerá. Ahora bien, existen determidos patrones que suelen acompañar la antesala de las crisis sistémicas. Entre los más destacados: (i) cambios sustanciales en el volumen de crédito y los precios de los activos; (ii) interrupciones severas en la intermediación financiera y en el suministro de financiamiento externo a diversos actores en la economía; y/o, (iii) problemas a gran escala en los balances (de empresas, hogares, intermediarios financieros y soberanos).\footnote{CLAESSENS y KÖSE (2013) proporcionan un análisis de la literatura existente que analiza los elementos recurrentes de las crisis financieras, y destacan que en la mayoría de los análisis disponibles se observa que las crisis financieras suelen estar precedidas de incrementos acelerados del precio de los activos, periodos de fuerte expansión económica y acumulación de endeudamiento excesivo por parte de uno o varios agentes, que hace a la economía vulnerable a una crisis de confianza. 
}

Como ya conocemos porque razones puede ocasionarse una crisis en alguno de los segmentos del Sector Financiero, debemos entender cómo y por qué una crisis contenida puede pasar a ser una crisis financiera sistémica. 

\subsubsection{Teoría del Ciclo Conducutual: MINSKY (1989)}
Para verlo vamos a presentar en primer lugar la Teoría del Ciclo Conductual de MINSKY, sobre el cual sumaremos diferentes aportaciones que nos permitan entender los canales y mecanismos a través de los cuales se desbordan los efectos de un colapso financiero. 

La visión de MINSKY (1982, p. 101) sobre el capitalismo queda mejor resumida en su aforismo: \textbf{La estabilidad (...) es desestabilizadora}. Para él, \textbf{esta aparente paradoja surge de la forma en que el comportamiento humano influye en la economía al mismo tiempo que la propia economía influye en el comportamiento humano}.\footnote{KING (2013, p. 220) resume claramente este proceso: la cautela conduce a la confianza y luego a la exuberancia, antes de que todo colapse debido a préstamos arriesgados que resulta imposible devolver (por parte de consumidores, propietarios de viviendas o empresas). 
}\textsuperscript{,}\footnote{Como se puede intuir, la teoría de la inversión de KEYNES constituye un punto de apoyo fundamental para MINSKY (1975). Según KEYNES, la inversión empresarial era volátil porque los responsables de la toma de decisiones no podían determinar racionalmente cuál era la mejor inversión posible, dado que desconocían cómo sería el futuro cuando los bienes producidos por una nueva fábrica llegasen finalmente a los consumidores. Por ello, buscan señales, especialmente señales sobre lo que están haciendo otras empresas (KEYNES, p. 156, utiliza la famosa analogía del concurso de belleza: cada participante intenta averiguar quién creerán los demás que es la mejor inversión, en lugar de elegir a lo que el mismo considera). No obstante, MINSKY, que coincidía con KEYNES en la importancia del factor psicológico, difería de él en dos aspectos:
\begin{lnum}
    \item \textbf{Inestabilidad financiera $\neq$ Inestabilidad real}. En primer lugar, a diferencia de KEYNES, MINSKY considera que la inestabilidad económica surge de las decisiones tomadas en Wall Street para financiar nuevas inversiones. Esto también tiene raíces psicológicas y sociales, al igual que la propia decisión de inversión. Cuando Wall Street es optimista y financia una gran cantidad de nuevas inversiones, estas generan beneficios en una economía en auge y los financiadores obtienen elevados rendimientos. Por el contrario, cuando Wall Street se vuelve pesimista y cauteloso, los fondos son difíciles de conseguir y la inversión no se lleva a cabo, lo que provoca pérdidas tanto para las empresas como para los financieros de Wall Street debido al deterioro de las condiciones económicas;
    \item \textbf{Límites de la Política Económica}. En segundo lugar, KEYNES creía que una buena política gubernamental podía evitar una crisis importante y también remediar los problemas económicos. MINSKY era más escéptico. Consideraba que el gobierno podía reducir la probabilidad de una crisis económico-financiera; sin embargo, las tendencias inherentes al comportamiento humano hacían imposible evitar completamente una crisis.
\end{lnum}
} 

Bajo esa lógica, MINSKY considera que los agentes no solo tienen incentivos económicos (apalancamiento), sino también psicológicos y sociales para endeudarse.\footnote{Cuando reina el optimismo, el apalancamiento aumenta tanto para los hogares como para las empresas. Cuanto mayor sea mi apalancamiento, mayores serán proporcionalmente mis ganancias derivadas de la revalorización de los activos. Existe una gran diferencia entre comprar una vivienda aportando un 10\% de entrada o un 50\%. Si mi vivienda costó 300.000 dólares y duplica su valor, obtengo una ganancia de 300.000 dólares. Si aporté un 50\%, es decir, 150.000 dólares, la rentabilidad sobre mi inversión inicial es del 100\%. Si aporté únicamente un 10\%, es decir, 30.000 dólares, mi ganancia es nueve veces superior a la inversión inicial, o un 900\%.
} De esta manera, podemos empezar por distinguir tres tipos de financiación, que deben entenderse como un cíclo evolutivo de tres etapas (al que volveremos más adelante):
\begin{la}
    \item \textbf{Financiación de cobertura} (hedge financing). La financiación de cobertura se produce cuando los flujos de ingresos esperados son suficientes para devolver íntegramente un préstamo. Se trata de una forma segura y conservadora de crédito, como un préstamo para la compra de un automóvil que se amortiza completamente en cinco años;
    \item \textbf{Financiación especulativa} (speculative finance). La financiación especulativa se produce cuando los flujos de ingresos esperados son suficientes para pagar los intereses del préstamo, pero no para amortizar el principal. En este caso, el préstamo debe renovarse o reestructurarse periódicamente;
    \textbf{Financiación Ponzi} (Ponzi finance). La financiación Ponzi se produce cuando los flujos de ingresos esperados ni siquiera son suficientes para pagar los intereses del préstamo. En esta situación, el endeudamiento aumenta continuamente y no existe ninguna posibilidad realista de devolver el préstamo.
\end{la}

Para MINSKY, cabe esperar que durante las expansiones económicas, a medida que personas y empresas buscan mayores rendimientos, incrementan su endeudamiento mediante operaciones de cobertura. No obstante, el optimismo reinante desplaza poco a poco las economías desde una situación en la que predominaba la financiación de cobertura (hedge finance), hacia otra dominada por la financiación especulativa (speculative finance). En algún momento (punto de MINSKY), los prestatarios ya no pueden ni siquiera pagar los intereses de sus préstamos, mucho menos amortizar el principal (financiación PONZI).\footnote{Además, este ciclo se ve acompañado de mayor complacencia regulatoria. La razón es sencilla: a medida que las rentabilidad económicas aumentan y la economía crece a un ritmo mayor, existen razones suficientes para creer que las reivindicaciones del Sistema Financiero respecto a los límites impuestos por la regulación enfrentarán cada vez menor oposición. (Tésis evidentemente validada en los 80's).
} Los niveles de deuda son tan elevados que también las instituciones financieras comienzan a temer las consecuencias de los impagos y endurecen sus criterios de concesión de crédito.

Por tanto, estos cambios en la estructura financiera de empresas y hogares desplazan a las economías desde la estabilidad, hacia la cuasi-estabilidad, y posteriormente hacia la inestabilidad (MINSKY, 1964). En último término, la restricción de crédito acaba generando un sin fin de males económicos: (i) quiebra bancarias por insolvencia; (ii) deterioros sistémicos del balance; y, (ii) pánicos bancarios, en el que las personas acuden masivamente a retirar su dinero antes de que el banco quiebre y sea demasiado tarde para recuperarlo; etc. 

Además, existe un factor crucial que permite explicar porque estas consecuencias son inevitables: \textbf{el elevado nivel de apalancamiento}. Los instrumentos de Gestión del riesgo requieren implícitamente una situación razonablemente sana de la economía. Cuando el apalancamiento es tal que cualquier flexibilización de las condiciones de liquidez pudiera agravar las circunsatancias de los prestatarios/prestamistas, estamos \textbf{atados de manos} ante una debacle financiera. 

\subsubsection{Teoría de la deuda-deflación}
Una vez entendemos cómo podemos llegar a este punto de crisis, independientemente del segemento del Sisema Financiero en el que primero se reflejen las consecuencias de esta dinámica, ¿qué podemos esperar que suceda? ¿cómo se convierte esta crisis de liquidez en una crisis financiera sistémica?

Para verlo vamos a introducir progresivamente lo que se conoce como la Teoría del mecanismo deuda-deflación, inicialmente propuesto por FISHER (1933).\footnote{
Tras la crisis de 1929, cuando se estudia por primera vez de manera moderna la retroalimentación entre la deflación y las tensiones financieras, siendo el precursor de la teoría deuda-deflación FISHER (1933). Es decir, Fisher no limita la direccionalidad de los efectos desde la economía real (deflación) a la financiera (crisis), sino que estudia de manera sistemática un efecto retroalimentación o \textit{feedback}. Como se verá a continuación, trabajos posteriores ampliaron los canales estudiados para entender esta interrelación.
}

\imagenfit{Fig4.png}{Canales del mecanismo deuda-deflación}{BIS (2005)}

En una situación de sobreendeudamiento, es comprensible suponer que los agentes busquen reducir la deuda mediante la redención. Esta redención generalizada, no obstantem conduciría a: (i) tensiones financieras; (ii) contracción de la oferta monetaria; y, (iii) una desaceleración de la velocidad de circulación del dinero. ¿Por qué? La venta masiva de activos generaría una abrupta y sostenida caída de los precios de los activos. Si estos fondos fuesen además utilizados para redimir las deudas, se cancelaría progresivamente la función de banca fraccionaria y, en consecuencia, disminuiría la velocidad de circulación del dinero por una menor liquidez. 

\imagenfit{Fig2.png}{Canal de FISHER}{BIS (2005)}

La combinación de las tensiones financieras junto con el reequilibrio de los mercados monetarios degenera en deflación, que asu vez aumento el valor real del endeudamiento de los agentes: el canal de FISHER cierra el bucle de la deuda-deflación.\footnote{FISHER entiende la deflación como \textit{la raíz de casi todos los males}, y explica así el efecto retroalimentación en su versión desestabilizadora, hablando de una liquidación de deudas que \textit{se derrota a sí misma}: \textit{[S]i el sobreendeudamiento con el que comenzamos era lo suficientemente grande, la liquidación de deudas no puede seguir el ritmo de la caída de precios que provoca. En ese caso, la liquidación se derrota a sí misma. Si bien disminuye la cantidad de dólares adeudados, puede que no lo haga tan rápido como aumenta el valor de cada dólar adeudado. Entonces, el propio esfuerzo de los individuos por disminuir su carga de deudas la aumenta, debido al efecto masivo del pánico por liquidar en aumentar el valor de cada dólar adeudado. Entonces tenemos la gran paradoja que, sostengo, es el principal secreto de la mayoría, si no todas, las grandes depresiones: Cuanto más pagan los deudores, más deben. Cuanto más se inclina el barco económico, más tiende a inclinarse. No está tendiendo a enderezarse, sino que se está volcando. Pero si el sobreendeudamiento no es lo suficientemente grande como para hacer que la liquidación se derrote a sí misma, la situación es diferente y más simple. Es entonces más análoga a un equilibrio estable; cuanto más se balancea el barco, más tenderá a enderezarse.}

PIGOU criticó la visión de FISHER resaltando que, mientras que FISHER destacaba únicamente el papel de los deudores, debían tenerse en cuenta también las ganancias para los acreedores. ¿No deberían simplemente cancelarse las fuerzas opuestas? Años más tarde, TOBIN (1975, 1980) argumentó que el canal de FISHER predomina, porque los deudores generalmente tienen una mayor propensión al gasto que los acreedores. Por tanto, la deflación contaría con un efecto retroalimentación adicional, a través de su impacto desigual en deudores y acreedores y por tanto de su efecto en la demanda agregada: \textit{La agregación no importaría si pudiéramos estar seguros de que las propensiones marginales al gasto desde la riqueza fueran iguales para acreedores y deudores. Pero si la propensión al gasto fuera sistemáticamente mayor para los deudores, aunque sea en pequeña medida, el efecto Pigou sería arrasado por este efecto Fisher}. TOBIN (1975) también sugirió que la deflación reduce el valor del capital y la equidad, haciendo que tanto el consumo como la inversión sean menos atractivos.
} Años más tarde, TOBIN amplío las aportaciones de FISHER argumentando que, como los deudores tenían una mayor propensión de gasto que los acreedores, debía añadirse un nuevo bucle de retroalimentación vía gasto: conecta el gasto de nuevo con la deflación vía redistribución (ya incluido en la figura). 

MINSKY (1982) amplío la Teoría deuda-deflación más allá del mercado monetario, incorporando el mercado de activos financieros. El autor considera que el efecto de las ventas masivas en los precios de los activos: (i) dan lugar a nuevas ventas masivas; y, (ii) contribuyen a empeorar la deflación, a través de un efecto riqueza negativo.

\imagenfit{Fig3.png}{Canales de MINSKY}{BIS (2005)}

MINSKY enfatizó en el efecto circular de estas ventas masivas, pues la caída en los precios de los activos podrían ser tan grande que la venta de activos no pudiera liquidar los fondos suficientes como para redimir las deudas. Esta situación generaría ventas adicionales, reforzando el proceso de caídas en los precios y pérdida de patrimonio neto del deudor. En otras palabras, cuando la venta por angustia reduce los precios de los activos, las pérdidas resultantes agravan el endeudamiento y pueden llevar a más ventas por angustia. Por otro lado, y en cuanto al segundo canal, argumenta que la caída en los precios de los activos refuerza la deflación a través del gasto agregado. 

Sin embargo, a diferencia de TOBIN, esta repercusión sobre el gasto no se produciría por motivos estrictamente monetarios (la propia deflación y sus efectos redistributivos) sino por la caída de valor del ahorro de los hogares. Es decir, un efecto riqueza negativo que afectaría tanto a deudores como a acreedores.\footnote{Sobre esto último, cuando los compromisos de pago no se pueden cumplir incluso tras la liquidación de activos, las pérdidas de los deudores se transfieren a sus acreedores mediante un incumplimiento o refinanciamiento de los contratos. Tales pérdidas de capital inducidas resultan en una contracción adicional del consumo e inversión más allá de la causada por la disminución inicial del ingreso.
}

Por último, una tercera aportación realizada por BERNANKE (1983, en linea con su propuesta del \textbf{acelerador financiero}) centra la atención en la intermediación bancaria. Las garantías y el patrimonio neto son márgenes de seguridad que se incorporan explícita (en el caso de las primeras) o implícitamente (en el caso del segundo) en los acuerdos de préstamo y normalmente protegen al sistema bancario de las consecuencias de la caída de los precios.\footnote{\textit{Los problemas bancarios de 1930-33 perturbaron el proceso de asignación de crédito al crear cambios grandes e imprevistos en los canales de flujo de crédito. [...] más las quiebras reales, forzaron una contracción del papel del sistema bancario en la intermediación de crédito.} BERNANKE (1983)
} 

Sin embargo, caídas de precios lo suficientemente grandes afectan al sistema bancario hacen que una situación de sobre-endeudamiento y deflación se refuerce a través de la intermediación bancaria: 
\begin{la}
    \item \textbf{Solvencia: Racionamiento de crédito}. En primer lugar, los impagos afectan a la solvencia del balance de las entidades, que se ven obligadas a racionar el crédito para evitar entrar en causa de quiebra;
    \item \textbf{Costes de información: Racionamiento de crédito}. En segundo lugar, estos impagos implican también incrementos en los costes de intermediación en lo referido a la recopilación de información. De nuevo, por lógica económica, ello da lugar a un refuerzo del racionamiento del crédito. 
\end{la}

En suma, la contracción del crédito inevitablemente conduce a una caída del gasto agregado, vía disminución de la oferta monetaria, lo que agrava las tensiones deflacionistas, y por tanto un efecto circular o recursivo que se suma a los formulados por FISHER y MINSKY. De hecho, esta contracción del crédito agrava doblemente su situación por el empeoramiento de las condiciones dentro del propio sistema financiero. La razón es sencilla: si la oferta de crédito es limitada, aunque se ponga freno al sobre-endeudamiento, la única vía que se deja a los deudores para obtener liquidez para redimir sus deudas es liquidando activos, lo que acaba agravando los efectos observados en los círculos de FISHER y MINSKY. De manera sintetizada:

\imagenfit{Fig4.png}{Canales del mecanismo deuda-deflación}{BIS (2005)}

¿Cuándo está espiral degenerada? De nuevo, la respuesta posiblemente más acertada es la que nos proporciona MINSKY (1983). Este estancamiento económico continúa hasta que la deuda se reduce a niveles manejables, los agentes pueden volver a gastar con mayor libertad y las instituciones financieras se recuperan y sanean su balance, de modo que pueden volver a conceder crédito sin restricciones importantes... Y vuelta a empezar.








\clearpage
\section{Conclusión} 


\subsection{Recapitulación}


\subsection{Extensiones y relación con otras partes del Temario}



\subsection{Opinión}

\subsubsection{Idea final (Salida o cierra)}

\clearpage
\pagenumbering{gobble}
\MyBibliography

\href{https://www.hbs.edu/behavioral-finance-and-financial-stability/data/Pages/global.aspx?utm_source=chatgpt.com}{\textit{Data -- Global Crises Data by Country}. Harvard Business School}


% ANEXOS
\clearpage
\anexo{\textbf{Preguntas Test}}
%\input{\TemaCode/questions.tex} 




\clearpage
\anexo{Ciclo conductual del Sistema Financiero: Una historia de interrelación regulatoria y exuberancia irracional (MINSKY)}\label{anx:ciclo_minsky}

Todo lo aquí presentado es una traducción y ligera adaptación del paper: \href{https://sabeconomics.org/journal/RePEc/beh/JBEPv1/articles/JBEP-2-1-5.pdf}{\textit{Hyman Minsky and behavioral finance}. PRESSMAN (2018)}

La visión de Minsky (1982, p. 101) sobre el capitalismo queda mejor resumida en su conciso aforismo: \textbf{La estabilidad (...) es desestabilizadora}. Esta aparente paradoja surge de la forma en que el comportamiento humano influye en la economía al mismo tiempo que la propia economía influye en el comportamiento humano. KING (2013, p. 220) resume claramente este proceso: la cautela conduce a la confianza y luego a la exuberancia, antes de que todo colapse debido a préstamos arriesgados que resulta imposible devolver (por parte de consumidores, propietarios de viviendas o empresas).

Para MINSKY, una economía estable lleva a las personas a esperar estabilidad en el futuro y las impulsa a asumir mayores riesgos con la esperanza de obtener una mayor rentabilidad de su dinero. A esto debe añadírsele la tendencia al comportamiento gregario que tienen los agentes (\textbf{efecto rebaño});\footnote{Existen numerosas razones que se citan para naturalizar este efecto. Destacamos las siguientes:
\begin{lnum}
    \item \textbf{Bounded rationality}. En primer lugar, hasta que no realizamos una acción y observamos sus consecuencias, resulta difícil saber cuáles serán los resultados de las distintas alternativas posibles, por lo que es imposible seguir plenamente los dictados de la racionalidad económica;
    \item \textbf{Presión social y comportamiento adaptativo}. En segundo lugar, existe una presión social hacia la conformidad y los seres humanos son animales sociales. Además, no hay que olvidar que el comportamiento gregario tiene valor adaptativo: por eso los animales viven en manadas y los peces nadan en bancos. Las personas con una tendencia a seguir al grupo tuvieron mayores probabilidades de sobrevivir en el pasado remoto y transmitieron esa inclinación a sus descendientes;
    \item \textbf{Carga de las consecuencias: incentivo}. Finalmente, siempre es más fácil equivocarse cuando todos los demás también se equivocan. Si todos actúan de la misma manera, cuando las cosas salen mal resulta sencillo eludir responsabilidades afirmando que \textit{nadie lo sabía}.
\end{lnum}
} como resultado de esta propensión conductual, cada vez más individuos asumen riesgos crecientes porque observan que otros hacen lo mismo. Cabe resaltar, no obstante, que esta conducta \textbf{no es, \textit{per se}, irracional}. Pensémoslo. Este tipo de toma de decisiones tiene sentido porque la rentabilidad de cualquier inversión dependerá del estado general de la economía, o de la demanda agregada, cuando la nueva inversión se traduzca en bienes y servicios puestos a la venta. Cuando muchas empresas invierten y contratan más trabajadores, y cuando estos gastan sus ingresos, la mayoría de las nuevas inversiones resultarán rentables. Por el contrario, si las demás empresas no están invirtiendo ni contratando, la inversión realizada por una empresa concreta perderá dinero porque sus ventas serán bajas.

En definitiva, por una cosa o la otra, lo natural será que los agentes se acaben comportando como los demás; y, \textbf{sobretodo, elegirán en función de lo que creen que harán otros}.

Como se puede intuir, la teoría de la inversión de KEYNES constituye un punto de apoyo fundamental para MINSKY (1975).\footnote{Según KEYNES, la inversión empresarial era volátil porque los responsables de la toma de decisiones no podían determinar racionalmente cuál era la mejor inversión posible, dado que desconocían cómo sería el futuro cuando los bienes producidos por una nueva fábrica llegasen finalmente a los consumidores. Por ello, buscan señales, especialmente señales sobre lo que están haciendo otras empresas (KEYNES, p. 156, utiliza la famosa analogía del concurso de belleza: cada participante intenta averiguar quién creerán los demás que es la mejor inversión, en lugar de elegir a lo que el mismo considera).
} Aspectos sociales y psicológicos que, como se sabe, llevaron a KEYNES a concluir que el capitalismo era inherentemente inestable. No obstante, MINSKY, que coincidía con Keynes en este punto, difería de él en dos aspectos:
\begin{lnum}
    \item \textbf{Inestabilidad financiera $\neq$ Inestabilidad real}. En primer lugar, a diferencia de KEYNES, MINSKY considera que la inestabilidad económica surge de las decisiones tomadas en Wall Street para financiar nuevas inversiones. Esto también tiene raíces psicológicas y sociales, al igual que la propia decisión de inversión. Cuando Wall Street es optimista y financia una gran cantidad de nuevas inversiones, estas generan beneficios en una economía en auge y los financiadores obtienen elevados rendimientos. Por el contrario, cuando Wall Street se vuelve pesimista y cauteloso, los fondos son difíciles de conseguir y la inversión no se lleva a cabo, lo que provoca pérdidas tanto para las empresas como para los financieros de Wall Street debido al deterioro de las condiciones económicas;
    \item \textbf{Límites de la Política Económica}. En segundo lugar, KEYNES creía que una buena política gubernamental podía evitar una crisis importante y también remediar los problemas económicos. MINSKY era más escéptico. Consideraba que el gobierno podía reducir la probabilidad de una crisis económico-financiera; sin embargo, las tendencias inherentes al comportamiento humano hacían imposible evitar completamente una crisis.
\end{lnum}

Habida cuenta de lo expuesto, podemos adentrarnos en el núcleo teórico del Ciclo propuesto por MINSKY.

$\mathbf{\text{Ciclo Financiero de Minsky}}$

Bajo esa lógica conductual, los agentes tienen no solo incentivos económicos (apalancamiento), sino también psicológicos y sociales para endeudarse.\footnote{Cuando reina el optimismo, el apalancamiento aumenta tanto para los hogares como para las empresas. Cuanto mayor sea mi apalancamiento, mayores serán proporcionalmente mis ganancias derivadas de la revalorización de los activos. Existe una gran diferencia entre comprar una vivienda aportando un 10\% de entrada o un 50\%. Si mi vivienda costó 300.000 dólares y duplica su valor, obtengo una ganancia de 300.000 dólares. Si aporté un 50\%, es decir, 150.000 dólares, la rentabilidad sobre mi inversión inicial es del 100\%. Si aporté únicamente un 10\%, es decir, 30.000 dólares, mi ganancia es nueve veces superior a la inversión inicial, o un 900\%.
} Por tanto, cabe esperar que durante las expansiones económicas, a medida que personas y empresas buscan mayores rendimientos, incrementan su endeudamiento. Conforme aumentan las cargas de deuda, resulta más difícil devolver el dinero prestado. 

En algún momento, los prestatarios ya no pueden ni siquiera pagar los intereses de sus préstamos, mucho menos amortizar el principal. En este contexto, los niveles de deuda son tan elevados que también las instituciones financieras comienzan a temer las consecuencias de los impagos y endurecen sus criterios de concesión de crédito. Lo que habrá la puerta a un sin número de azotes con proyección nominal: (i) quiebra bancarias por insolvencia; (ii) deterioros sistémicos del balance; y, (ii) pánicos bancarios, en el que las personas acuden masivamente a retirar su dinero antes de que el banco quiebre y sea demasiado tarde para recuperarlo; etc. En definitiva, el ciclo alcanza un punto en el que una severa contracción económica deviene ineitable; lo que, a su vez, agrava las cargas de deuda debido a la caída de los ingresos. Su traslación real es ahora también inevitable: (i) quiebras empresariales; (ii) pesimismo respecto al futuro; (iii) reducción del gasto. En definitiva: una fuerte desaceleración económica, e incluso puede desembocar en una crisis económica. 

Este estancamiento económico continúa hasta que la deuda se reduce a niveles manejables, los agentes pueden volver a gastar con mayor libertad y las instituciones financieras se recuperan y sanean su balance, de modo que pueden volver a conceder crédito sin restricciones importantes... Y vuelta a empezar.

Como vemos, gran parte del problema radica en \textbf{la forma en que la psicología influye sobre el comportamiento}. Así que, si parece tan lógico, ¿por qué no se tuvo, al menos hasta cierto punto, en cuenta? MINSKY (1975, 1982, 1986) identifica varias razones por las que los economistas han descuidado esta característica del capitalismo:
\begin{lnum}
    \item \textbf{Estabilidad: La Edad de oro}. En primer lugar, las políticas económicas y la regulación gubernamental de las instituciones financieras estabilizaron la economía mundial durante la mitad del siglo XX. Esto llevó a creer que la estabilidad se había convertido en una característica permanente del capitalismo y generó desinterés por estudiar su inestabilidad inherente;
    \item \textbf{Consenso teórico y \textit{naturalismo} teórico}. En segundo lugar, los economistas llegaron a esperar que, dado que comprendían el funcionamiento de la economía, el gobierno sería capaz de rescatarla frente a cualquier crisis, mitigando los daños y garantizando que cualquier problema fuese de corta duración;
    \item \textbf{Efecto rebaño: retroalimentación teórica}. En tercer lugar, el comportamiento gregario o social llevó tanto a economistas como a financieros a creer que el capitalismo era estable porque todos los demás también lo creían y, encima, FAMA (1970) lo validaba con su Hipótesis de los Mercados Eficientes. Fascinados por la idea de la racionalidad individual y, al mismo tiempo, comportándose de forma gregaria, los economistas llegaron a aceptar la hipótesis de los mercados eficientes. Y, además de desviar la atención de la inestabilidad inherente al capitalismo, la hipótesis de los mercados eficientes tuvo una importante implicación en términos de política económica: se consideró innecesaria la regulación gubernamental de las finanzas, ya que Wall Street podía regularse eficazmente por sí mismo.
\end{lnum}

El resultado: \textbf{un esfuerzo concertado en Estados Unidos para desregular la banca y las finanzas a partir de la década de 1980} (la desregulación, Europa 80-90, Reino Unido inicia en 1986 (\textit{Big Bang})). No obstante, rechazando la visión convencional de unas finanzas guiadas por inversores racionales y considerando, en cambio, que la psicología de los inversores es volátil y está fuertemente influida por el comportamiento ajeno, MINSKY intentó explicar cómo el comportamiento humano real genera problemas económicos y financieros. Distinguiendo tres tipos de financiación, que deben entenderse como un cíclo evolutivo de tres etapas (al que volveremos más adelante):
\begin{la}
    \item \textbf{Financiación de cobertura} (hedge financing). La financiación de cobertura se produce cuando los flujos de ingresos esperados son suficientes para devolver íntegramente un préstamo. Se trata de una forma segura y conservadora de crédito, como un préstamo para la compra de un automóvil que se amortiza completamente en cinco años;
    \item \textbf{Financiación especulativa} (speculative finance). La financiación especulativa se produce cuando los flujos de ingresos esperados son suficientes para pagar los intereses del préstamo, pero no para amortizar el principal. En este caso, el préstamo debe renovarse o reestructurarse periódicamente;
    \textbf{Financiación Ponzi} (Ponzi finance). La financiación Ponzi se produce cuando los flujos de ingresos esperados ni siquiera son suficientes para pagar los intereses del préstamo. En esta situación, el endeudamiento aumenta continuamente y no existe ninguna posibilidad realista de devolver el préstamo.
\end{la}

MINSKY pensaba que el optimismo desplazaba a las economías desde una situación en la que predominaba la financiación de cobertura (hedge finance), hacia otra dominada por la financiación especulativa (speculative finance) y, finalmente, hacia una situación en la que predominaba la financiación Ponzi. Estos cambios en la estructura financiera de empresas y hogares también desplazan a las economías desde la estabilidad hacia una cuasi-estabilidad y posteriormente hacia la inestabilidad (MINSKY, 1964). 

$\mathbf{\textbf{Desregulación}:}\qquad \text{Implicaciones Económicas}$

¿Se tuvieron en cuenta las aportaciones de MINSKY? La respuesta es clara y evidente: no. El proceso de desregulación puede resumirse cómo sigue.
En primer lugar, la Garn–St. Germain Depository Institutions Act de 1982 desreguló las cooperativas de crédito (credit unions) y las asociaciones de ahorro y préstamo (savings and loans), y permitió la existencia de bancos no bancarios (non-bank banks), como compañías hipotecarias, prestamistas de día de pago (payday lenders) y fondos de cobertura (hedge funds) que captan dinero de individuos adinerados y utilizan esos fondos para conceder préstamos. Estos bancos en la sombra (shadow banks) crecieron rápidamente, sin verse limitados por muchas de las restricciones impuestas a los bancos tradicionales. Sus depositantes o inversores obtenían mayores rendimientos porque estas entidades concedían préstamos arriesgados a tipos de interés elevados. La desregulación de estas instituciones incrementó la presión para desregular también a los grandes bancos comerciales, de modo que pudieran competir con los bancos en la sombra, sujetos a una regulación mucho más laxa.

La Riegel-Neal Interstate Bank Efficiency Act de 1994 eliminó las restricciones a la banca interestatal. Esto impulsó el surgimiento de megainstituciones financieras que se volvieron demasiado grandes para quebrar (too big to fail). Sabiendo que el gobierno tendría que rescatarlas si corrían peligro de colapsar, los grandes bancos pudieron asumir riesgos aún mayores, ya que la principal consecuencia negativa de una política crediticia agresiva (la quiebra) quedaba mitigada por la promesa implícita de un rescate gubernamental.

Finalmente, en noviembre de 1999, el presidente Clinton firmó la ley Gramm-Leach-Bliley, derogando así la Glass-Steagall Act de 1933. Esta reforma del New Deal había establecido el seguro de depósitos y había limitado los riesgos que los bancos comerciales podían asumir con los depósitos asegurados. Bajo Glass-Steagall, un banco podía aceptar depósitos y conceder préstamos, o bien dedicarse a la venta de valores financieros, pero no podía realizar ambas actividades simultáneamente. Si un banco concedía un préstamo, debía mantenerlo en su balance, ya que tenía prohibido venderlo. Además, Glass-Steagall exigía que los bancos ofrecieran únicamente hipotecas convencionales a tipo fijo.

La derogación de Glass-Steagall fomentó la concesión de préstamos de alto riesgo. Los bancos podían aprobar préstamos, agrupar muchos de ellos en paquetes y vender posteriormente esos paquetes en el mercado. Su único incentivo era prestar. Cada préstamo generaba una comisión para el banco y, puesto que posteriormente se vendía, el banco no tenía que preocuparse por la capacidad del prestatario para devolverlo. Esto llevó a los bancos a promover hipotecas subprime, préstamos sin verificación de ingresos, \textit{préstamos mentirosos} (liar loans, en los que se aconsejaba a los prestatarios falsear información sobre ingresos y patrimonio para obtener una hipoteca), préstamos con tipos iniciales reducidos (teaser rates) y préstamos sin entrada inicial.

A los prestatarios se les decía que los precios de la vivienda siempre subirían y que podrían refinanciar sus hipotecas antes de que expirara el período promocional y aumentaran bruscamente sus cuotas mensuales.

Posteriormente, estos préstamos se empaquetaban en títulos respaldados por hipotecas (mortgage-backed securities) y se vendían a inversores desprevenidos, que carecían tanto del tiempo como de la disposición necesarios para examinar todas las hipotecas incluidas en el paquete, además de los conocimientos técnicos requeridos para hacerlo. En lugar de ello, los inversores confiaban en que la diversificación proporcionaría seguridad. Algunas hipotecas individuales podrían acabar en ejecución hipotecaria, pero el conjunto del paquete probablemente seguiría generando pagos de intereses mientras la mayoría de los prestatarios continuaran pagando sus cuotas.

Los inversores se sentían aún más tranquilos debido a la máxima calificación AAA otorgada a estos títulos hipotecarios, sin saber que las agencias de calificación eran remuneradas por las mismas empresas que empaquetaban y vendían los activos, por lo que tenían incentivos económicos para concederles la mejor nota posible. También aquí entró en juego el comportamiento gregario: todos creían que los títulos respaldados por hipotecas eran seguros porque los demás también lo creían y los compraban.

MINSKY (1982, 1986) se opuso firmemente a la desregulación financiera y temía sus consecuencias. Su argumento contra la desregulación se basaba en su visión de la psicología y el comportamiento humano. MINSKY consideraba que las expectativas eran irracionales y cíclicas, y que inducían cambios en los comportamientos de préstamo y endeudamiento que terminaban generando una crisis económica.

Una crisis financiera vuelve conservadores a los bancos en materia de crédito. Sin embargo, con el paso del tiempo las cosas cambian. A medida que los préstamos se devuelven, aumenta el optimismo. Los prestamistas llegan a convencerse de que prestar es una actividad segura y rentable. Los análisis tradicionales del riesgo comienzan a considerarse excesivamente conservadores porque no ha habido problemas financieros recientes. En consecuencia, las instituciones financieras conceden préstamos cada vez más arriesgados y animan a los prestatarios a asumir mayores niveles de deuda (MINSKY, 1982, p. 123). El optimismo también lleva a consumidores y empresas a endeudarse más y el comportamiento gregario no hace más que reforzar este círculo.

Para el sector bancario existe además una razón práctica para comportarse de forma gregaria. Si un único banco experimenta dificultades financieras, puede permitirse que quiebre. Pero cuando muchas entidades están en riesgo de colapsar, el gobierno se ve obligado a intervenir para evitar una crisis económica. Además, dado que el gobierno asegura los depósitos bancarios (hasta 250.000 dólares por persona y por banco), los contribuyentes terminan asumiendo el coste tanto si los bancos son rescatados como si quiebran y arrastran a la economía en su caída.

$\textbf{Clasificación de MINSKY}: \qquad \text{¿Puede explicar la Teoría de MINSY la crisis (2008)?}$

La clasificación tripartita que hemos visto antes puede aplicarse fácilmente a la crisis financiera y económica de comienzos del siglo XXI. Los prestatarios de cobertura eran aquellos que tenían una hipoteca convencional que amortizaban completamente en treinta años. Los prestatarios especulativos poseían préstamos a tipo variable con un tipo inicial reducido (teaser rate). Esto les permitía acceder a hipotecas más grandes y viviendas de mayor valor, pero generaba problemas de flujo de caja cuando el préstamo se reajustaba a un tipo de interés más elevado; en consecuencia, los prestatarios tenían que refinanciar periódicamente. Esto no representaba un problema mientras los precios de la vivienda continuaran aumentando. Las hipotecas exóticas cuyo principal crecía con el tiempo constituyen buenos ejemplos de financiación Ponzi.

La transición desde la financiación de cobertura hacia la financiación especulativa y posteriormente hacia la financiación Ponzi generó una enorme expansión del mercado inmobiliario. La facilidad para obtener crédito impulsó los precios de la vivienda y reforzó la creencia de que dichos precios solo podían aumentar (ignorando convenientemente la fuerte caída de los precios inmobiliarios durante la Gran Depresión). El auge se transformó en una burbuja especulativa a medida que un número creciente de hipotecas era concedido a personas incapaces de devolverlas. Esto se vio favorecido por el fuerte aumento de los precios de los activos, que podían utilizarse como garantía para nuevos préstamos.
 
Los problemas comenzaron cuando los precios de la vivienda primero se estabilizaron y posteriormente descendieron en 2007. Esto significó que los propietarios ya no podían refinanciar sus hipotecas ni extraer capital de sus viviendas para incrementar su consumo. El colapso del mercado de títulos de tasa subastada en febrero de 2008, puso de manifiesto la creciente reticencia de los inversores a mantener hipotecas o títulos respaldados por hipotecas.\footnote{El mercado de títulos de tasa subastada, o auction-rate securities, era un mercado en el que los paquetes de hipotecas eran subastados mensualmente al mejor postor
} Además, dejó atrapados a numerosos inversores. Se les había prometido que estas inversiones les proporcionarían liquidez. Sin embargo, a partir de febrero de 2008 ya no podían vender los activos que poseían porque no existían compradores para dichos títulos (LEE, 2008).

La quiebra de Lehman Brothers (2008), fuertemente apalancada en instrumentos de titulización, constituyó la proverbial gota que colmó el vaso. Las instituciones financieras comenzaron a acumular liquidez y se volvieron reacias a conceder préstamos. Incapaces de refinanciarse, los propietarios dejaron de pagar sus hipotecas. Los precios de la vivienda cayeron, y tanto los bancos como los hogares se encontraron en graves dificultades financieras. Con menos crédito y menos gasto, la economía entró en una profunda recesión.

$\textbf{Conclusiones}:\qquad \text{Implicaciones de Política Económica}$

Para MINSKY, la inestabilidad inherente del capitalismo y la posibilidad permanente de una crisis financiera tenían importantes implicaciones de política económica.

MINSKY insistió continuamente en que el Estado debía regular las instituciones financieras. La regulación era necesaria para reducir la especulación y la asunción excesiva de riesgos, así como para impedir que la euforia especulativa condujera a ciclos de auge y caída que terminaran dañando la economía. MINSKY (1982) consideraba que las regulaciones financieras del New Deal, como la Glass-Steagall Act, favorecían la estabilidad económica. Estas regulaciones estaban diseñadas para evitar que las instituciones financieras asumieran riesgos excesivos y especularan con dinero depositado por los ahorradores y asegurado por el gobierno. Según MINSKY, el desmantelamiento de estas regulaciones conduciría a una crisis financiera semejante a la Gran Depresión. Las presiones ejercidas por el sector financiero para lograr la desregulación durante los períodos de bonanza económica debían ser resistidas. Sin embargo, también reconocía la dificultad de hacerlo. Pensaba que las presiones psicológicas y políticas propias de los buenos tiempos terminarían provocando una reducción de la regulación financiera, preparando así el terreno para la siguiente crisis.

Coincidiendo con KEYNES, MINSKY (1982) sostenía que el Estado podía reducir tanto la frecuencia como la gravedad de las crisis económicas. Para ello, los bancos centrales debían actuar como prestamistas de última instancia durante las crisis, tal como hizo la Federal Reserve System durante la Gran Recesión. De lo contrario, la quiebra de una institución financiera pondría en peligro a otras instituciones, haciéndolas aún más reacias a prestar. Además, temiendo perder sus ahorros, los ciudadanos podrían acudir masivamente a los bancos para retirar su dinero; dado que la mayor parte de los depósitos bancarios se presta a terceros, los bancos no pueden satisfacer simultáneamente todas las retiradas y podrían acabar quebrando. Con menos crédito y menos gasto, el resultado inevitable sería una fuerte contracción económica. No obstante, MINSKY reconocía que cuanto más eficaces fueran los bancos centrales en evitar una crisis, más creerían todos que una nueva crisis era imposible, lo que conduciría a mayores niveles de apalancamiento y, en consecuencia, a un mayor riesgo de futuras crisis.

MINSKY también pensaba que eran necesarias respuestas gubernamentales pragmáticas para limitar los daños derivados de una crisis financiera. Defendía una política fiscal keynesiana activa, en la que el gobierno actuara como empleador de última instancia. MINSKY (1986, cap. 13) modeló su propuesta de empleo sobre los programas de la Works Progress Administration (WPA) y del Civilian Conservation Corps (CCC). Quería que el gobierno ofreciera empleos públicos remunerados al salario mínimo a cualquier persona incapaz de encontrar trabajo en el sector privado. También consideraba que esto incentivaría al sector privado a ofrecer salarios superiores al mínimo legal. Todo ello contribuiría a estimular el gasto de los consumidores y el crecimiento económico, mientras que unos salarios más elevados aliviarían las cargas de deuda que soportan los hogares (véase WRAY, 2016).

MINSKY (1986) también señaló que el crecimiento del sector público como proporción de la economía ayudaba a estabilizar la actividad económica porque la inversión pública (infraestructuras, edificios, investigación y desarrollo) es más estable que la inversión empresarial, y porque un sector público más amplio incorpora estabilizadores automáticos más potentes, como los seguros de desempleo. También aquí temía que el éxito económico y el bajo desempleo generaran presiones para reducir el tamaño y la influencia del Estado con el fin de favorecer la iniciativa privada.

Sin embargo, aunque el Estado puede reducir la probabilidad y la gravedad de las crisis, no puede eliminarlas completamente. Debido a las tendencias psicológicas y conductuales descritas anteriormente, MINSKY pensaba que las economías capitalistas siempre estarían sujetas a algún tipo de crisis. El problema, según él, es que un buen desempeño económico vuelve a las personas complacientes y excesivamente confiadas; llegan a creer que su éxito se debe exclusivamente a su propia brillantez y que el gobierno obstaculiza la obtención de logros aún mayores. Las presiones de Wall Street sobre los políticos para reducir la regulación conducen a una menor supervisión y a una mayor especulación, especialmente cuando todos pueden señalar la ausencia de crisis durante un largo período de tiempo. Esto abre la puerta a la siguiente crisis.

Su análisis de las causas psicológicas y sociales de las crisis financieras convierte a MINSKY en uno de los padres fundadores de las finanzas conductuales (behavioral finance), una rama de la economía y las finanzas que estudia las decisiones financieras no como el resultado de una racionalidad perfecta, sino como consecuencia de cómo se comportan realmente las personas en el mundo real.

La obra del premio Nobel SHILLER (2016), una de las principales figuras de este campo, se apoya firmemente en los fundamentos establecidos por MINSKY. SHILLER (2003) ha criticado la hipótesis de los mercados eficientes desarrollada por FAMA y ha identificado numerosas formas en que los mercados financieros se comportan de manera irracional. A diferencia de MINSKY, SHILLER (2012) ha intentado diseñar mecanismos e incentivos que fomenten comportamientos más racionales por parte de las instituciones financieras. Sin embargo, si los mercados financieros son inherentemente irracionales porque la irracionalidad forma parte de la naturaleza humana, entonces la conclusión de política económica de MINSKY sigue siendo válida: son necesarias regulaciones gubernamentales estrictas y rigurosamente aplicadas para mantener bajo control al sistema financiero y evitar que las finanzas provoquen una gran crisis económica.


\anexo{Riegos de Basilea: Instrumentos de Gestión de riesgos}\label{anx:riesgos_basilea}
A lo largo del tema hemos tratado los tres riesgos más importantes de entre los siete que señala el Comité de Basilea. Aquí presentaremos brevemente los otros cuatro tipos de riesgos que deben ser adecuadamente identificados, cuantificados y gestionados: operacional, de concentración, de contraparte, de apalancamiento excesivo, de modelo, de contagio y sistémico.

$\textbf{Riesgo operacional: Visión de conjunto}$

El cuarto es el riesgo operacional: pérdidas derivadas de procesos internos, personas, sistemas o eventos externos, incluyendo continuidad de negocio, tecnologías de información y resiliencia operativa. El riesgo operacional se observa mediante pérdidas internas, eventos externos, indicadores de control, continuidad operativa y \textbf{escenarios de estrés}. 

El riesgo operacional es aquél que puede provocar pérdidas como resultado de errores humanos en la propia entidad (en la gestión y administración de cuentas o en la documentación de contratos, en la operatividad, controles inadecuados, incumplimiento de normativa, incumplimientos de contratos), fraude externo (uso fraudulento de tarjetas de crédito, robos, atracos, etc.) o interno (apropiación indebida de activos de la organización de forma intencional, actividades no autorizadas), procesos internos inadecuados o defectuosos, fallos tecnológicos que, así como daños a activos físicos como resultado de accidentes, siniestros naturales o siniestros provocados. Incluye por lo tanto el riesgo legal y excluye el riesgo estratégico o de negocio y/o el riesgo reputacional.


El riesgo operacional es aquél que puede provocar pérdidas como resultado de errores humanos en la propia entidad (en la gestión y administración de cuentas o en la documentación de contratos, en la operatividad, controles inadecuados, incumplimiento de normativa, incumplimientos de contratos), fraude externo (uso fraudulento de tarjetas de crédito, robos, atracos, etc.) o interno (apropiación indebida de activos de la organización de forma intencional, actividades no autorizadas), procesos internos inadecuados o defectuosos, fallos tecnológicos que, así como daños a activos físicos como resultado de accidentes, siniestros naturales o siniestros provocados. Incluye por lo tanto el riesgo legal y excluye el riesgo estratégico o de negocio y/o el riesgo reputacional.


El IMF codifica esta observación mediante los Financial Soundness Indicators, que estandarizan indicadores de capitalización, calidad de activos, rentabilidad, liquidez, sensibilidad al riesgo de mercado y exposición sectorial para evaluar solidez financiera. Como se ha comentado, la labor de los supervisores va más allá de la comprobación del cumplimiento de la normativa. Un ejercicio que evidencia esto son los tests de estrés o pruebas de resistencia. Germinalmente, se trataba de herramientas internas de las propias entidades para ser conocedoras de sus riesgos (generalmente, de crédito, de mercado y de liquidez). Con el paso del tiempo, se ha avanzado tanto a hacer obligatorio su empleo por parte de las entidades, como a que los propios supervisores los lleven a cabo en el marco de su labor. Además, recientemente, y pese a la dificultad computacional y de modelización que implica, han comenzado a proliferar los tests de estrés sistémicos, centrados en las interconexiones dentro del sistema bancario o incluso incorporando también a los sectores NBFI. 

Por último, existen test de estrés ad-hoc, tales como test de estrés del Covid-19, del cambio climático o de la resiliencia operativa digital. Los escenarios que se analizan incluyen evoluciones futuras hipotéticas de variables de la economía real y sector financiero, tanto de la economía doméstica como extranjera, tales como crecimiento del PIB, tasa de desempleo, consumo, tipos de interés a corto y largo plazo, precio de distintos activos, tipos de cambio de divisas o medidas de volatilidad. 

Los escenarios de estrés han de ser severos pero plausibles. Su severidad refleja la aversión al riesgo de las autoridades y su habilidad para intervenir, donde el escenario ha de ser creíble. Así, tal y como recomiendan PARLATORE y PHILIPPON (2022), los supervisores deberían concentrarse en los factores de riesgo que crean pueden ser relevantes en el futuro y donde los costes asociados a su intervención no sean exorbitantes. PARLASCA (2021) resalta la inconsistencia temporal del supervisor, donde pone de manifiesto las ventajas de ser exigente y rígido ex-ante para incentivar un comportamiento prudente, mientras que la estrategia por parte del supervisor ex-post, una vez que se ha producido la crisis o la situación de estrés, es inversa de manera que prefiere ser reflexivo y flexible para apaciguar a los mercados sobre la situación económica y financiera del banco. 

No todos los ejercicios de tests de estrés que se realizan se publican ya que su publicación al mercado puede tener consecuencias no deseadas. La comunicación de los resultados de los tests de estrés fomenta la disciplina de mercado y puede promover que las entidades se apresuren a aumentar capital, pero también puede crear situación de pánico con reacciones excesivas e ineficientes, por lo que hay que considerar si se comunica y qué información y con qué granularidad se comunica al mercado.

$\textbf{Riesgo de contraparte}$

El riesgo de contraparte surge de las transacciones de un banco con socios comerciales, en lugar de prestatarios. El riesgo de contraparte se define como el riesgo de que la contraparte pueda incumplir antes de la liquidación final de los flujos de efectivo. En caso de producirse el banco sufriría una pérdida si el valor de mercado de la transacción es positivo (es decir, si la contraparte le debe al banco) en el momento del incumplimiento. 
    
Este tipo de riesgo puede aparecer incluso si la contraparte cumple totalmente con lo acordado, como resultado de la disminución del valor de los activos debido al deterioro de la calidad crediticia de la contrapartida. A diferencia de la exposición de un banco al riesgo de crédito a través de un préstamo, donde la exposición es unilateral y sólo el banco prestamista se enfrenta al riesgo de pérdida, el riesgo de crédito de contraparte crea un riesgo de pérdida bilateral. El valor de mercado de la transacción puede ser positivo o negativo tanto para el banco como para su contraparte, es incierto y puede variar con el movimiento de los factores subyacentes del mercado. 
    
Las operaciones de derivados son ejemplos de operaciones que dan lugar a un riesgo de contraparte. La caída de Lehman Brothers en 2008 marcó un punto de inflexión en la evaluación de la relevancia del riesgo de crédito de contraparte, al provocar la pérdida de cientos de miles de millones de dólares, lo que llevó al BCBS junto con los bancos a trabajar para establecer cámaras de compensación y contrapartes centrales (CCP) para los mercados de derivados, con el fin de reducir su riesgo y permitir negociar con la confianza de que los contratos se liquidarán en su totalidad.




\end{document}